\documentclass[11pt,letterpaper]{article}
\usepackage[T1]{fontenc}
\usepackage[utf8]{inputenc}
\usepackage{lmodern}
\usepackage[margin=1in,headheight=15pt]{geometry}
\usepackage{amsmath,amssymb,amsthm,mathtools,mathrsfs,aliascnt}
\usepackage{microtype,booktabs,tabularx,array,longtable}
\usepackage{xcolor,fancyhdr,titlesec,needspace,enumitem}
\definecolor{Ink}{HTML}{243747}
\definecolor{Accent}{HTML}{37666D}
\usepackage[colorlinks=true,linkcolor=blue!35!black,citecolor=green!30!black,urlcolor=blue!45!black]{hyperref}
\usepackage{xurl}
\usepackage[nameinlink,noabbrev]{cleveref}
\titleformat{\section}{\Large\bfseries\color{Ink}}{\thesection}{.65em}{}
\titleformat{\subsection}{\large\bfseries\color{Accent}}{\thesubsection}{.65em}{}
\pagestyle{fancy}\fancyhf{}
\fancyhead[L]{\small Hilbert geometry at surreal scales}
\fancyhead[R]{\small Merged research report}
\fancyfoot[C]{\thepage}
\setlength{\parskip}{.26em}
\setlength{\parindent}{1.1em}
\setlength{\emergencystretch}{3em}
\widowpenalty=10000\clubpenalty=10000
\setlist{itemsep=.15em,topsep=.3em,leftmargin=2em}
\numberwithin{equation}{section}
\newtheorem{theorem}{Theorem}[section]
\crefname{theorem}{Theorem}{Theorems}
\newaliascnt{proposition}{theorem}
\newtheorem{proposition}[proposition]{Proposition}
\aliascntresetthe{proposition}
\crefname{proposition}{Proposition}{Propositions}
\newaliascnt{lemma}{theorem}
\newtheorem{lemma}[lemma]{Lemma}
\aliascntresetthe{lemma}
\crefname{lemma}{Lemma}{Lemmas}
\newaliascnt{corollary}{theorem}
\newtheorem{corollary}[corollary]{Corollary}
\aliascntresetthe{corollary}
\crefname{corollary}{Corollary}{Corollaries}
\theoremstyle{definition}
\newaliascnt{definition}{theorem}
\newtheorem{definition}[definition]{Definition}
\aliascntresetthe{definition}
\crefname{definition}{Definition}{Definitions}
\newaliascnt{example}{theorem}
\newtheorem{example}[example]{Example}
\aliascntresetthe{example}
\crefname{example}{Example}{Examples}
\newaliascnt{question}{theorem}
\newtheorem{question}[question]{Question}
\aliascntresetthe{question}
\crefname{question}{Question}{Questions}
\theoremstyle{remark}
\newaliascnt{remark}{theorem}
\newtheorem{remark}[remark]{Remark}
\aliascntresetthe{remark}
\crefname{remark}{Remark}{Remarks}
\newaliascnt{warning}{theorem}
\newtheorem{warning}[warning]{Warning}
\aliascntresetthe{warning}
\crefname{warning}{Warning}{Warnings}
\crefname{section}{Section}{Sections}
\crefname{subsection}{Section}{Sections}
\crefname{appendix}{Appendix}{Appendices}
\crefname{table}{Table}{Tables}
\crefname{equation}{}{}
\crefformat{equation}{(#2#1#3)}
\crefrangeformat{equation}{(#3#1#4)--(#5#2#6)}
\crefmultiformat{equation}{(#2#1#3)}{ and~(#2#1#3)}{, (#2#1#3)}{ and~(#2#1#3)}
\crefname{enumi}{item}{items}
\newcommand{\R}{\mathbb R}
\newcommand{\C}{\mathbb C}
\newcommand{\Q}{\mathbb Q}
\newcommand{\Z}{\mathbb Z}
\newcommand{\N}{\mathbb N}
\newcommand{\kk}{\mathbb k}
\newcommand{\No}{\mathbf{No}}
\newcommand{\F}{F}
\newcommand{\K}{K}
\newcommand{\HH}[1]{\mathscr H_\Gamma(#1)}
\newcommand{\BB}[2]{\mathscr B_\Gamma(#1,#2)}
\newcommand{\BH}{\mathscr B_\Gamma(H,H)}
\newcommand{\Bpos}[2]{\mathscr B^{>0}_\Gamma(#1,#2)}
\newcommand{\OO}{\mathcal O}
\newcommand{\mm}{\mathfrak m}
\newcommand{\supp}{\operatorname{supp}}
\newcommand{\ran}{\operatorname{ran}}
\newcommand{\Dom}{\operatorname{Dom}}
\newcommand{\dom}{\operatorname{dom}}
\newcommand{\mul}{\operatorname{mul}}
\newcommand{\res}{\operatorname{res}}
\newcommand{\st}{\operatorname{st}}
\newcommand{\Graph}{\operatorname{Graph}}
\newcommand{\rank}{\operatorname{rank}}
\newcommand{\ind}{\operatorname{ind}}
\newcommand{\Span}{\operatorname{span}}
\newcommand{\End}{\operatorname{End}}
\newcommand{\GL}{\operatorname{GL}}
\newcommand{\diag}{\operatorname{diag}}
\newcommand{\Det}{\operatorname{Det}}
\newcommand{\coker}{\operatorname{coker}}
\newcommand{\Un}{\mathsf U}
\newcommand{\ip}[2]{\langle #1,#2\rangle}
\newcommand{\norm}[1]{\lVert #1\rVert}
\newcommand{\abs}[1]{\lvert #1\rvert}
\newcommand{\coeff}[2]{[t^{#1}]#2}
\newcommand{\sH}{\mathop{\sum}\nolimits^{\mathrm H}}
\newcommand{\cH}{\mathcal H}
\newcommand{\cG}{\mathcal G}
\newcommand{\Spl}{\operatorname{Spl}}
\newcommand{\src}[1]{\textup{\textsf{\small[#1]}}}
\newcommand{\mergetag}{\textup{\textsf{\small[merge]}}}
\newcommand{\replabel}[1]{\nolinkurl{#1}}
\newcommand{\fn}[1]{\nolinkurl{#1}}
\newcommand{\sSOT}{\mathop{\sum}\nolimits^{\mathrm{SOT}}}
\newcommand{\Bnn}{\mathscr B^{\ge0}_\Gamma(H,H)}
\newcommand{\Dg}{\mathscr D}
\newcommand{\fin}{\mathrm{fin}}
\newcommand{\cb}{\mathrm{cb}}
\newcommand{\Oz}{\mathbf{Oz}}
\newcolumntype{Y}{>{\raggedright\arraybackslash}X}
\hypersetup{pdftitle={Hilbert Geometry at Surreal Scales},pdfauthor={Merged research report for the Surreal project},pdfsubject={Surreal and surcomplex Hahn--Hilbert spaces; orthogonal projections; amplified and infinitesimal graphs; projection-lattice obstructions; infinite simultaneous unitary straightening; metric rigidity; Moore--Penrose inverses}}
\begin{document}
\hypersetup{pageanchor=false}
\begin{titlepage}
\centering
\vspace*{2mm}
{\large\scshape Merged research report for the Surreal project\par}
\vspace{7mm}
{\Huge\bfseries\color{Ink} Hilbert Geometry\\[.15em] at Surreal Scales\par}
\vspace{6mm}
{\Large Orthogonal splitting, amplified and infinitesimal graphs,\\[.2em]
projection-lattice obstructions, infinite unitary straightening,\\[.2em]
metric rigidity, and Fredholm least squares\par}
\vspace{6mm}
{\large Built from three manuscripts dated September 23, 2026\par}
\vspace{5mm}
\begin{minipage}{.96\textwidth}
\begin{abstract}\noindent
We develop Hilbert geometry over set-sized surreal and surcomplex Hahn
fields while retaining an ordinary Hilbert space at each coefficient. The
scalar-valued inner product is positive definite, its norm exists without a
divisible value group, and the valued space is spherically complete. These
facts do not imply the projection theorem.

Using the collection's automatic adjoint theorem, every orthogonally
complemented subspace is an infinitesimal unitary deformation of an ordinary
closed subspace, equivalently the graph of a positive-order Hahn series of
bounded operators over its residue. For an ordinary bounded $T$ and an
infinite Hahn scalar $a$, $\Graph(aT)$ is complemented if and only if
$\ran T$ is closed; for an ordinary closed densely defined $A$ and a positive
infinitesimal $q$, the graph of $qA$ is orthoclosed and spherically complete
but complemented if and only if $A$ is bounded. Both are instances of one
criterion: an infinitesimally scaled closed linear relation is complemented
exactly when its domain is closed. For a nontrivial value group the poset of
complemented subspaces is a lattice exactly when the coefficient space is
finite dimensional; a positive operator with three exponents has a
noncomplemented kernel, and a countable orthogonal family of rank-one
projections has no join; an orthogonal family of projections is conjugated
to its residues by one unitary exactly when its coefficient columns are
bounded and have well-ordered support. A positive Hahn metric with strictly
positive residue is isometric to the canonical one exactly when its residue
is uniformly positive.

A positive counterpart: every positive-order Hahn perturbation of an
ordinary Fredholm operator has an adjointable Moore--Penrose inverse, and a
finite Schur matrix determines its kernel, cokernel, index and exact inverse
valuation. Adjointable operators on the class of all surcomplex vectors have
one set-sized support.

This report is built from three manuscripts; it says which source proved
each result, prints shared results once, and keeps every source's
limitations. It is an AI-assisted draft: it has not been refereed,
and nothing in it has been formalized in Lean.
\end{abstract}
\end{minipage}
\vfill
\begin{minipage}{.96\textwidth}\small
\textbf{Sources and status.} Manuscripts 03 (the base text), 06 and 07
(added later), pinned to commits \texttt{9a385d3}, \texttt{a45d722} and
\texttt{fc27d87}; titles in \cref{hgeo:sub:provenance,hgeo:us:sub:source}. Results
already printed elsewhere in the collection are cited, not reprinted. The
construction, the automatic adjoint theorem, the direct-rotation and polar
formulas and the classical perturbation tools are credited, not claimed. The provenance of
every section is recorded in \cref{hgeo:app:provenance}; non-claims are
collected in \cref{hgeo:app:nonclaims}. Finite checks accompany the report
and check finite identities only. No named published conjecture is asserted
to be solved.
\end{minipage}
\end{titlepage}
\hypersetup{pageanchor=true}
\pagenumbering{roman}
\setcounter{tocdepth}{2}
\tableofcontents
\clearpage
\pagenumbering{arabic}

\section{The question, the model, and the results}
\label{hgeo:sec:intro}

Replacing real or complex vector components by surreal or surcomplex
numbers is not merely a change of notation. In infinite dimension one
must choose which coordinate families are vectors, which infinite sums
are legitimate, which topology defines completeness, and which operators
are admitted. Different choices produce genuinely different theories. This
report makes one such choice and proves both its strengths and its
unavoidable limitations.

\subsection{The model}\label{hgeo:sub:model}
We use the coefficient-Hilbert model already constructed in Part~I of the
collection's report \emph{Infinite-Dimensional Hahn Spectral Theory}
(directory \fn{docs/surcomplex/infinite-dimensional-hahn-spectral-theory},
label prefix \replabel{ihs:}) \cite{RepoSpectral}:
\[
 \cH=\HH{H}=H((t^\Gamma)),\qquad
 \F=\R((t^\Gamma)),\qquad \K=\C((t^\Gamma))=\F[i].
\]
Here $H$ is an ordinary Hilbert space and $\Gamma$ is a set-sized ordered
abelian group, nonzero unless the contrary is said. A vector is a formal
series of ordinary Hilbert vectors with well-ordered support. Through a
chosen normal-form embedding of the scalar field, its coordinates are actual
surcomplex numbers. There is also an entirely real version, with $H$ real
and $\F$ in place of $\K$; except where said otherwise the text is written
in the complex case and every statement and proof holds verbatim in the
real case. Its duality and completion are developed in the collection's
report \emph{Three Duals at Surreal Scales} (prefix \replabel{duals:})
\cite{RepoDuals}.

This model retains every ordinary vector of $H$, not only vectors with
finite coordinate support. It also retains ordinary bounded operators
as constant coefficients. It does \emph{not} retain ordinary Hilbert
norm convergence as the topology on the constant copy of $H$.
The distinction is crucial both for what works and for what fails.

\subsection{Main results}\label{hgeo:sub:results}
Call a $\K$-linear subspace $M\subseteq\cH$ \emph{orthogonally split}, or
\emph{orthogonally complemented}, when $\cH=M\oplus M^\perp$. This is
equivalent to every vector having a nearest point in $M$, and to $M$ being
the range of a self-adjoint idempotent (\cref{hgeo:prop:projection}).
\begin{enumerate}[label=\textup{(\arabic*)},leftmargin=2.2em]
\item \emph{Classification of complemented subspaces}
\src{03, 06}. Every split $M$ is $M=U\HH{M_0}$ for a unitary $U$ with
$v(U-I)>0$ and a unique ordinary closed $M_0=\res M$; equivalently it is the
graph of a positive-order Hahn series in $\BB{M_0}{M_0^\perp}$
(\cref{hgeo:thm:normalform,hgeo:cor:graphchart}). Given the collection's
rotation lemma \replabel{ihs:fr:lem:rotation}, this is a corollary of an
integrality lemma for projections (\cref{hgeo:lem:integralproj}). Closed
residue is necessary but not sufficient (\cref{hgeo:prop:closedresidue}).
\item \emph{Unitaries} \src{06}. Unitary orbits of projections are
classified by two ordinary dimensions (\cref{hgeo:cor:orbits}); every
adjointable unitary is $\widehat{U_0}\exp X$ uniquely
(\cref{hgeo:thm:unitary}); finite orthogonal decompositions straighten
simultaneously (\cref{hgeo:thm:finite}).
\item \emph{Two graph theorems and their common criterion} \src{03}, \src{06},
\mergetag. For ordinary bounded $T:H\to G$ and $a\in\K$,
\[
 \Graph(aT)\text{ is split}
 \quad\Longleftrightarrow\quad
 a=0\text{ or }v(a)\ge0\text{ or }\ran T\text{ is closed},
\]
and for $v(a)<0$ the residue is $\ker T\oplus\ran T$
(\cref{hgeo:thm:threshold}). For ordinary closed densely defined
$A$ and a positive infinitesimal monomial $q=t^\eta$, the graph $M_q(A)$ of
$qA$ is orthoclosed and spherically
complete, with residue $\Dom A\oplus0$, and is split if and only if $A$ is
bounded (\cref{hgeo:thm:unboundedgraph}). Neither theorem implies the other;
both splitting criteria follow from \cref{hgeo:thm:relation}: for a closed
linear relation $R$ and $v(b)>0$, the scaled relation $M_b(R)$ has residue
$\dom R\oplus\mul R$ and is split exactly when $\dom R$ is closed.
\item \emph{The projection poset} \src{03, 06}. For $\Gamma\ne0$,
$\Spl(\HH{H})$ is a lattice if and only if $\dim H<\infty$
(\cref{hgeo:thm:lattice}); a positive operator with exponents $0,\eta,2\eta$
has a noncomplemented kernel (\cref{hgeo:thm:positivekernel}); a countable
family of mutually orthogonal rank-one projections has no least upper bound,
and countable simultaneous straightening fails for two distinct reasons
(\cref{hgeo:thm:nojoin,hgeo:cor:nosimultaneous,hgeo:prop:incoherent}).
\item \emph{Metric rigidity} \src{06}. A positive metric
$\Theta=\widehat{\Theta_0}+\Theta_+$ with strictly positive residue is
isometric to the canonical one by some linear bijection, with no regularity
assumed, exactly when $\Theta_0$ is uniformly positive
(\cref{hgeo:thm:metric}).
\item \emph{Fredholm least squares} \src{03}. For $A=T_0+E$ with $T_0$
ordinary Fredholm and $v(E)>0$, with $m=\dim\ker T_0$, $n=\dim\ker T_0^*$
and $r$ the rank of a finite Schur matrix $S$,
\[
 \dim_\K\ker A=m-r,\qquad
 \dim_\K(\cG/\ran A)=n-r,\qquad
 \ind_\K A=m-n,
\]
$A^\dagger$ exists in the adjointable Hahn algebra, and for $r>0$
\[
 v(A^\dagger)=-\bigl(\nu_r(S)-\nu_{r-1}(S)\bigr),
 \qquad \nu_0(S)=0,
\]
where $\nu_j$ is the least valuation of a nonzero $j$-minor
(\cref{hgeo:thm:schur,hgeo:thm:MP,hgeo:thm:scale}). No Archimedean
hypothesis, divisibility, analytic convergence or bound on all coefficient
operator norms is imposed. These theorems contain the determinant-free
clauses of the collection's Hahn Fredholm alternative
\replabel{ihs:fr:thm:fredholm} (\cref{hgeo:rem:fredholm}).
\item \emph{The whole surreal class} \src{03, 06}. Working with all of $\No$
or $\No[i]$ in a class theory without global choice, every class-linear map
with an everywhere-defined class adjoint has one set-sized Hahn expansion in
ordinary bounded operators, so the whole theory applies to literal surreal
and surcomplex vector components (\cref{hgeo:thm:class,hgeo:cor:globalgeometry}).
\item \emph{Infinite simultaneous straightening} \src{07}. A set-indexed
orthogonal family of projections $P_i$ with residues $E_i$ has one unitary
$U$ with $P_i=UE_iU^*$ for all $i$ if and only if, at every exponent, the
canonical columns $\sum_i(P_i)_\gamma E_i$ synthesize to a bounded operator,
and their active support is well ordered. An explicit $U$ is given for
complete and incomplete families; the criterion is equivalent to a unique
normal completely bounded atomic calculus, and it gives joins. A complete
family with uniformly bounded coefficients at every order fails it, and one
of its subfamilies has no join
(\cref{hgeo:us:thm:main,hgeo:us:thm:calculus,hgeo:us:thm:uniform}).
\end{enumerate}

These statements are proved here, but proving a statement here is not a
claim that nobody has previously proved it. The sources name their own
candidate contributions and credit their inputs; the literature comparisons
were targeted and priority is not certified (\cref{hgeo:sec:provenance},
\cref{hgeo:app:nonclaims}).

\subsection{Conventions and renamed symbols}\label{hgeo:sub:conventions}
The two sources wrote the same objects with different letters, and some of
their letters collide with the collection's usage. The following convention
is fixed for the whole report; \cref{hgeo:tab:renames} lists every renaming.
\begin{itemize}
\item $\F=\R((t^\Gamma))$ is the ordered real Hahn field and
$\K=\C((t^\Gamma))=\F[i]$ the complex one, as in source~03, in the spectral
report (\replabel{ihs:}, where $\F$ is written $\mathrm{Fld}$) and in the
finite spectral-theory report ($F_\Gamma$, $K_\Gamma$). Source~06 had these
two letters swapped ($K$ real, $F$ complex); they are swapped back
throughout. $\kk\in\{\R,\C\}$ is the coefficient field of $H$.
\item $\HH{H}=H((t^\Gamma))$ is the whole Hahn--Hilbert space. Source~06
wrote $\mathcal E_\Gamma(H)$; that letter is not used here, because in
the three-duals report $E_\Gamma(V)=\Span_\K V\cong\K\otimes V$ is the
\emph{algebraic} scalar extension, in general a proper subspace of
$V((t^\Gamma))$.
\item $\BB{H}{G}=\mathcal B(H,G)((t^\Gamma))$ (source~06:
$\mathcal A_\Gamma(H,G)$; the spectral report writes $\mathcal A_\Gamma(H)$
for $\BH$), $v(A)=\min\supp A$ (06: $v_{\mathrm{op}}$), and
$\Bpos{H}{G}$ is the set of series supported in $\Gamma_{>0}$, zero
included. A hat denotes the constant lift of an ordinary operator.
\item \emph{Split} and \emph{orthogonally complemented} are synonyms; a
\emph{projection} is an everywhere-defined self-adjoint idempotent (06:
``projector''). \emph{Closed} means valuation closed; \emph{orthoclosed}
means $M=M^{\perp\perp}$. The two are never used interchangeably.
\item $q=t^\eta$ with $\eta\in\Gamma_{>0}$ is a fixed positive
infinitesimal monomial (06: $\tau=t^\delta$).
\item On $\ell^2=\ell^2(\N_{\ge1})$, $T=\diag(1/n)$ (06: $S$; the spectral
report writes $D$), $\Lambda=T^{-1}=\diag(n)$ with $\Dom\Lambda=\ran T$
(06: $D$), and $h=(1/n)_{n\ge1}\in\ell^2\setminus\ran T$.
\item $\Un(H)$ is the ordinary unitary group and $\Un(\cH)$ the group of
adjointable unitaries (06: $\mathcal U$; the spectral report uses
$\mathcal U(T)$ for the set of norm bounds of $T$).
\item The Fredholm blocks and auxiliary operators of source~03 are renamed
where they met the field letters or the rotation: see
\cref{hgeo:tab:renames}.
\end{itemize}
Inner products are linear in the first argument. The valuation of a vector
or an operator is the least exponent of its support; $v(0)=+\infty$.
\Cref{hgeo:us:sec:intro,hgeo:us:sec:criterion,hgeo:us:sec:consequences,%
hgeo:us:sec:calculus,hgeo:us:sec:uniform,hgeo:us:sec:support,%
hgeo:us:sec:surreal,hgeo:us:sec:questions}, added with source~07, fix further
conventions in \cref{hgeo:us:sub:conventions}.

\subsection{Provenance of this report}\label{hgeo:sub:provenance}
This report is built from two manuscripts written independently and dated
September 23, 2026, items 03 and 06 of the twenty-seventh batch placed in
commit \texttt{a4dcb91}; their shipped file prefixes are
\fn{03-orthogonal-splitting-} and \fn{06-hilbert-foundations-}.
A third manuscript, source~07 (batch~34, manuscript~09, pinned to
\texttt{fc27d87}, prefix \fn{07-unitary-synthesis-}), was written later
against this report and is added after the merged text, in
\cref{hgeo:us:sec:intro,hgeo:us:sec:criterion,hgeo:us:sec:consequences,%
hgeo:us:sec:calculus,hgeo:us:sec:uniform,hgeo:us:sec:support,%
hgeo:us:sec:surreal,hgeo:us:sec:questions}; its provenance and the choices
of that addition are in \cref{hgeo:us:sub:source}. The rest of this subsection
concerns the merge of sources~03 and~06.
Source~03 pins repository commit \texttt{9a385d3} (35 commits before that
placement) and source~06 pins \texttt{a45d722} (30 commits before it). None
of the reports cited below changed between the pins and the placement.

\begin{center}
\small
\renewcommand{\arraystretch}{1.15}
\begin{tabular}{>{\raggedright\arraybackslash}p{0.07\textwidth}>{\raggedright\arraybackslash}p{0.34\textwidth}>{\raggedright\arraybackslash}p{0.49\textwidth}}
\toprule
Source & Manuscript & Principal contribution to this report\\
\midrule
03 & \emph{Hilbert Geometry at Surreal Scales: Orthogonal Splitting,
Amplified Graphs, and Fredholm Least Squares} (29 pp.) &
Base text and structure; the residue normal form and graph chart; the
amplification threshold for bounded operators; the lattice dichotomy with
its finite-dimensional converse; the Fredholm, Moore--Penrose and
inverse-valuation theorems; scale coherence; the class descent in NBG
without global choice; fine-net stabilization.\\
06 & \emph{Hilbert Geometry at Surreal Scales: Orthogonal subspaces,
infinitesimal graphs, metric rigidity, and projection-lattice obstructions}
(25 pp.) &
The $c_{00}$ example; the strict-contraction functional; the classification
in its own form, with the unitary-orbit corollary; the unitary group and
finite simultaneous straightening; the closed densely defined graph theorem
and the distance cut; the positive operator with a noncomplemented kernel;
the missing countable join and the two obstructions to countable
straightening; metric rigidity; the global geometry corollary; the Sol\`er
framing.\\
\bottomrule
\end{tabular}
\end{center}

\paragraph{Where the merge had to choose.}
\begin{itemize}[leftmargin=*]
\item \emph{Base.} Source~03 is the base, because on every shared theorem
it has the weaker hypotheses: the lattice theorem for every
infinite-dimensional $H$ with the finite-dimensional converse (06:
one direction, for $\ell^2$), the class descent without global choice (06:
with global choice), and an arbitrary amplification scalar (06: a
monomial). The straightening theorem is identical in both.
\item \emph{Printed once.} The positive form, spherical completeness,
automatic adjoints, the integral-projection lemma, the normal form, the graph
chart, the no-meet construction (03's and 06's are the same up to swapping
coordinates), the class descent and fine discreteness are proved by both
sources and printed once, crediting both. Two genuinely different proofs of
nonnegativity of projections are both kept (\cref{hgeo:lem:integralproj}).
\item \emph{Printed by citation.} The automatic adjoint theorem and its
bound (\replabel{ihs:hh:thm:adjoint}, \replabel{ihs:hh:cor:autobounded}),
reproved by both sources, are cited; the rectangular form is obtained from
the square one by a two-line reduction \mergetag. So are the norm-attainment
criterion (\replabel{ihs:hh:thm:norm}), the coefficientwise Riesz criterion
(\replabel{duals:thm:riesz}), the direct rotation
(\replabel{ihs:fr:lem:rotation}), the finite Moore--Penrose and
determinantal-scale theorems of the finite spectral-theory report
(\replabel{thm:leastsquares}, \replabel{thm:scales},
\replabel{lem:DeltaInvariant}), and fine discreteness
(\replabel{a:thm:discrete}). Source proofs are kept only where they prove
more or proceed differently.
\item \emph{The two graph theorems} are related but neither implies the
other: 03's concerns bounded operators with arbitrary kernel and range and
also finite scalars; 06's concerns unbounded operators and also computes the
orthogonal complement. Both are printed with their proofs, and their common
splitting criterion is added by the merge with a complete proof
(\cref{hgeo:sub:relation}).
\end{itemize}
Results added by the merge carry the tag \mergetag{} and have complete
proofs. \Cref{hgeo:app:provenance} records, section by section, which source
proved what; \cref{hgeo:app:stale} lists the repository statements of the
sources that were corrected or completed.

\section{Scalars, supports, and the surreal interpretation}
\label{hgeo:sec:scalars}

\subsection{Standing conventions}
An ordered abelian group is written additively. A subset is well ordered
in its given order, not in the reverse order. For a vector space $V$ over
$\kk$,
\[
 V((t^\Gamma))=
 \left\{x=\sum_{\gamma\in\Gamma}x_\gamma t^\gamma:
 x_\gamma\in V,\ \supp x\text{ well ordered}\right\}.
\]
Set $v(0)=+\infty$ and $v(x)=\min\supp x$ otherwise. The exponent order is
the valuation order: positive exponents are small. The field $\F$ is
ordered by its leading nonzero coefficient, and complex conjugation on $\K$
acts on coefficients and fixes monomials. The valuation ring and its maximal
ideal are
\[
 \OO=\{a\in\K:v(a)\ge0\},\qquad
 \mm=\{a\in\K:v(a)>0\}.
\]
Zero is included. Residue is extraction of the coefficient of $t^0$; on
finite real scalars it is the standard part $\st$. A scalar is finite in the
ordered-field modulus precisely when it is in $\OO$. An infinite scalar has
negative valuation.

\begin{lemma}[Hahn--Neumann support calculus \src{06}]\label{hgeo:lem:support}
If $S,S'$ are well-ordered subsets of an ordered abelian group, then $S+S'$
is well ordered and, for every $\alpha$, the set
$\{(s,s')\in S\times S':s+s'=\alpha\}$ is finite. If $S\subseteq\Gamma_{>0}$
is well ordered, its generated additive monoid $S^*$ is well ordered and each
exponent has only finitely many representations by words in $S$, counting
all word lengths, apart from the unique empty word at zero.
\end{lemma}

The second assertion is the positive-support form of Neumann's lemma
\cite{Neumann}; with the first it is the spectral report's support lemma
\replabel{ihs:lem:support}, where a support-engine proof and attribution are
given \cite{RepoSpectral}. It is valid without an Archimedean hypothesis.
For the first finiteness statement, infinitely many distinct first
coordinates would have a strictly increasing subsequence, forcing a strictly
decreasing subsequence of second coordinates, contrary to well-ordering. The
same descending-sequence argument gives well-ordering of the sum set. The
positive-word assertion is the classical support theorem, imported rather
than inferred from convergence.

\begin{definition}[Strong Hahn summation]\label{hgeo:def:strong}
A set-indexed family $(x_j)_{j\in J}$ in $V((t^\Gamma))$ is strongly
Hahn summable if the union of its supports is well ordered and, at each
exponent, only finitely many members have a nonzero coefficient.
The sum $\sH_jx_j$ is formed by these finite coefficient sums. A linear map
is \emph{strong} if it preserves strongly summable families and their sums.
\end{definition}

This is not the strong operator topology and not ordinary convergence.
Modern treatments of formal summability and strongly linear maps include
\cite{Bagayoko,Freni}; the distinction is central in the three-duals report.
\Cref{hgeo:lem:support} justifies every formal multiplication and
positive-order power series used below. In particular, in a unital
coefficient algebra,
\begin{equation}\label{hgeo:eq:neumann}
 (I+B)^{-1}=\sH_{j\ge0}(-B)^j\qquad (v(B)>0).
\end{equation}
For noncommutative coefficient rings this Neumann inverse is the spectral
report's \replabel{ihs:lem:neumann}, whose Lean formalization is recorded in
\fn{docs/FORMALIZATION.md}. Binomial series $(I+B)^c$ for real $c$, and
$e^B$, $\log(I+B)$, are also defined by \cref{hgeo:lem:support}
(\replabel{ihs:cor:binomial}; see \cref{hgeo:sub:calculus}). These are
identities of Hahn series. They do not presuppose topological convergence of
the finite partial sums.

\begin{example}[Summable families not converging through finite sums
\src{03, 06}]\label{hgeo:ex:partialsums}
Two families show the phenomenon.
\begin{enumerate}[label=\textup{(\alph*)}]
\item The family $t^{n/(n+1)}$, $n\ge1$, in $\C((t^\Q))$ is strongly
summable. Its exponents increase to $1$. The tail of every finite partial
sum still has valuation below $1$, so the partial sums do not converge in
the valuation topology \src{03}.
\item Let $\Gamma=\R\times\R$ with lexicographic order and $\delta=(0,1)$.
The strong sum $\sum_{n\ge0}t^{n\delta}$ exists and equals
$(1-t^\delta)^{-1}$. Its finite partial sums have errors of valuation
$(0,N+1)$, which never exceeds $(1,0)$, so they do not converge either. None
of the positive-order operator power series below is justified by assuming
$n\delta\to+\infty$ \src{06}.
\end{enumerate}
This distinction is already emphasized in the collection's foundations and
duality reports \cite{RepoFound,RepoDuals}.
\end{example}

\subsection{Actual surreal and surcomplex components}
Fix an order-preserving additive embedding
$\iota:\Gamma\hookrightarrow(\No,+)$. For example, the usual
$\Z,\Q,\R$ can be embedded using their standard surreal copies.
Conway normal forms give the embedding
\begin{equation}\label{hgeo:eq:embedding}
 \sum_\gamma r_\gamma t^\gamma
 \longmapsto
 \sum_\gamma r_\gamma\omega^{-\iota(\gamma)}.
\end{equation}
The support on the right is reverse well ordered in its exponents,
exactly as a surreal normal form requires. The right-hand sum is a surreal
normal form, not the limit of its finite partial sums in the full fine
topology. Addition, multiplication, order, and strong summation are
preserved; adjoining $i$ gives an embedding into $\No[i]$
\cite{Conway,RepoFound}. If $\Gamma$ is divisible, $\F$ is real closed and
$\K$ algebraically closed; this familiar fact is used only where it is
mentioned, and the norm construction below needs less.

For a fixed ordinary orthonormal basis $(e_j)$ of $H$, every vector in
$\HH{H}$ therefore has coordinates in a specified subfield of
$\No[i]$. \Cref{hgeo:sec:coordinates} gives the exact coordinate
admissibility condition.

\subsection{Why the whole surreal class is not silently a set field}
A nonzero vector space over all of $\No$ or $\No[i]$ cannot have a
set-sized underlying collection: multiplication of any fixed nonzero
vector by all scalars would inject a proper class into that set.
Thus a literal global theory requires class-sized objects or a
universe-relative formulation. The main theorems below are set-sized
and can be transported by \eqref{hgeo:eq:embedding}; \cref{hgeo:sec:global}
states a separate class-theoretic extension.

For any specified set of scalar normal forms, their supports lie in a
common set-sized additive subgroup of surreal exponents. Thus particular
calculations take place in a common Hahn workspace. This does not mean
that every class-sized problem has been reduced to a single set field,
nor that the full surreal fine topology equals the intrinsic topology of
that workspace.

\section{The Hahn--Hilbert space and its completeness}
\label{hgeo:sec:construction}

\subsection{The inner product and its norm}
For $x,y\in\HH{H}$ put
\begin{equation}\label{hgeo:eq:inner}
 \ip{x}{y}=\sum_{\alpha\in\Gamma}
   \Bigl(\sum_{\beta+\gamma=\alpha}\ip{x_\beta}{y_\gamma}_H\Bigr)t^{\alpha}.
\end{equation}
At each exponent only finitely many coefficient pairs contribute
(\cref{hgeo:lem:support}). The ordinary Hilbert inner products are evaluated
\emph{before} the Hahn convolution, possibly as ordinary infinite Hilbert
sums; these are two different levels of summation. This is the spectral
report's coefficient-Hilbert inner product.

\begin{proposition}[Positive-definite Hahn extension
\src{03, 06}]\label{hgeo:prop:inner}
The form \eqref{hgeo:eq:inner} is Hermitian and positive definite. If
$x\ne0$, $\delta=v(x)$ and $u=x_\delta$, then
\[
 \ip{x}{x}=t^{2\delta}\norm{u}_H^2(1+r),\qquad r=0\text{ or }v(r)>0,
\]
and $\ip{x}{x}$ has a unique positive square root in $\F$, regardless of
whether $\Gamma$ is divisible. With $\norm{x}=\sqrt{\ip{x}{x}}$,
\[
 v(\norm{x})=v(x),\qquad v(\ip{x}{x})=2v(x)\quad(x\ne0).
\]
Cauchy--Schwarz, the triangle inequality, the parallelogram identity, and
finite Pythagoras hold over $\F$.
\end{proposition}

This is \replabel{ihs:hh:prop:inner} of the spectral report (and
\replabel{duals:prop:norm} of the three-duals report); both sources reprove
it. For the reader's convenience: the only pair contributing at $2\delta$ is
$(\delta,\delta)$, so the leading coefficient is $\norm{u}_H^2>0$; the root
is $t^\delta\norm{u}_H\sH_{j\ge0}\binom{1/2}{j}r^j$, a strong sum by
\cref{hgeo:lem:support} whose square is $\ip{x}{x}$ by the binomial identity;
uniqueness follows from $(a-b)(a+b)=0$ in an ordered field. Expanding the
nonnegative $\ip{x-by}{x-by}$ with $b=\ip{x}{y}/\ip{y}{y}$ gives
Cauchy--Schwarz, and the triangle inequality follows. No supremum or limiting
argument is used. The same construction gives $\abs a=\sqrt{a\bar a}$ for
$a\in\K$.

\begin{warning}\label{hgeo:warn:roots}
The assertion is not that every positive element of $\F$ has a square
root. For example, $t$ has no square root when $\Gamma=\Z$.
Squared norms have even leading valuation, which is exactly why all roots
needed here do exist.
\end{warning}

\subsection{Valuation topology and spherical completeness}
A neighborhood basis of $0$ is $\{x:v(x)>\gamma\}$, $\gamma\in\Gamma$.
It is the same topology as that generated by all positive $\F$-valued norm
radii: if $v(x)>v(r)$ then $\norm{x}<r$, while a monomial radius of exponent
strictly greater than $\gamma$ forces $v(x)>\gamma$. In rank one, when
$\Gamma\subseteq\R$, $d(x,y)=\exp(-v(x-y))$, with $d(x,x)=0$, is a
real-valued ultrametric for the same topology. It is not the $\F$-valued
Hilbert norm itself. This is the topology of a fixed workspace, not the
topology induced by all positive tolerances in the full surreal field
(\cref{hgeo:prop:globalnets}).

\begin{theorem}[Spherical completeness of the coefficient space
\src{03, 06}]\label{hgeo:thm:complete}
For every $\kk$-vector space $V$, with no completeness or topology assumed on
$V$, every nonempty nested set of balls
\[
 B(a,\gamma)=\{x:v(x-a)\ge\gamma\}
       \quad(a\in V((t^\Gamma)),\ \gamma\in\Gamma)
\]
has nonempty intersection. In particular $V((t^\Gamma))$ is complete
for valuation-Cauchy nets. Thus $\HH{H}$ is spherically complete even
when its orthogonal projection theorem fails.
\end{theorem}

This is the three-duals report's \replabel{duals:prop:spherical}; source~06
credits it there and both sources reprove it. In brief: a ball fixes all
coefficients strictly below its radius exponent; nested balls impose
compatible prescriptions; the vector $x$ with the prescribed coefficients and
zeros elsewhere has well-ordered support, because any subset $S'$ of it and
any $s\in S'$ give, below $s$, a subset of the well-ordered support of one
ball's centre; hence $x$ lies in every ball. A Cauchy net has compatible
stabilized truncations, and the same argument produces its limit. The
argument uses no norm convergence in $V$. Spherical completeness here is for
valuation balls only, not for arbitrary ordered-field norm balls.

The last sentence is informative: completeness here does not force
ordinary completeness on coefficient subspaces. If $V\subset H$ is
any ordinary linear subspace, even one that is not ordinarily closed,
then $V((t^\Gamma))$ is valuation closed in $H((t^\Gamma))$: a coefficient
outside $V$ cannot be changed by a sufficiently small valuation
perturbation. This is one source of the projection failure.

\begin{proposition}[A first nonprojectable complete subspace
\src{06}]\label{hgeo:prop:c00}
Let $H=\ell^2(\N,\kk)$ and $M=c_{00}((t^\Gamma))\subseteq\HH{H}$, where
$c_{00}$ is the space of finitely supported sequences. Then $M$ is proper,
valuation closed, and spherically complete, but $M^\perp=0$. No point
outside $M$ has a nearest point in $M$.
\end{proposition}
\begin{proof}
Closedness is the preceding paragraph and spherical completeness is
\cref{hgeo:thm:complete} with $V=c_{00}$. The constants $e_n$ belong to $M$,
so orthogonality to $M$ forces every coordinate of every coefficient to
vanish; thus $M^\perp=0$. If $x\notin M$ and $m\in M$, then $r=x-m\ne0$ has
some $\ip{r}{e_n}\ne0$, and
\[
 \norm{x-(m+\ip{r}{e_n}e_n)}^2
 =\norm{r}^2-\abs{\ip{r}{e_n}}^2<\norm{r}^2.
\]
The correction belongs to $M$, excluding a nearest point.
\end{proof}

This subspace is \emph{not} orthoclosed: $M^{\perp\perp}=\HH{H}$. It is
therefore not yet a witness to the failure of orthomodularity;
\cref{hgeo:thm:unboundedgraph} supplies an orthoclosed one. The same
phenomenon, a closed proper subspace with zero orthogonal complement,
appears in the spectral report's closed-range theorem
\replabel{ihs:hh:thm:closedrange} and in the three-duals report's closed
hyperplane \replabel{duals:thm:hyperplane}.

\subsection{Finite-dimensional Hilbert geometry survives}
\begin{proposition}[\src{03}]\label{hgeo:prop:finite}
Every finite-dimensional $\K$-linear subspace of $\HH{H}$ has an
orthonormal basis and an orthogonal projection. If
$u_1,\ldots,u_d$ is such a basis, then
\[
 P_Mx=\sum_{j=1}^d\ip{x}{u_j}u_j.
\]
The coefficient-zero vectors $(u_j)_0$ form an ordinary orthonormal
family, and the residue of $M$ has ordinary dimension $d$.
\end{proposition}
\begin{proof}
Finite Gram--Schmidt uses only positive squared norms and their roots
from \cref{hgeo:prop:inner}. The displayed formula is then a
self-adjoint idempotent onto $M$. Since $\norm{u_j}=1$, one has
$v(u_j)=0$, and taking constant coefficients of the orthogonality
identities gives ordinary orthonormality.
For $x=\sum a_ju_j$, positivity gives
$v(x)=\min_jv(a_j)$; there is no cancellation of its least coefficient
because the leading coefficient vectors of the $u_j$ are orthonormal.
Thus integral vectors in $M$ have precisely the residue span of those
$d$ vectors.
\end{proof}

\section{Coordinates, two summation levels, and duality}
\label{hgeo:sec:coordinates}

\subsection{An exact coordinate-space description}
Choose an ordinary orthonormal basis indexed by a set $I$, identifying
$H$ with $\ell^2(I,\kk)$. The coordinate map sends
$x=\sum x_\gamma t^\gamma$ to
$a_i=\sum_\gamma(x_\gamma)_i t^\gamma$.

\begin{proposition}[Admissible surcomplex coordinate arrays
\src{03, 06}]\label{hgeo:prop:coordinates}
The coordinate map identifies $\HH{\ell^2(I)}$ with the families
$(a_i)\in\K^I$ satisfying both conditions:
\[
 \bigcup_{i\in I}\supp a_i\text{ is well ordered},\qquad
 \bigl(\coeff{\gamma}{a_i}\bigr)_{i\in I}\in\ell^2(I)
       \quad\text{for every }\gamma.
\]
The inner product is
\begin{equation}\label{hgeo:eq:twosums}
 [t^\delta]\ip{x}{y}
 =\sum_{\alpha+\beta=\delta}
     \left(\sum_{i\in I}(x_\alpha)_i
                         \overline{(y_\beta)_i}\right).
\end{equation}
The outer sum is finite and each inner sum is ordinarily absolutely
convergent.
\end{proposition}
\begin{proof}
The conditions are precisely that the coefficient rows define ordinary
Hilbert vectors and that their nonzero support is a Hahn support; the
support of the vector is exactly the displayed union.
They provide the inverse coordinate map. Cauchy--Schwarz in the ordinary
$\ell^2(I)$ proves absolute convergence of the inner sum; the support
lemma proves finiteness of the outer sum.
\end{proof}

\begin{example}[Why a single formal coordinate sum is inadequate
\src{03, 06}]
The constant vector $x=(1/n)_{n\ge1}$ belongs to $\HH{\ell^2}$ and has
its usual real squared norm $\sum_{n\ge1}n^{-2}$. But the scalar family
$(1/n^2)_{n\ge1}$ is not strongly Hahn summable: infinitely many nonzero
terms occur at exponent zero. Formula~\eqref{hgeo:eq:twosums} uses ordinary
Hilbert summation at a fixed coefficient, not that inadmissible formal sum.
Conversely, $x=\sum_{n\ge1}e_n t^{n/(n+1)}$ is a legitimate Hahn vector even
though its monomial partial sums do not converge in valuation.
\end{example}

Thus ``square summable surcomplex coordinates'' must include a summation
rule. Merely requiring finite partial squared sums to be bounded is
far too weak: an infinite scalar already bounds every ordinary integer,
so the constant family of coordinate values $1$ would pass that test.
The space is also generally larger than $\K\otimes_\kk H$, which contains only
series whose coefficient vectors lie in one finite-dimensional subspace; the
exact completion distinctions are treated in the three-duals report, and no
density of the algebraic scalar extension is assumed here.

\subsection{What becomes of an orthonormal basis}
The constant $e_i$ remain orthonormal and separate all vectors by inner
products. They need not be a valuation-topological Schauder basis.
A nonzero ordinary tail, however small its ordinary Hilbert norm, has
valuation zero. Its size therefore does not tend to zero at all surreal
scales. Basis synthesis in this model is \emph{layered}: ordinary
Hilbert synthesis within each coefficient, then Hahn assembly across
exponents. Parseval holds coefficientwise in the sense of
\eqref{hgeo:eq:twosums}, not as an unrestricted strong sum over $i$.

\subsection{Riesz representation requires more than valuation continuity}
\label{hgeo:sub:riesz}
The three-duals report \cite{RepoDuals} proves a detailed classification;
the spectral report has its own counterexample
(\replabel{ihs:hh:prop:riesz}). Two elementary obstructions are enough
here. Both start from an everywhere-defined but ordinarily discontinuous
linear functional on an infinite-dimensional $H$, obtained from a Hamel
basis.

\emph{An infinitely bounded functional \src{03}.} Let $\ell:H\to\kk$ be such
a functional and set $L\bigl(\sum x_\gamma t^\gamma\bigr)=\sum\ell(x_\gamma)
t^\gamma$. It respects strong sums and satisfies $v(Lx)\ge v(x)$, hence is
valuation continuous. It cannot equal $\ip{x}{y}$ for a Hahn vector $y$:
testing constants forces all nonconstant coefficients of $y$ to vanish and
represents $\ell$ by an ordinary Hilbert vector, a contradiction. For a
fixed $\eta>0$, this functional even satisfies $\abs{Lx}\le t^{-\eta}\norm{x}$.
The infinite bound dominates each individual ordinary coefficient ratio.
Therefore the usual phrase ``bounded functional'' does not restore Riesz
unless the meaning of boundedness is strengthened.

\begin{proposition}[Even strict contractions need not be representable
\src{06}]\label{hgeo:prop:badfunctional}
If $H$ is infinite dimensional and $q=t^\eta$ with $\eta>0$, there is a
$\K$-linear strong functional $L:\HH{H}\to\K$ such that
\[
 \abs{Lx}<c\norm{x}\qquad(x\ne0)
\]
for every positive ordinary real $c$, but $L$ is not representable.
\end{proposition}
\begin{proof}
Choose an orthonormal sequence $(e_n)$ and a linear functional
$f:H\to\kk$ with $f(e_n)=n$, extended algebraically from its span by a Hamel
basis. It is unbounded. Extend $f$ coefficientwise to $\widehat f$, and put
$L=q\widehat f$. Coefficientwise extension is strong, and
$v(Lx)\ge v(x)+\eta>v(\norm x)$, so the displayed inequality follows from
leading exponents, independently of the size of the ordinary coefficient.
If $Lx=\ip{x}{y}$, testing constants $u$ and taking the coefficient at
$\eta$ would give $f(u)=\ip{u}{y_\eta}_H$, a bounded functional, which is
impossible.
\end{proof}

By contrast, a globally Hahn-supported series of ordinary continuous
functionals is represented, coefficient by coefficient, by a unique Hahn
vector. Precisely: a $\K$-linear functional is representable if and only if
it is strong and each coefficient functional $u\mapsto\coeff{\gamma}{Lu}$ on
the ordinary $H$ is continuous. This restricted Riesz theorem is the
three-duals report's \replabel{duals:thm:riesz} (the identification of the
represented dual with $H'((t^\Gamma))$); source~06 reproves it
(its Proposition~5.2) with the Baire support lemma
\cref{hgeo:lem:baire} below. Neither obstruction above is claimed as new.

\subsection{Least field-valued norm bounds}
Every ordinary bounded operator lifts to a constant Hahn operator
(\cref{hgeo:sec:adjoints}), and a norm bound means $C\in\F_{\ge0}$ with
$\norm{\widehat Ax}\le C\norm{x}$ for all $x$.

\begin{theorem}[Least field-valued bound; repository theorem \src{06}]
\label{hgeo:thm:norm}
Let $0\ne A\in\mathcal B(H,G)$. Its constant lift $\widehat A$ has a least
nonnegative $\F$-valued norm bound if and only if $A$ attains its ordinary
operator norm, and then the least bound is $\norm A$. If $\norm A$ is not
attained, a nonnegative finite $C\in\F$ is a bound exactly when
$\st C\ge\norm A$, every positive infinite scalar is a bound, and the set of
bounds has no least element.
\end{theorem}
\begin{proof}
For $H=G$ this is the spectral report's norm-attainment criterion
\replabel{ihs:hh:thm:norm}, whose proof describes all upper bounds exactly as
stated. For the rectangular case \mergetag, apply it to
$T=\bigl(\begin{smallmatrix}0&0\\A&0\end{smallmatrix}\bigr)$ on $H\oplus G$:
$\norm T=\norm A$, $T$ attains its norm exactly when $A$ does, and
$\norm{\widehat T(x,y)}=\norm{\widehat Ax}$ with
$\norm{x}\le\norm{(x,y)}$, so $\widehat T$ and $\widehat A$ have the same
bounds (test $y=0$ for the converse).
\end{proof}

The zero operator has least nonnegative bound zero and is excluded from the
nonzero criterion. For example, $Ae_n=(1-1/(n+1))e_n$ is positive,
self-adjoint and adjointable after lifting, but has no least $\F$-valued
operator-norm bound (the spectral report's \replabel{ihs:hh:ex:noleast}).
Thus the order bound of \cref{hgeo:cor:vop} is not a construction of an
$\F$-valued Banach $C^*$-algebra norm \src{06}.

\section{The intrinsic algebra of adjointable operators}
\label{hgeo:sec:adjoints}

\subsection{Operators with a common Hahn support}
For ordinary Hilbert spaces $H,G$, put
\[
 \BB{H}{G}=\mathcal B(H,G)((t^\Gamma)).
\]
An element $A=\sum A_\gamma t^\gamma$ has one well-ordered support of
nonzero ordinary bounded operators. It acts on all Hahn vectors by
\[
 (Ax)_\delta=\sum_{\gamma+\alpha=\delta}A_\gamma x_\alpha.
\]
The support lemma makes these coefficient sums finite. Composition is
Hahn convolution and
$A^*=\sum A_\gamma^*t^\gamma$. There is no requirement that the numbers
$\norm{A_\gamma}$ be uniformly bounded as $\gamma$ varies, and no growth
condition. Write $v(A)=\min\supp A$ for $A\ne0$, and $v(0)=+\infty$.
Then
\begin{equation}\label{hgeo:eq:opvaluation}
 v(Ax)\ge v(A)+v(x),\qquad v(AB)\ge v(A)+v(B).
\end{equation}
An operator is called \emph{integral} if its valuation is nonnegative.
An integral operator with an integral two-sided inverse preserves vector
valuations exactly. Multiplication on either side by such an operator
also preserves the valuation of any nonzero operator: apply
\eqref{hgeo:eq:opvaluation} first to the operator and then to its inverse.

The apparent restriction to globally supported series is intrinsic to
adjoints, rather than an optional technical convenience. The key step is a
Baire lemma, isolated by source~06; it is Step~2 of the spectral report's
proof of \replabel{ihs:hh:thm:adjoint}.

\begin{lemma}[A Baire support lemma \src{06}]\label{hgeo:lem:baire}
Let $(A_\gamma)_{\gamma\in J}$ be bounded operators from an ordinary
Hilbert space $H$ to another, where $J$ is a subset of an ordered set. If
$\{\gamma:A_\gamma u\ne0\}$ is well ordered for every $u\in H$, then
$\{\gamma:A_\gamma\ne0\}$ is well ordered.
\end{lemma}
\begin{proof}
Otherwise dependent choice gives a strictly decreasing sequence
$\gamma_1>\gamma_2>\cdots$ of indices of nonzero operators. Each
$H\setminus\ker A_{\gamma_n}$ is open and dense, because a proper closed
subspace has empty interior. Baire's theorem gives a $u$ in their
intersection, contradicting the hypothesis on its support.
\end{proof}

\begin{theorem}[Automatic adjoint structure; repository input
\src{03, 06}]\label{hgeo:thm:adjoint}
Let $A:\HH{H}\to\HH{G}$ and $B:\HH{G}\to\HH{H}$ be everywhere-defined
$\K$-linear maps satisfying
\[
 \ip{Ax}{y}=\ip{x}{By}\qquad(x\in\HH{H},\ y\in\HH{G}).
\]
There is a unique member of $\BB{H}{G}$ producing the given map, and
$B=A^*$ coefficientwise. Conversely every member of $\BB{H}{G}$ is
adjointable. In particular every everywhere-defined adjointable map is
strong, valuation continuous, and order bounded. No continuity, strongness
or common-support assumption is needed in the forward direction.
\end{theorem}
\begin{proof}
The square case $H=G$ is \replabel{ihs:hh:thm:adjoint} of the spectral
report, and both sources reprove it with the same three steps (bounded
coefficients by the closed graph theorem, one well-ordered support by
\cref{hgeo:lem:baire}, identification on all vectors by testing against
constants); they are not reprinted. The rectangular case reduces to the
square one \mergetag. On $\HH{H\oplus G}=\HH{H}\oplus\HH{G}$ put
$C(x,y)=(By,Ax)$. Then
\[
 \ip{C(x,y)}{(x',y')}=\ip{By}{x'}+\ip{Ax}{y'}
 =\ip{y}{Ax'}+\ip{x}{By'}=\ip{(x,y)}{C(x',y')},
\]
so $C$ is an everywhere-defined self-adjoint map, and the square theorem
gives $C=\sum_\gamma C_\gamma t^\gamma$ with $C_\gamma=C_\gamma^*$ bounded on
$H\oplus G$ and well-ordered support. The lower-left blocks $A_\gamma$ of the
$C_\gamma$ give the expansion of $A$, their upper-right blocks give that of
$B$, and $C_\gamma=C_\gamma^*$ identifies the latter with $A_\gamma^*$.
Uniqueness follows by evaluating on constants and extracting coefficients.
The converse is coefficient comparison. Strongness, continuity and the order
bound are \replabel{ihs:hh:cor:autobounded} applied to $C$, since
$\norm{Ax}=\norm{C(x,0)}$. For real $H,G$ the same proof applies verbatim, as
both sources observe.
\end{proof}

\begin{corollary}[Valuation and automatic boundedness
\src{03, 06}]\label{hgeo:cor:vop}
For $0\ne A\in\BB{H}{G}$,
\[
 v(Ax)\ge v(A)+v(x),\qquad
 v(A^*)=v(A),\qquad
 v(A^*A)=2v(A).
\]
If $s=v(A)$ and $c$ is an ordinary real with $c>\norm{A_s}$, then
$\norm{Ax}\le ct^s\norm{x}$ for all $x$.
\end{corollary}
\begin{proof}
The first two assertions are the support bound and the coefficientwise
adjoint. The leading coefficient of $A^*A$ is $A_s^*A_s\ne0$. The bound is
\replabel{ihs:hh:cor:autobounded}: for nonzero $x$ with leading term
$t^\delta u$, the coefficient at $2(s+\delta)$ of
$c^2t^{2s}\norm{x}^2-\norm{Ax}^2$ is $c^2\norm{u}^2-\norm{A_su}^2>0$, also
when $A_su=0$, since then the leading output moves to a higher exponent.
\end{proof}

\begin{remark}
The theorem uses ordinary Hilbert completeness in its coefficient
closed-graph and Baire arguments. Spherical completeness of the Hahn
space alone does not supply those arguments for arbitrary coefficient
vector spaces. Nor does the theorem classify all valuation-continuous
maps: the three-duals report gives strictly larger classes.
\end{remark}

\subsection{Positive-order operator calculus}\label{hgeo:sub:calculus}
If $B\in\Bpos{H}{H}$, then
\begin{equation}\label{hgeo:eq:formalcalculus}
 (I+B)^{-1}=\sum_{n\ge0}(-B)^n,
 \quad (I+B)^c=\sum_{n\ge0}\binom cn B^n,
 \quad e^B=\sum_{n\ge0}\frac{B^n}{n!},
 \quad \log(I+B)=\sum_{n\ge1}\frac{(-1)^{n+1}B^n}{n}
\end{equation}
($c\in\R$) are legitimate strong Hahn operator series \src{06}. Their support
lies in the additive monoid of $\supp B$. At a fixed exponent there are
finitely many words, so the usual one-variable formal identities survive
substitution, even in the noncommutative coefficient algebra. No two
noncommuting variables are ever substituted into an identity that presupposes
commutativity. If $B=B^*$, then $I+B$ and $(I+B)^{1/2}$ are self-adjoint and
positive definite, each $I$ plus positive order; for $x$ with leading
coefficient $u$ at exponent $\delta$, the quadratic form of either operator
has leading coefficient $\norm u^2$ at exponent $2\delta$. This is
\emph{formal} smallness, not a real operator-norm estimate.

\subsection{Finite-rank operators and self-adjoint inverses}
For fixed Hahn vectors $u,w$, the map $x\mapsto\ip{x}{w}u$ is
adjointable, with adjoint $y\mapsto\ip{y}{u}w$. Expanding its two
vector supports gives a Hahn series of bounded finite-rank coefficient
operators. Finite sums of these maps include every finite-dimensional
orthogonal projection constructed in \cref{hgeo:prop:finite}.

If a self-adjoint everywhere-defined map $B$ is bijective, then its
algebraic inverse is self-adjoint: writing $x=Bu,y=Bv$ gives
\[
 \ip{B^{-1}x}{y}=\ip{u}{Bv}=\ip{Bu}{v}
 =\ip{x}{B^{-1}y}.
\]
\Cref{hgeo:thm:adjoint} then puts $B^{-1}$ in the Hahn operator
algebra. This specific observation is used below, rather than an unproved
general bounded-inverse theorem over $\K$ \src{03}.

\section{Orthogonal splitting and its residue normal form}
\label{hgeo:sec:splitting}

\subsection{Three equivalent formulations}

\begin{definition}\label{hgeo:def:split}
A $\K$-linear subspace $M\subseteq\cH$ is \emph{orthogonally split}, or
\emph{orthogonally complemented}, if $\cH=M\oplus M^\perp$, where
$M^\perp=\{x:\ip{x}{m}=0\text{ for every }m\in M\}$.
The word ``split'' below always means this orthogonal property. $M$ is
\emph{orthoclosed} if $M=M^{\perp\perp}$.
\end{definition}

\begin{proposition}[Projection and best approximation \src{03}]
\label{hgeo:prop:projection}
For a subspace $M\subseteq\cH$, the following are equivalent:
\begin{enumerate}[label=\textup{(\roman*)}]
\item $M$ is orthogonally split;
\item $M$ is the range of an everywhere-defined self-adjoint idempotent;
\item every $x\in\cH$ has a nearest point in $M$, with distances
compared in the ordered field $\F$.
\end{enumerate}
The projection and all nearest points are unique. A split subspace is
valuation closed, orthoclosed, and closed under all strong sums of its
vectors.
\end{proposition}
\begin{proof}
A decomposition $x=m+n$ with $m\in M,n\in M^\perp$ is unique by
positivity. Sending $x$ to $m$ defines a linear self-adjoint idempotent.
Conversely the range and kernel of a self-adjoint idempotent are
orthogonal, and its kernel equals the orthogonal complement of its
range.

For such a decomposition and any $m'\in M$, Pythagoras gives
\[
 \norm{x-m'}^2=\norm{n}^2+\norm{m-m'}^2.
\]
Thus $m$ is the unique nearest point. Conversely, let $m$ be nearest
to $x$ and put $r=x-m$. If $n\in M$ and $\ip{r}{n}\ne0$, then $n\ne0$
and the scalar $a=\ip{r}{n}/\ip{n}{n}$ gives
\[
 \norm{r-an}^2=\norm{r}^2-
          \frac{\abs{\ip{r}{n}}^2}{\norm{n}^2}<\norm{r}^2,
\]
contradicting minimality. Hence $r\in M^\perp$ and the splitting
follows. Finally $M=\ker(I-P)$ for its projection $P$; automatic
adjoint structure makes $P$ continuous and strong. These two
properties prove the two closure assertions, and $M^{\perp\perp}=
\ran P$ because $M^\perp=\ran(I-P)$.
\end{proof}

\subsection{Every projection has ordinary residue}

\begin{lemma}[Integral projections \src{03, 06}]\label{hgeo:lem:integralproj}
Every projection $P$ in $\BH$ is integral. If $P\ne0$,
then $v(P)=0$. Its residue $P_0$ is an ordinary orthogonal projection,
and $P-P_0$ has strictly positive order or is zero.
\end{lemma}
\begin{proof}[Proof \src{03}]
Let $s=v(P)$ when $P\ne0$. Its nonzero leading coefficient $P_s$ is
ordinary self-adjoint, so $P_s^2\ne0$: otherwise
$\norm{P_su}^2=\ip{P_s^2u}{u}=0$ for all $u$.
Consequently $v(P^2)=2s$. The identity $P^2=P$ forces $2s=s$, hence
$s=0$. Taking residues gives $P_0^2=P_0=P_0^*$.
\end{proof}
\begin{proof}[Second route to $v(P)\ge0$ \src{06}]
Pythagoras gives $\norm{Px}\le\norm{x}$ for all $x$. If $s=v(P)$ were
negative, a constant $u$ with $P_su\ne0$ would give $v(Pu)=s<0=v(u)$, so
$\norm{Pu}$ would exceed $\norm u$ by an infinite factor, contradicting this
inequality.
\end{proof}

For any subspace define its \emph{integral residue space} by
\begin{equation}\label{hgeo:eq:residue}
 \res M=\{x_0:x\in M,\ v(x)\ge0\}\subseteq H.
\end{equation}
This is an ordinary linear subspace, without a presumed ordinary closedness
property.

\begin{corollary}[Closed-residue obstruction \src{03, 06}]\label{hgeo:cor:residue}
If $M$ is split with projection $P$, then
\[
 \res M=\ran P_0.
\]
In particular its residue space is ordinarily closed in $H$.
\end{corollary}
\begin{proof}
For integral $x=Px$ one has $x_0=P_0x_0$. Conversely, for constant
$u\in H$, the vector $Pu$ is integral, belongs to $M$, and has residue
$P_0u$. The range of an ordinary orthogonal projection is closed.
\end{proof}

Ordinary closedness in this corollary is a substantially stronger
restriction than valuation closedness of $M$. It detects the
graph obstructions of \cref{hgeo:sec:graphs,hgeo:sec:lattice} without an
appeal to spectral theory.

\subsection{A global normal form using the classical direct rotation}

\begin{theorem}[Residue normal form \src{03, 06}]\label{hgeo:thm:normalform}
A subspace $M\subseteq\HH{H}$ is orthogonally split if and only if
there are an ordinary closed subspace $M_0\subseteq H$ and a Hahn
unitary $U\in\BH$ satisfying $v(U-I)>0$ or $U=I$, such that
\begin{equation}\label{hgeo:eq:normalform}
 M=U\HH{M_0}.
\end{equation}
Here $M_0=\res M$ is unique. Given the projection $P$ onto $M$ and
$P_0$ onto $M_0$, a canonical choice is
\begin{align}
 W&=PP_0+(I-P)(I-P_0)=I+(P-P_0)(2P_0-I),\label{hgeo:eq:rotationW}\\
 U&=W\bigl(I-(P-P_0)^2\bigr)^{-1/2},\label{hgeo:eq:rotation}
\end{align}
the inverse square root being the formal binomial series about $I$. Then
$UP_0=PU$ and $P=UP_0U^*$. Hence every split subspace is unitarily
isometric to the Hahn lift of an ordinary closed subspace, by a unitary whose
residue is the identity.
\end{theorem}
\begin{proof}
Let $M$ be split with projection $P$. By \cref{hgeo:thm:adjoint}, $P\in\BH$,
and by \cref{hgeo:lem:integralproj}, $P_0$ is an ordinary orthogonal
projection with $v(P-P_0)>0$ (including $P=P_0$). These are exactly the
hypotheses of the spectral report's support-controlled direct rotation
\replabel{ihs:fr:lem:rotation}, with $P^\circ=P_0$: it gives
$U^*U=UU^*=I$, $UP_0=PU$ and $v(U-I)>0$ for the unitary
\eqref{hgeo:eq:rotation}. Both sources reprove that lemma (by direct
multiplication, $W^*W=WW^*=I-(P-P_0)^2$ and $WP_0=PW$; the square
$(P-P_0)^2$ commutes with both projections, so its binomial inverse square
root commutes with $W,W^*,P,P_0$); the second expression for $W$ in
\eqref{hgeo:eq:rotationW} is source~06's. It follows that $P=UP_0U^*$ and,
since the constant projection $P_0$ acts coefficientwise,
$\ran P_0=\HH{M_0}$ with $M_0=\ran P_0$; so \eqref{hgeo:eq:normalform} holds.
\Cref{hgeo:cor:residue} gives $M_0=\res M$, proving uniqueness. Conversely,
the operator $UP_0U^*$ is an orthogonal projection onto the subspace in
\eqref{hgeo:eq:normalform}.
\end{proof}

The formula \eqref{hgeo:eq:rotation} is the classical direct rotation of
projections, the Kato--Nagy two-projection intertwiner; see Simon's accounts
and attributions to earlier work \cite{Simon,SimonKato} and Kato
\cite{Kato}. The spectral report uses it for isolated finite spectral
clusters. The point here is that automatic adjoints and integral
projections place \emph{every} split Hahn subspace in its domain of formal
applicability. Neither the formula nor the rotation lemma is a new result;
the unitary \eqref{hgeo:eq:rotation} is canonical as a formula but is not the
only near-identity unitary conjugating $P_0$ to $P$.

\subsection{Exact graph coordinates for a fixed residue}

\begin{corollary}[The positive-order graph chart \src{03, 06}]
\label{hgeo:cor:graphchart}
Fix an ordinary closed $M_0\subseteq H$ and decompose
$H=M_0\oplus M_0^\perp$. The split subspaces with residue $M_0$ are in
bijection with
\[
 X\in\BB{M_0}{M_0^\perp},\qquad X=0\text{ or }v(X)>0,
\]
by $X\mapsto\Graph X=\{(z,Xz):z\in\HH{M_0}\}$. For $J=(I,X)^{\mathsf T}$ the
projection is
\begin{equation}\label{hgeo:eq:graphproj}
 P_X=J(I+X^*X)^{-1}J^*
 =\begin{pmatrix}R_X&R_XX^*\\XR_X&XR_XX^*\end{pmatrix},
 \qquad R_X=(I+X^*X)^{-1}.
\end{equation}
The operator $X$ is uniquely determined by the subspace.
\end{corollary}
\begin{proof}
Take the near-identity unitary $U$ in \cref{hgeo:thm:normalform}
and write its blocks relative to the displayed decomposition.
The block $U_{00}=I+\text{positive-order terms}$ is invertible by
\eqref{hgeo:eq:neumann}. Its image of $M_0$ is therefore the graph of
$X=U_{10}U_{00}^{-1}$, a positive-order operator with bounded
coefficients. A graph determines its map uniquely.

Conversely $I+X^*X$ has a Neumann inverse of residue $I$. Formula
\eqref{hgeo:eq:graphproj} is self-adjoint and idempotent, since
$J^*J=I+X^*X$; its range is exactly $\Graph X$. Its residue is the
projection onto the first summand. All products and inverses have
support in the positive monoid generated by the supports of $X,X^*$.
\end{proof}

One may read the bijection as an exact formal chart on the Grassmannian of
split subspaces; no differential or manifold structure is asserted by that
terminology \src{06}. This is a classification of all split subspaces, not
of all valuation-closed subspaces.

\begin{proposition}[Closed residue is not sufficient \src{06}, completed
\mergetag]\label{hgeo:prop:closedresidue}
Let $H$ be infinite dimensional, let $f:H\to\kk$ be a discontinuous linear
functional, let $\widehat f$ be its coefficientwise extension, and put
\[
 M=\{(x,q\widehat f(x)):x\in\HH{H}\}\subseteq\HH{H\oplus\kk}.
\]
Then $M$ is valuation closed and closed under strong sums, its residue
$\res M=H\oplus0$ is closed, but $M^\perp=0$ and $M\ne\HH{H\oplus\kk}$. So
$M$ is neither split nor orthoclosed.
\end{proposition}
\begin{proof}
The map $\Phi(x,y)=y-q\widehat f(x)$ is linear, strong and satisfies
$v(\Phi(x,y))\ge\min\{v(y),v(x)+\eta\}$, so it is valuation continuous and
$M=\ker\Phi$ is closed and strong-sum closed. An integral $(x,q\widehat f(x))$
has $v(x)\ge0$, and its second component has valuation at least
$v(x)+\eta>0$; its residue is $(x_0,0)$, and every $(u,0)$, $u\in H$, is the
residue of the constant choice $x=u$. So $\res M=H\oplus0$. If $(w,c)\perp M$,
pairing with $(u,qf(u))$ for constant $u$ and extracting the coefficient at
$\gamma+\eta$ gives
\[
 \ip{u}{w_{\gamma+\eta}}_H+f(u)\,\overline{c_\gamma}=0\qquad(u\in H).
\]
If some $c_\gamma\ne0$, then $f(u)=-\ip{u}{w_{\gamma+\eta}}_H/
\overline{c_\gamma}$ would be continuous; so $c=0$, and then every
coefficient of $w$ is orthogonal to $H$, so $w=0$. Finally $(0,1)\notin M$.
A split subspace with zero complement would be everything, and
$M^{\perp\perp}=\HH{H\oplus\kk}\ne M$.
\end{proof}

Source~06 remarks that ``closed residue alone is insufficient'', pointing to
coefficient maps such as those of \cref{hgeo:prop:badfunctional}; the
proposition supplies the details. It settles negatively the case left open
by source~03, which makes no claim that an arbitrary valuation-closed
subspace with closed residue automatically has a positive-order graph
chart: this one has closed residue $H\oplus0$ and no chart. The
bounded-coefficient hypothesis in \cref{hgeo:cor:graphchart} is part of the
classification, not a disposable regularity condition. The subspace $M$ is
not orthoclosed; whether every \emph{orthoclosed} subspace with closed
residue splits is not decided here (\cref{hgeo:q:orthoclosed}).

\section{Unitary structure and finite simultaneous straightening}
\label{hgeo:sec:unitary}

\begin{corollary}[Unitary orbit classification \src{06}]\label{hgeo:cor:orbits}
Two projections on $\HH{H}$ are unitarily conjugate if and only if their
ordinary residue projections have the same range Hilbert dimension and the
same kernel Hilbert dimension. These dimensions are cardinalities of ordinary
orthonormal bases, not Hamel dimensions.
\end{corollary}
\begin{proof}
A unitary and its inverse are contractions, so by the argument of the second
route in \cref{hgeo:lem:integralproj} their supports are nonnegative. Taking
residues in $U^*U=UU^*=I$ gives an ordinary unitary $U_0$, and $UPU^*=P'$
gives $U_0P_0U_0^*=P'_0$, which preserves the two dimensions. Conversely,
equality of the two dimension pairs gives an ordinary unitary between the
residue decompositions. Compose its constant lift with the two straightening
unitaries of \cref{hgeo:thm:normalform}.
\end{proof}

\begin{theorem}[A split unitary group and its infinitesimal logarithm
\src{06}]\label{hgeo:thm:unitary}
Every adjointable unitary on $\HH{H}$ has a unique factorization
\[
 U=\widehat{U_0}\exp X,\qquad U_0\in\Un(H),\quad
 X\in\Bpos{H}{H},\quad X^*=-X.
\]
Consequently residue gives a split exact sequence of groups
\[
 1\longrightarrow\Un_{1,\Gamma}(H)
 \longrightarrow\Un(\HH{H})
 \overset{\res}{\longrightarrow}\Un(H)\longrightarrow1,
\]
where $\Un_{1,\Gamma}(H)=\{U\in\Un(\HH H):U-I\in\Bpos{H}{H}\}$. Exponential
and logarithm are inverse bijections between the positive-order skew-adjoint
series and $\Un_{1,\Gamma}(H)$.
\end{theorem}
\begin{proof}
As in \cref{hgeo:cor:orbits}, $U$ and $U^*$ have nonnegative supports, and
$U_0^*U_0=U_0U_0^*=I$. Put $V=\widehat{U_0}^*U=I+Y$ with $Y$ of positive
order, and set $X=\log V$. Since $V$ is unitary, formal one-variable
identities give
\[
 X^*=\log V^*=\log V^{-1}=-\log V=-X.
\]
Conversely $X^*=-X$ implies $(\exp X)^*\exp X=\exp(-X)\exp X=I$.
Uniqueness follows from the residue and the formal inverse identities for
exponential and logarithm. Constant lifting is a group-theoretic section.
\end{proof}

The exponential parameter is not an additive group parameter:
multiplication of noncommuting infinitesimal unitaries corresponds to the
formal Baker--Campbell--Hausdorff law. No real norm convergence of that law
is asserted. General formal exponential correspondences are established
material \cite{Bagayoko}; the statement above specifies the exact unitary
group in the present Hilbert setting.

\begin{theorem}[Finite orthogonal decompositions straighten simultaneously
\src{06}]\label{hgeo:thm:finite}
Let $P_1,\ldots,P_r$ be pairwise orthogonal projections with
$\sum_{j=1}^rP_j=I$, and put $Q_j=\widehat{(P_j)_0}$. Then a single
near-identity adjointable unitary $U$ satisfies
\[
 UQ_jU^*=P_j\qquad(1\le j\le r).
\]
An explicit choice is
\[
 W=\sum_{j=1}^rP_jQ_j,\qquad U=W(W^*W)^{-1/2}.
\]
The same conclusion holds for a finite orthogonal family not summing to
$I$, after adjoining its complementary projection.
\end{theorem}
\begin{proof}
The $Q_j$ form an ordinary orthogonal partition of $I$, and $W_0=I$.
Hence $W$ is formally invertible and $W^*W=I+\text{positive order}$.
Orthogonality gives $WQ_j=P_jQ_j=P_jW$. Taking adjoints,
$Q_jW^*=W^*P_j$, so $Q_jW^*W=W^*P_jW=W^*WQ_j$: the operator $W^*W$ commutes
with each $Q_j$, and so does its inverse square root. Therefore
$UQ_j=P_jU$. Moreover $U^*U=I$, and invertibility of $U$ implies $UU^*=I$.
Its residue is $I$.
\end{proof}

For $r=2$ this is \cref{hgeo:thm:normalform} with a different unitary. The
finite qualification is essential: a countable family may fail to have even
an orthogonal join, and countable simultaneous straightening can fail for two
independent reasons (\cref{hgeo:sub:countable}). Source~07 gives the exact
infinite criterion (\cref{hgeo:us:thm:main}), of which this theorem is the
finite case (\cref{hgeo:us:rem:finite}).

\section{Amplified and infinitesimal graphs}
\label{hgeo:sec:graphs}

\subsection{The normal equation is the whole obstruction}

\begin{proposition}[General adjointable graph criterion \src{03}]
\label{hgeo:prop:graphcriterion}
Let $A\in\BB{H}{G}$. Its graph is valuation closed and closed under
strong Hahn sums. It is orthogonally split if and only if
$I+A^*A$ is surjective on $\HH{H}$. In that case this operator is
invertible in $\BB{H}{H}$ and the graph projection is
\[
 P=\begin{pmatrix}I\\A\end{pmatrix}
       (I+A^*A)^{-1}\begin{pmatrix}I&A^*\end{pmatrix}.
\]
\end{proposition}
\begin{proof}
The graph is the kernel of the continuous strong map
$(x,y)\mapsto y-Ax$, which gives both closure properties.
The operator $B=I+A^*A$ is injective because
\[
 \ip{Bx}{x}=\norm{x}^2+\norm{Ax}^2>0\quad(x\ne0).
\]
If the graph is split, project $(f,0)$ onto a point $(x,Ax)$.
Orthogonality of the residual to every $(u,Au)$ gives $Bx=f$.
As $f$ was arbitrary, $B$ is onto. Conversely, if $B$ is onto it is
bijective and self-adjoint, so the observation at the end of
\cref{hgeo:sec:adjoints} gives $B^{-1}$ in the adjointable algebra.
Writing $Jx=(x,Ax)$, the identity $J^*J=B$ verifies that
$JB^{-1}J^*$ is an orthogonal projection with range $\Graph A$.
\end{proof}

Positivity of $I+A^*A$ does not alone imply its surjectivity. This is
exactly the step in an ordinary Hilbert-space graph argument that
cannot be imported without further hypotheses.

\begin{lemma}[Graphs of integral operators split \src{03, 06}]
\label{hgeo:lem:finitegraph}
If $B\in\BB{V}{W}$ is integral, then its graph is split in $\HH{V\oplus W}$,
and the residue of its projection is the ordinary projection onto
$\Graph B_0$.
\end{lemma}
\begin{proof}
The constant coefficient of $I+B^*B$ is $I+B_0^*B_0$, an ordinary bounded
positive operator bounded below by $I$, hence boundedly invertible.
Factoring out this constant reduces the inverse of $I+B^*B$ to a
positive-order Neumann series, and \cref{hgeo:prop:graphcriterion} applies.
Taking residues in the projection formula gives the last assertion.
\end{proof}

\subsection{Finite and infinite scalar amplification of bounded operators}

\begin{theorem}[Exact amplification threshold \src{03}]\label{hgeo:thm:threshold}
Let $T:H\to G$ be an ordinary bounded operator and $a\in\K$.
Then
\begin{equation}\label{hgeo:eq:threshold}
 \Graph(aT)\text{ is split}
 \quad\Longleftrightarrow\quad
 a=0\text{ or }v(a)\ge0\text{ or }\ran T\text{ is ordinarily closed}.
\end{equation}
For nonzero infinite $a$, one has the exact residue formula
\begin{equation}\label{hgeo:eq:graphresidue}
 \res\Graph(aT)=\ker T\oplus\ran T\subseteq H\oplus G.
\end{equation}
When the graph is split, its projection has residue the ordinary
projection onto $\Graph(a_0T)$ in the finite case, and onto
$\ker T\oplus\ran T$ in the infinite case.
\end{theorem}
\begin{proof}
If $a=0$ the graph is horizontal. If $v(a)\ge0$, then $aT$ is integral with
residue $a_0T$, $a_0=[t^0]a$, and \cref{hgeo:lem:finitegraph} proves
splitting and the finite-case residue assertion.

Now let $v(a)<0$. Set $b=a^{-1}$ and $\eta'=v(b)>0$, with
$b_{\eta'}\ne0$. A graph vector $(x,y)$ is integral exactly when both
components are integral. It satisfies $Tx=by$. The constant
coefficient gives $Tx_0=0$. The coefficient at $\eta'$ gives
\[
 T x_{\eta'}=b_{\eta'} y_0,
\]
because $y$ has no negative exponents and $b$ has least exponent
$\eta'$. Hence $(x_0,y_0)\in\ker T\oplus\ran T$.
Conversely, for any $k\in\ker T,u\in H$,
\[
 x=k+bu,\qquad y=Tu
\]
is an integral graph vector with residue $(k,Tu)$. This proves
\eqref{hgeo:eq:graphresidue}. If the graph is split,
\cref{hgeo:cor:residue} makes this residue space ordinarily
closed, forcing $\ran T$ to be closed in $G$.

For the converse assume $\ran T$ closed. Write
\[
 H=N\oplus X,\quad N=\ker T,\qquad
 G=Y\oplus Y^\perp,\quad Y=\ran T.
\]
The ordinary operator $T_1=T|_X:X\to Y$ is a bounded isomorphism,
by the ordinary bounded inverse theorem. Let $C_1=T_1^*T_1$ on $X$;
it is boundedly invertible. On $N$, $B=I+\bar a a T^*T$ equals $I$.
On $X$ an inverse is
\begin{equation}\label{hgeo:eq:infinitegraphinverse}
 B_X^{-1}=(\bar a a)^{-1}C_1^{-1}
      \bigl(I+(\bar a a)^{-1}C_1^{-1}\bigr)^{-1}.
\end{equation}
The inner correction has positive valuation $2\eta'$, so its inverse
is Hahn defined. This supplies a two-sided inverse to $B$ and proves
splitting. The residue projection is the ordinary projection onto
\eqref{hgeo:eq:graphresidue}, by \cref{hgeo:cor:residue}.
\end{proof}

\begin{remark}[Geometric interpretation \src{03}]
At a finite nonzero residue $a_0$, the leading geometry is the ordinary
sloped graph. If $v(a)>0$, it is horizontal. At an infinite scale and
with closed range, the nonkernel portion of the graph becomes vertical:
its residue is $\ker T\oplus\ran T$. A nonclosed ordinary range
cannot be the residue of an orthogonal projection, so precisely this
case obstructs the transition. The threshold depends on the sign of
$v(a)$, not on a metric magnitude or a discrete value-group step. A
failure of ordinary closed range becomes a failure of surreal best
approximation after \emph{any} infinite amplification, however small its
negative valuation.
\end{remark}

\begin{example}[An explicit complete graph with no nearest point \src{03}]
\label{hgeo:ex:badgraph}
Let $H=G=\ell^2(\N_{\ge1})$, $T=\diag(1/n)$ and $A=q^{-1}T$.
The range of $T$ is not ordinarily closed: it contains all
finite-support vectors but does not contain $h=(1/n)_{n\ge1}$.
Indeed $Tx=h$ would require $x=(1,1,\ldots)\notin\ell^2$.
The graph of $A$ is therefore not split.

There is also a direct witness to failure of best approximation.
The vector $z=(h,0)$ has no nearest point in the graph. Such a point
$(x,Ax)$ would solve the normal equation
\begin{equation}\label{hgeo:eq:badnormal}
 (q^2I+T^2)x=q^2h.
\end{equation}
Because $T^2$ is injective, the left side has valuation $v(x)$ for
nonzero $x$. Thus $v(x)=2\eta$. Its coefficient there would satisfy
$T^2x_{2\eta}=h$, requiring
$x_{2\eta}=(n)_{n\ge1}$, again not a Hilbert vector.
The contradiction proves the assertion without a compact-operator
spectral theorem.

This graph is valuation closed, strong-sum closed, and complete as a
closed subspace of the complete valued space in \cref{hgeo:thm:complete}.
These properties do not supply a nearest point. Source~03 asserts failure of
an attained minimum here and needs no assertion about a distance infimum;
for the flipped point $(0,h)$ the distances have no infimum at all
(\cref{hgeo:prop:distancecut}).
\end{example}

\begin{corollary}[Enlarging the scale group cannot repair this graph \src{03}]
\label{hgeo:cor:norepair}
Let $\Gamma\hookrightarrow\Delta$ be an ordered group embedding,
and extend coefficients and exponents to the $\Delta$ Hahn spaces.
For ordinary $T$ and $v(a)<0$, the extended graph is split exactly
when the original one is split, namely when $\ran T$ is ordinarily
closed. This holds also for noncofinal embeddings of value groups.
\end{corollary}
\begin{proof}
The image of $a$ still has negative valuation. \Cref{hgeo:thm:threshold}
applies over both groups with the same ordinary range condition.
\end{proof}

\subsection{Infinitesimal graphs of closed unbounded operators}
Let $H,G$ be ordinary Hilbert spaces and let $A:\Dom A\subseteq H\to G$ be a
closed densely defined ordinary operator. The domain is given its inherited
Hilbert inner product, not its graph norm, when forming $\Dom A((t^\Gamma))$.
Write $A_\Gamma$ for coefficientwise application on that space and put
\begin{equation}\label{hgeo:eq:unboundedgraph}
 M_q(A)=\{(x,qA_\Gamma x):
                 x\in\Dom A((t^\Gamma))\}
 \subseteq\HH{H\oplus G}.
\end{equation}
Although $A$ can be unbounded in ordinary norm, $A_\Gamma$ does not lower
valuation on its domain. Thus the graph is infinitesimally close to the
horizontal domain in a valuation sense. It is not the graph of an
everywhere-defined bounded-coefficient Hahn operator unless $A$ is bounded.

\begin{theorem}[Exact infinitesimal-graph criterion \src{06}]
\label{hgeo:thm:unboundedgraph}
For every closed densely defined $A$ and every $\eta\in\Gamma_{>0}$, with
$q=t^\eta$:
\begin{enumerate}[label=\textup{(\roman*)}]
\item $M_q(A)$ is a strong $\K$-linear subspace, valuation closed and
spherically complete.
\item Its orthogonal complement is
\begin{equation}\label{hgeo:eq:graphperp}
 M_q(A)^\perp
 =\{(-qA_\Gamma^*y,y):y\in\Dom A^*((t^\Gamma))\},
\end{equation}
where $A_\Gamma^*$ means coefficientwise application of the ordinary
unbounded adjoint $A^*$. In particular $M_q(A)$ is orthoclosed.
\item Its residue is exactly $\res M_q(A)=\Dom A\oplus\{0\}$.
\item It is orthogonally split if and only if $\Dom A=H$,
equivalently if and only if the closed operator $A$ is bounded.
\end{enumerate}
\end{theorem}
\begin{proof}
\emph{Strong linearity and completeness.} Coefficientwise application of
any linear map preserves strong sums on its domain: supports only shrink,
and each coefficient sum is finite. Scalar convolution is likewise
compatible. For $x\ne0$,
\begin{equation}\label{hgeo:eq:graphisometry}
 v\bigl((x,qA_\Gamma x)\bigr)=v(x),
\end{equation}
because the second component has valuation at least $v(x)+\eta$.
The graph parametrization is therefore a valuation isometry from
$\Dom A((t^\Gamma))$, which is spherically complete by
\cref{hgeo:thm:complete}, even though $\Dom A$ is usually
incomplete in its inherited ordinary norm.

\emph{Orthogonality.} If $(u,w)$ is orthogonal to the graph, testing a
constant $k\in\Dom A$ and extracting exponent $\gamma$ gives
\[
 \ip{k}{u_\gamma}_H+\ip{Ak}{w_{\gamma-\eta}}_G=0
 \qquad(k\in\Dom A).
\]
By the definition of the ordinary densely defined adjoint, this says
$w_{\gamma-\eta}\in\Dom A^*$ and
$A^*w_{\gamma-\eta}=-u_\gamma$. Varying $\gamma$ proves
\eqref{hgeo:eq:graphperp}; the reverse inclusion follows from finite Hahn
coefficient regrouping. Taking a second orthogonal complement, by the same
computation with $-A^*$ in place of $A$ and the coordinates exchanged,
replaces $A$ by $A^{**}$. The ordinary theorem $A^{**}=A$ for closed densely
defined operators \cite{Kato} gives orthoclosedness. Orthogonal complements
are valuation closed because pairings with a fixed vector are continuous;
this proves the remaining closedness assertion in (i).

\emph{Residue and complementation.} By \eqref{hgeo:eq:graphisometry}, an
integral graph vector has $v(x)\ge0$ and second component of positive
valuation. Its residue is $(x_0,0)$ with $x_0\in\Dom A$. Every such residue
is realized by a constant $x$. If the graph were split,
\cref{hgeo:cor:residue} would make $\Dom A$ closed in $H$. It is dense,
so $\Dom A=H$, and the ordinary closed graph theorem makes $A$ bounded.
Conversely, if $A$ is bounded, $q\widehat A$ has positive order and
\cref{hgeo:cor:graphchart} gives its orthogonal graph projection.
\end{proof}

\begin{corollary}[Explicit absence of best approximation \src{06}]
\label{hgeo:cor:best}
If $A$ is unbounded and $h'\in H\setminus\Dom A$, the constant vector
$(h',0)$ has no nearest point in $M_q(A)$.
\end{corollary}
\begin{proof}
By the proof of \cref{hgeo:prop:projection}, a nearest point $m$ in a linear
subspace must leave an orthogonal residual. Suppose $(h',0)=m+n$ with
$m\in M_q(A)$ and $n\in M_q(A)^\perp$. Pythagoras makes both vectors integral.
By \cref{hgeo:thm:unboundedgraph}(iii) and \eqref{hgeo:eq:graphperp}, their
residues lie in $\Dom A\oplus0$ and $0\oplus\Dom A^*$ respectively. Hence
$h'\in\Dom A$, a contradiction.
\end{proof}

The theorem is an explicit mechanism behind the abstract Sol\`er
obstruction (\cref{hgeo:sub:literature}): infinitesimal scaling preserves
the ordinary domain defect in the residue, while valuation completeness
ignores the ordinary metric incompleteness of that domain.

\begin{proposition}[Exact lower cut of the distances \src{06}]
\label{hgeo:prop:distancecut}
Let $H=G=\ell^2(\N_{\ge1})$, $\Lambda=\diag(n)$ with
$\Dom\Lambda=\{x:\sum n^2\abs{x_n}^2<\infty\}=\ran T$, and
$z=(h,0)$ with $h=(1/n)$. The nonnegative lower bounds of
$\{\norm{z-m}:m\in M_q(\Lambda)\}$ are precisely zero and the positive
infinitesimals of $\F$. In particular, no distance infimum exists in $\F$.
\end{proposition}
\begin{proof}
$\Lambda$ is closed, self-adjoint and unbounded ($\norm{\Lambda e_n}=n$), and
$h\notin\Dom\Lambda$. If a graph vector $m$ has negative valuation, then
$z-m$ has negative valuation and infinite norm. If it is integral, its
constant coefficient is $(x_0,0)$ with $x_0\in\Dom\Lambda$, so
$h-x_0\ne0$, and $\norm{z-m}$ has positive ordinary standard part. Thus every
positive infinitesimal is a lower bound. Conversely let $h^{(N)}$ be the
finite truncation of $h$. For $m_N=(h^{(N)},q\Lambda h^{(N)})$,
\[
 \st\norm{z-m_N}=\norm{h-h^{(N)}}_H\longrightarrow0
\]
in the ordinary real sense. No positive bound with positive standard part
can be a lower bound, and no infinite positive bound either. The set of zero
and the positive infinitesimals has no greatest member, since doubling a
positive infinitesimal gives a larger one and $q$ is one. Therefore the
distance set has no infimum.
\end{proof}

The cut is exactly the one of the three-duals report's
\replabel{duals:thm:distance}, proved there for the closed hyperplane
$M_\Gamma$ of \replabel{duals:thm:hyperplane}, with the same argument
(ordinary density of the residue plus valuation zero of every integral
difference). Source~06 does not cite it. The new point of
\cref{hgeo:prop:distancecut} is only that the subspace is now orthoclosed;
the hyperplane has zero orthogonal complement and is not.

\subsection{The common criterion: scaled closed linear relations}
\label{hgeo:sub:relation}
The two graph theorems are related, but neither implies the other.
\Cref{hgeo:thm:threshold} treats bounded $T$ with arbitrary kernel and
range, and also finite scalars; \cref{hgeo:thm:unboundedgraph} treats
unbounded $A$ and also computes the orthogonal complement. After exchanging
the coordinates, both are graphs of a closed linear relation multiplied by a
positive infinitesimal in the second coordinate.

A \emph{closed linear relation} $R\subseteq H\oplus G$ is a closed linear
subspace. Its domain and multivalued part are
\[
 \dom R=\{x:(x,y)\in R\text{ for some }y\},\qquad
 \mul R=\{y:(0,y)\in R\};
\]
$\mul R$ is closed, while $\dom R$ need not be. Coefficientwise,
$R((t^\Gamma))\subseteq\HH{H\oplus G}$ consists of the Hahn vectors all of
whose coefficients lie in $R$. For $b\in\K$ with $v(b)>0$ put
\[
 M_b(R)=\{(x,by):(x,y)\in R((t^\Gamma))\}\subseteq\HH{H\oplus G}.
\]

\begin{theorem}[Scaled closed relations \mergetag]\label{hgeo:thm:relation}
Let $R\subseteq H\oplus G$ be a closed linear relation and $b\in\K$ with
$v(b)>0$. Then
\begin{equation}\label{hgeo:eq:relationresidue}
 \res M_b(R)=\dom R\oplus\mul R,
\end{equation}
and the following are equivalent: $M_b(R)$ is split; $\res M_b(R)$ is
ordinarily closed; $\dom R$ is ordinarily closed. When they hold,
\begin{equation}\label{hgeo:eq:relationsplit}
 M_b(R)=\Graph\bigl(b\widehat{A_R}\bigr)\oplus\bigl(0\oplus\HH{\mul R}\bigr),
\end{equation}
an orthogonal sum, where $A_R:\dom R\to(\mul R)^\perp$ is the bounded
operator part of $R$.
\end{theorem}
\begin{proof}
\emph{Residue.} Let $(x,by)\in M_b(R)$ be integral, with
$(x,y)\in R((t^\Gamma))$. Then $v(x)\ge0$ and $(x_0,y_0)\in R$, so
$x_0\in\dom R$. The residue of the second component is
$\coeff{0}{(by)}=\sum_{\alpha+\beta=0}b_\alpha y_\beta$ with
$\alpha\in\supp b\subseteq\Gamma_{>0}$, so only coefficients $y_\beta$ with
$\beta<0$ occur. For $\beta<0$, $x_\beta=0$, so $(0,y_\beta)\in R$ and
$y_\beta\in\mul R$. Hence $\res M_b(R)\subseteq\dom R\oplus\mul R$.
Conversely, for $u\in\dom R$ choose $(u,y)\in R$; the constant pair gives the
integral vector $(u,by)$ with residue $(u,0)$. For $m\in\mul R$, the vector
$(0,b^{-1}m)$ lies in $R((t^\Gamma))$, since each of its coefficients is a
multiple of $(0,m)$, and gives $(0,m)\in M_b(R)$ with residue $(0,m)$.
The residue space is linear, so it contains $\dom R\oplus\mul R$.

\emph{Necessity.} If $M_b(R)$ is split, \cref{hgeo:cor:residue} makes
$\dom R\oplus\mul R$ closed in $H\oplus G$, and then
$\dom R=\{u:(u,0)\in\dom R\oplus\mul R\}$ is closed in $H$. Conversely,
$\mul R$ is always closed, so $\dom R\oplus\mul R$ is closed whenever
$\dom R$ is.

\emph{Sufficiency.} Let $H_1=\dom R$ be closed and $G_1=(\mul R)^\perp$.
For $u\in H_1$ the set $\{y:(u,y)\in R\}$ is a coset of $\mul R$; let
$A_Ru$ be its unique element in $G_1$. Then $A_R:H_1\to G_1$ is linear, and
its graph is $R\cap(H\oplus G_1)$, which is closed. As $H_1$ is a Hilbert
space, the closed graph theorem makes $A_R$ bounded. Every $(u,y)\in R$ is
$(u,A_Ru)+(0,y-A_Ru)$ with $y-A_Ru\in\mul R$, and this decomposition is given
by bounded ordinary projections, so it holds coefficientwise:
$R((t^\Gamma))=\{(u,\widehat{A_R}u)\}+\bigl(0\oplus\mul R((t^\Gamma))\bigr)$.
Multiplying the second coordinates by $b$, and using that $b$ is invertible,
gives \eqref{hgeo:eq:relationsplit}. The first summand is the graph of the
positive-order operator $b\widehat{A_R}\in\BB{H_1}{G_1}$, which is split in
$\HH{H_1\oplus G_1}$ by \cref{hgeo:cor:graphchart}; extending its projection
by zero on $\HH{(H_1\oplus G_1)^\perp}$ gives a projection on
$\HH{H\oplus G}$. The second summand is the constant split subspace
$\HH{0\oplus\mul R}$, which is orthogonal to $\HH{H\oplus G_1}$. The sum of
two projections with orthogonal ranges is the projection onto the sum
(\cref{hgeo:prop:poset}), so $M_b(R)$ is split.
\end{proof}

\begin{corollary}[Both graph theorems \mergetag]\label{hgeo:cor:bothgraphs}
\begin{enumerate}[label=\textup{(\alph*)}]
\item For bounded $T:H\to G$ and $v(a)<0$, the exchange of coordinates
$(x,y)\mapsto(y,x)$ maps $\Graph(aT)$ onto $M_{a^{-1}}(R)$ with
$R=\{(Tx,x):x\in H\}\subseteq G\oplus H$, for which $\dom R=\ran T$ and
$\mul R=\ker T$. \Cref{hgeo:thm:relation} gives the residue
\eqref{hgeo:eq:graphresidue} and the infinite case of
\cref{hgeo:thm:threshold}.
\item For closed densely defined $A$, $M_q(A)=M_q(\Graph A)$ with
$\dom=\Dom A$ and $\mul=0$; \cref{hgeo:thm:relation} gives
\cref{hgeo:thm:unboundedgraph}(iii) and (iv).
\end{enumerate}
The two cases meet exactly when $T$ is bounded and injective with dense
range and $A=T^{-1}$: then $\Graph(q^{-1}T)$ and $M_q(T^{-1})$ are exchanged
by the coordinate flip. In particular, for $T=\diag(1/n)$ the subspace of
\cref{hgeo:ex:badgraph} is the flip of $M_q(\Lambda)$.
\end{corollary}
\begin{proof}
For (a), $\Graph(aT)$ flipped is $\{(aTx,x)\}$; substituting $x=a^{-1}x'$
gives $\{(Tx',a^{-1}x')\}=M_{a^{-1}}(R)$, because $R((t^\Gamma))=
\{(Tx',x'):x'\in\HH{H}\}$. The flip is a constant unitary, so it preserves
splitting and maps residues by the ordinary flip. For (b), the coefficients
of $(x,A_\Gamma x)$ lie in $\Graph A$ exactly when those of $x$ lie in
$\Dom A$. The last assertion is (a) with $\ker T=0$ and
$\{(Tx,x)\}=\Graph(T^{-1})$.
\end{proof}

The relation theorem does not replace the two source theorems: it gives
neither the finite-scalar clause of \cref{hgeo:thm:threshold} nor the
orthogonal complement and orthoclosedness in
\cref{hgeo:thm:unboundedgraph}(ii).

\section{The projection poset: meets, kernels and countable joins}
\label{hgeo:sec:lattice}

\subsection{What survives: orthocomplements and orthogonal sums}
Let $\Spl(\cH)$ be the collection of split $\K$-linear subspaces,
ordered by inclusion. Its bottom and top are $0$ and $\cH$. Sending a
projection to its range is an order isomorphism from the poset of
projections, ordered by $P\le Q\iff PQ=QP=P$, onto $\Spl(\cH)$ (source~06
works with that poset of projections). Orthocomplementation, $P\mapsto I-P$,
is an order-reversing involution, by the defining orthogonal decompositions.

\begin{proposition}[Orthomodular poset structure \src{03}]\label{hgeo:prop:poset}
If $M,N\in\Spl(\cH)$ and $M\perp N$, then $M\oplus N$ is split
and is their least upper bound in $\Spl(\cH)$. If $M\subseteq N$,
then $N\cap M^\perp$ is split and
\[
 N=M\oplus(N\cap M^\perp).
\]
Consequently $\Spl(\cH)$ is an orthomodular poset.
\end{proposition}
\begin{proof}
Let $P_M,P_N$ be the unique projections. Orthogonality gives
$P_MP_N=P_NP_M=0$, so $P_M+P_N$ is a self-adjoint idempotent with
range $M\oplus N$. Every subspace containing both contains their sum,
which proves the least-upper-bound assertion.

If $M\subseteq N$, then $P_NP_M=P_M$ and, taking adjoints,
$P_MP_N=P_M$. The operator $P_N-P_M$ is a self-adjoint idempotent.
Its range is exactly $N\cap M^\perp$, proving both splitting and the
displayed orthomodular identity. These properties, together with
orthocomplementation and the bounds $0,\cH$, are the asserted
orthomodular-poset axioms. They do not assert that arbitrary binary
meets or joins exist.
\end{proof}

By induction, finite mutually orthogonal families have joins, their sums.
\Cref{hgeo:thm:nojoin} shows that countable ones need not.

\subsection{Two good subspaces with a bad intersection}

\begin{theorem}[Exact finite-dimensional/lattice dichotomy \src{03};
one direction also \src{06}]\label{hgeo:thm:lattice}
Let $\Gamma\ne0$ be any set-sized ordered abelian group. For an
ordinary real or complex Hilbert space $H$,
\[
 \Spl(\HH{H})\text{ is a lattice}
 \quad\Longleftrightarrow\quad \dim H<\infty.
\]
If $H$ is infinite dimensional, there are two split subspaces whose
intersection is not split and which have no greatest lower bound in
$\Spl(\HH{H})$. Their orthogonal complements have no least upper bound.
\end{theorem}
\begin{proof}
First work in $\HH{H_1\oplus H_1\oplus H_1}$ with
$H_1=\ell^2(\N_{\ge1})$. With $T=\diag(1/n)$ and $q=t^\eta$, define
\begin{align*}
 M&=\{(x,y,0):x,y\in\HH{H_1}\},\\
 N&=\{(x,y,Tx-qy):x,y\in\HH{H_1}\}.
\end{align*}
The first is a constant orthogonal summand. The second is the graph
of the integral adjointable operator
\begin{equation}\label{hgeo:eq:rowB}
 B(x,y)=Tx-qy,
\end{equation}
which is split by \cref{hgeo:lem:finitegraph}. Their intersection is
\begin{equation}\label{hgeo:eq:intersection}
 M\cap N=\ker B\oplus0=\{(x,q^{-1}Tx,0):x\in\HH{H_1}\}.
\end{equation}
Its residue space is $0\oplus\ran T\oplus0$, by \eqref{hgeo:eq:graphresidue}
and $\ker T=0$. That space is not ordinarily closed. \Cref{hgeo:cor:residue}
therefore proves that the intersection is not split.

Suppose a greatest lower bound $L$ of $M,N$ existed in the split
poset. For every nonzero $z\in M\cap N$, the one-dimensional space
$\K z$ is split by finite-dimensional orthogonal projection, and is
a common lower bound. Greatestness implies $\K z\subseteq L$ for
all such $z$. It follows that $M\cap N\subseteq L$, whereas the
lower-bound condition gives $L\subseteq M\cap N$. This would make
the nonsplit intersection itself split, a contradiction. An
order-reversing involution exchanges meets and joins, so
$M^\perp,N^\perp$ have no least upper bound.

Every infinite-dimensional Hilbert space contains a closed isometric
copy of $H_1\oplus H_1\oplus H_1$, obtained by partitioning a countable
orthonormal family into three infinite families and taking their
closed spans. Extend the two projections by zero on the ordinary
orthogonal complement. They give split subspaces in $\HH{H}$ with
the same intersection. Its residue is still nonclosed in $H$, and
the preceding no-meet argument applies unchanged. All choices here
are set-sized.

Finally, if $\dim H=d<\infty$, then $\HH{H}$ is isometric to
$\K^d$ (or $\F^d$ in the real case). Every linear subspace has a
finite orthonormal basis by \cref{hgeo:prop:finite}, hence is
split. Ordinary algebraic intersection and sum are meet and join
in the full subspace lattice, so the split poset is a lattice.
\end{proof}

Source~06 proves the direction ``not a lattice'' for $\ell^2$ (its
Theorem~9.3) with the same construction: its row operator
$(x,y)\mapsto\tau x-Sy$ is \eqref{hgeo:eq:rowB} with the two input
coordinates exchanged, and it transports the example to $\ell^2$ by an
ordinary unitary $\ell^2\oplus\ell^2\oplus\ell^2\cong\ell^2$. The theorem does
not conflict with \cref{hgeo:thm:normalform}: that theorem classifies each
split subspace individually, and intersections need not remain in the
classified family; taking residues does not preserve intersections, even
when both subspaces split.

\begin{remark}[Relation to existing non-Archimedean obstructions \src{03}]
The general phenomenon of nonorthomodularity in other non-Archimedean
Hilbert-like models is not new; Aguayo--Nova give such results for
$c_0$-type spaces \cite{AN}, and Sol\`er's theorem excludes orthomodularity
over these fields altogether (\cref{hgeo:sub:literature}). Here the
coefficient-Hilbert model, the complete classification of split subspaces,
and the precise criterion ``finite ordinary coefficient dimension'' identify a
particular obstruction. In particular, the poset of \emph{split} subspaces
remains orthomodular, but ceases to be a lattice. This is more specific than
saying that some closed subspaces lack orthogonal complements.
\end{remark}

At the level of inner-product geometry, $\Gamma=0$ gives the ordinary
Hilbert space and its complete closed-subspace lattice even in infinite
dimension. The earlier comparison of valuation and norm topologies
assumed $\Gamma\ne0$ and is not being applied to this degenerate case.
Thus nontrivial scale structure is an essential hypothesis of the
dichotomy, not just a notational preference.

\subsection{A positive operator with a noncomplemented kernel}

\begin{theorem}[An adjointable positive operator with a noncomplemented kernel
\src{06}]\label{hgeo:thm:positivekernel}
For every nonzero $\Gamma$, the adjointable positive self-adjoint operator on
$\HH{\ell^2\oplus\ell^2}$
\begin{equation}\label{hgeo:eq:positiveL}
 L=B^*B=
 \begin{pmatrix}
 T^2&-qT\\ -qT&q^2I
 \end{pmatrix},
\end{equation}
with $B$ as in \eqref{hgeo:eq:rowB}, has a valuation-closed, orthoclosed,
spherically complete kernel that is not orthogonally split. It is order
bounded and has only the three exponents $0,\eta,2\eta$.
\end{theorem}
\begin{proof}
All coefficients in \eqref{hgeo:eq:positiveL} are ordinary bounded
operators, so the Hahn operator is adjointable and order bounded.
Positivity follows from $\ip{Lz}{z}=\ip{Bz}{Bz}$. This identity and
anisotropy give $\ker L=\ker B=\{(x,q^{-1}Tx)\}=\Graph(q^{-1}T)$, which by
\cref{hgeo:cor:bothgraphs} is the flip of $M_q(\Lambda)$. The flip is a
constant unitary, so \cref{hgeo:thm:unboundedgraph} applies.
\end{proof}

Thus this obstruction is already polynomial in a single infinitesimal; it is
not caused by a divergent formal inverse or a complicated value group. It
also shows that the adjointable operator algebra is not a Rickart
$*$-algebra: not every element has a projection-generated right
annihilator \src{06}. Put $H=\ell^2\oplus\ell^2$. If the right annihilator
$\{X\in\BH:LX=0\}$ of $L$ were $\{PY:Y\in\BH\}$ for a projection $P$, then
$LP=0$ would imply $\ran P\subseteq\ker L$. For each
$x\in\ker L$, the rank-one operator $z\mapsto\ip{z}{e}x$, for a fixed unit
vector $e$, belongs to that annihilator and would force $x\in\ran P$. Thus
$\ran P=\ker L$, a contradiction.

\subsection{A mutually orthogonal countable family with no join}
\label{hgeo:sub:countable}
Let $(e_n)$ and $(f_n)$ be the standard orthonormal bases of the two
summands of $\ell^2\oplus\ell^2$. Define
\[
 v_n=e_n+nqf_n,\qquad w_n=-nqe_n+f_n,
\]
and let $P_n$ and $P_n'$ be the rank-one projections onto $\K v_n$ and
$\K w_n$. In the $n$th two-dimensional block,
\begin{equation}\label{hgeo:eq:Pn}
 P_n=\frac1{1+n^2q^2}
 \begin{pmatrix}1&nq\\nq&n^2q^2\end{pmatrix},\qquad
 P_n'=\frac1{1+n^2q^2}
 \begin{pmatrix}n^2q^2&-nq\\-nq&1\end{pmatrix}.
\end{equation}
They vanish outside that block. The $P_n$ are mutually orthogonal,
and $P_mP_n'=0$ for all $m,n$.

\begin{theorem}[No countable orthogonal join at one scale \src{06}]
\label{hgeo:thm:nojoin}
For every nonzero $\Gamma$, the family $(P_n)_{n\ge1}$ in
\eqref{hgeo:eq:Pn} has no least upper bound among the projections on
$\HH{\ell^2\oplus\ell^2}$, although $I$ is an upper bound. Every
individual projection has support in the same well-ordered monoid
$\N\eta$. Thus common well-ordered scalar support does not restore
countable projection completeness.
\end{theorem}
\begin{proof}
The coordinate relations $\ip{(x,y)}{w_n}=0$ for all $n$ say precisely
$y_n=nqx_n$ for the $n$th coordinates. Reconstructing Hilbert coefficients
coefficient by coefficient, as in \cref{hgeo:cor:bothgraphs}, gives
\[
 \bigcap_{n\ge1}\ker P_n'=M_q(\Lambda).
\]
Suppose $Q$ were the least upper bound of all $P_n$. Since $I-P_n'$ is an
upper bound for every $n$, minimality gives $Q\le I-P_n'$ for every $n$,
so $\ran Q\subseteq M_q(\Lambda)$. On the other hand, $v_n\in\ran Q$
for all $n$, and $\res v_n=(e_n,0)$. By \cref{hgeo:cor:residue}, the ordinary
subspace $\res\ran Q$ is closed. It contains every $(e_n,0)$, hence all of
$\ell^2\oplus0$. But it must lie in $\res M_q(\Lambda)=\Dom\Lambda\oplus0$,
which is proper. This is impossible. The support assertion follows by
expanding $(1+n^2q^2)^{-1}$.
\end{proof}

\begin{corollary}[Finite straightening has no countable analogue \src{06}]
\label{hgeo:cor:nosimultaneous}
There is no adjointable unitary simultaneously sending the constant
rank-one projections $Q_n$ onto $\K e_n$ to all the $P_n$ in
\eqref{hgeo:eq:Pn}.
\end{corollary}
\begin{proof}
If such a unitary $U$ existed, it would have nonnegative support
(\cref{hgeo:cor:orbits}), and its residue would preserve each constant line
$\kk e_n$, since $(P_n)_0=Q_n$. Write $U=\sum_{\gamma\ge0}U_\gamma
t^\gamma$. The coefficient at $\eta$ of $UQ_nU^*$ is
$\sum_{\alpha+\beta=\eta}U_\alpha Q_nU_\beta^*$ with $\alpha,\beta\in\supp U$;
the fixed well-ordered support of $U$ allows only finitely many such pairs,
including pairs of two strictly positive exponents summing to $\eta$, and
their operator norms are bounded independently of $n$, since $\norm{Q_n}=1$.
Thus $\sup_n\norm{\coeff{\eta}{UQ_nU^*}}<\infty$. Yet \eqref{hgeo:eq:Pn}
gives $P_n$ an exponent-$\eta$ coefficient of norm $n$. This contradiction
proves the claim.
\end{proof}

The proof retains possible intermediate exponents of $U$; discarding them
would leave a gap at a nondiscrete value group. The obstruction here is
\emph{ordinary boundedness of coefficient operators}, not support shape.

\begin{proposition}[Incoherent supports prevent simultaneous conjugation
\src{06}]\label{hgeo:prop:incoherent}
Suppose $\Gamma$ contains $\Q\eta$. Let $\widetilde P_n$ project onto
$\K(e_n+t^{\eta/n}f_n)$. Each $\widetilde P_n$ is a legitimate projection,
but no adjointable unitary conjugates the constant projections $Q_n$ onto
$\K e_n$ simultaneously to the $\widetilde P_n$.
\end{proposition}
\begin{proof}
If $S=\supp U$ is the common well-ordered operator support, then
$\supp(UQ_nU^*)\subseteq S+S$ for every $n$. The set $S+S$ is well
ordered, so it cannot contain the strictly decreasing sequence
$\eta,\eta/2,\eta/3,\ldots$, whereas the off-diagonal coefficient of
$\widetilde P_n$ at $\eta/n$ is nonzero.
\end{proof}

Thus infinite orthogonal synthesis has two genuinely distinct requirements:
a coherent common support and bounded ordinary coefficient synthesis.
Finite families automatically meet both; arbitrary countable families do
not. Sources~03 and~06 prove them as necessary conditions only. \mergetag{}
Source~07 shows that, in the form of bounded synthesized coefficients and
well-ordered active support, they are together also sufficient
(\cref{hgeo:us:thm:main}), and that neither can be dropped, even under
uniform bounds on the individual coefficients
(\cref{hgeo:us:thm:uniform,hgeo:us:thm:support}); this answers
\cref{hgeo:q:synthesis}.

\section{Metric rigidity and the exact positivity threshold}
\label{hgeo:sec:metric}

The preceding classification concerns the canonical Hahn extension of an
ordinary Hilbert inner product. There are other positive definite forms on
the same vector space. Even valuation completeness and a positive surreal
coercivity bound do not force such a form to be isometrically standard. The
relevant threshold is an \emph{ordinary real} lower bound on its residue
operator \src{06}.

Let $H\ne0$, and let
\begin{equation}\label{hgeo:eq:metricTheta}
 \Theta=\widehat{\Theta_0}+\Theta_+\in\BH,\qquad \Theta=\Theta^*,\quad
 v(\Theta_+)>0\text{ or }\Theta_+=0,
\end{equation}
where $\Theta_0\in\mathcal B(H)$ is ordinary self-adjoint and strictly
positive in the pointwise sense $\ip{\Theta_0u}{u}_H>0$ for $u\ne0$. Strict
positivity here does not mean uniform positivity. The form
\begin{equation}\label{hgeo:eq:innerTheta}
 \ip{x}{y}_\Theta=\ip{\Theta x}{y}
\end{equation}
is positive definite: if $x_\delta$ is the first nonzero coefficient of
$x$, its squared length has first term $t^{2\delta}\ip{\Theta_0x_\delta}
{x_\delta}_H$. Its positive square root exists by the same binomial
construction as in \cref{hgeo:prop:inner}, and
\begin{equation}\label{hgeo:eq:metricval}
 v(\norm{x}_\Theta)=v(x)=v(\norm{x}).
\end{equation}
In particular the valued-space completeness properties of
\cref{hgeo:thm:complete} are unchanged. This concerns valuation balls; it
does not upgrade the assertion to arbitrary ordered-field norm balls.

\begin{theorem}[Exact isometric rigidity threshold \src{06}]
\label{hgeo:thm:metric}
For $\Theta$ as in \eqref{hgeo:eq:metricTheta}, the following are equivalent:
\begin{enumerate}[label=\textup{(\roman*)}]
\item There is a $\K$-linear bijection $J:\HH{H}\to\HH{H}$ with
\[
 \ip{Jx}{Jy}=\ip{x}{y}_\Theta\qquad(x,y\in\HH{H}).
\]
No continuity, strongness, or adjointability of $J$ is assumed.
\item The residue $\Theta_0$ is uniformly positive: for some ordinary real
$c>0$, $\ip{\Theta_0u}{u}_H\ge c\norm{u}_H^2$ for every $u\in H$.
\end{enumerate}
When (ii) holds, an isometry is explicitly
\begin{equation}\label{hgeo:eq:metricJ}
 J=(I+\Theta')^{1/2}\widehat{\Theta_0^{1/2}},\qquad
 \Theta'=\widehat{\Theta_0^{-1/2}}\,\Theta_+\,\widehat{\Theta_0^{-1/2}}.
\end{equation}
\end{theorem}
\begin{proof}
Assume (ii). Ordinary positive operator theory gives bounded
$\Theta_0^{1/2}$ and $\Theta_0^{-1/2}$. The self-adjoint $\Theta'$ has
positive order, so its formal binomial square root and inverse square root
exist. The operator in \eqref{hgeo:eq:metricJ} is bijective and satisfies
\[
 J^*J=\widehat{\Theta_0^{1/2}}(I+\Theta')\widehat{\Theta_0^{1/2}}=\Theta.
\]
This proves (i), without a commutation assumption on $\Theta_0$ and
$\Theta_+$.

Conversely suppose (i). Since $J$ is onto, for every $z$ we have
\[
 \ip{Jx}{z}=\ip{\Theta x}{J^{-1}z}=\ip{x}{\Theta J^{-1}z}.
\]
Thus $J$ has an everywhere-defined adjoint $J^*=\Theta J^{-1}$, and
\cref{hgeo:thm:adjoint} applies: $J$ is a Hahn series of ordinary bounded
operators. Since $J^*J=\Theta$ and $v(\Theta)=0$, \cref{hgeo:cor:vop} gives
$2v(J)=0$, hence $v(J)=0$. Its residue satisfies
\begin{equation}\label{hgeo:eq:Jresidue}
 J_0^*J_0=\Theta_0.
\end{equation}
The strict positivity assumption makes $J_0$ injective.

It is also surjective. For a nonzero constant vector $u$, choose $x$ with
$Jx=u$. If $\delta=v(x)<0$, then the leading term of $Jx$ is
$t^\delta J_0x_\delta\ne0$, a contradiction. If $\delta>0$, the output
cannot have a nonzero constant coefficient. Thus $\delta=0$ and
$J_0x_0=u$. The ordinary bounded inverse theorem yields a bounded inverse for
$J_0$; \eqref{hgeo:eq:Jresidue} consequently gives
$\Theta_0\ge\norm{J_0^{-1}}^{-2}I$. This proves (ii).
\end{proof}

The argument does not assume that an algebraically bijective adjointable
operator automatically has an adjointable inverse. The adjoint of $J$ was
constructed directly, and the ordinary inverse theorem was used only after
proving that $J_0$ is onto.

\begin{corollary}[Coercive equivalent norms need not be isometric \src{06}]
\label{hgeo:cor:coercivemetric}
Let $T=\diag(1/n)$ on $\ell^2$. Then
\[
 \Theta=\widehat T^{\,2}+q^2I
\]
defines a positive definite form on $\HH{\ell^2}$ with
\begin{equation}\label{hgeo:eq:coercivenorm}
 q\norm{x}\le\norm{x}_\Theta\le 2\norm{x}\qquad(x\in\HH{\ell^2}),
\end{equation}
yet there is no surjective $\K$-linear isometry from this space to the
canonical Hahn--Hilbert space. Both have the same valuation and are
spherically complete.
\end{corollary}
\begin{proof}
The lower bound follows from
$\ip{\Theta x}{x}=\norm{\widehat Tx}^2+q^2\norm{x}^2$.
The ordinary contraction $T$ remains a contraction after Hahn extension:
$I-T^2$ has an ordinary positive square root, so the same identity of
positive forms holds over $\K$. Hence
$\norm{x}_\Theta^2\le(1+q^2)\norm{x}^2<4\norm{x}^2$ for $x\ne0$.
But $\Theta_0=T^2$ is not uniformly positive, since
$\ip{T^2e_n}{e_n}=n^{-2}$. Apply \cref{hgeo:thm:metric} and
\eqref{hgeo:eq:metricval}.
\end{proof}

The operator $T^2+q^2I$ is the spectral report's coercive nonsurjective
operator $D^2+s^2I$ of \replabel{ihs:hh:thm:coercive}. What the criterion adds
is an exact obstruction to \emph{all} bijective linear isometries of the
corresponding positive forms, including maps whose regularity was not
assumed. At an ordinary uniformly positive residue, positive-order
perturbations are removable by \eqref{hgeo:eq:metricJ}; below that
threshold, infinitesimal coercivity does not repair the missing ordinary
lower bound.

\section{Fredholm perturbations reduce to a finite defect matrix}
\label{hgeo:sec:schur}

\subsection{The hypotheses and the reduction}
Let $T_0:H\to G$ be an ordinary bounded Fredholm operator. This means
that its ordinary range is closed and that its kernel and cokernel
are finite dimensional. Write
\[
 N=\ker T_0,\quad X=N^\perp,\qquad
 Y=\ran T_0,\quad N_*=Y^\perp=\ker T_0^*.
\]
Let $m=\dim N$, $n=\dim N_*$. The restriction
$B_0=T_0|_X:X\to Y$ is a bounded isomorphism. Zero summands are permitted;
identity and inverse formulas on a zero space are interpreted in their
unique empty-block sense.

Suppose
\[
 A=T_0+E\in\BB{H}{G},\qquad E=0\text{ or }v(E)>0.
\]
Relative to $H=X\oplus N$, $G=Y\oplus N_*$, write
\begin{equation}\label{hgeo:eq:blocks}
 A=\begin{pmatrix}B&A_{12}\\A_{21}&A_{22}\end{pmatrix}.
\end{equation}
Here $B=B_0+\text{positive-order terms}$ is invertible by the
Neumann formula, and the other three blocks have positive order
or are zero. Define the finite \emph{Schur defect matrix}
\begin{equation}\label{hgeo:eq:schur}
 S=A_{22}-A_{21}B^{-1}A_{12}:\HH{N}\longrightarrow\HH{N_*}.
\end{equation}
Choosing ordinary orthonormal bases identifies it with an
$n\times m$ matrix over $\K$. Its entries have positive valuation
or are zero.

\begin{theorem}[Exact finite-defect reduction \src{03}]\label{hgeo:thm:schur}
With these hypotheses, the integral block operators
\begin{equation}\label{hgeo:eq:LR}
 L=\begin{pmatrix}I&0\\-A_{21}B^{-1}&I\end{pmatrix},\qquad
 R=\begin{pmatrix}I&-B^{-1}A_{12}\\0&I\end{pmatrix}
\end{equation}
have integral inverses, are near the identity, and satisfy
\begin{equation}\label{hgeo:eq:schuridentity}
 LAR=\begin{pmatrix}B&0\\0&S\end{pmatrix}=:\mathcal D.
\end{equation}
For $r=\rank_\K S$,
\begin{align}
 \dim_\K\ker A&=m-r,\label{hgeo:eq:kerdim}\\
 \dim_\K(\HH{G}/\ran A)&=n-r,\label{hgeo:eq:cokerdim}\\
 \ind_\K A&=m-n=\ind_\C T_0.\label{hgeo:eq:index}
\end{align}
All identities are exact Hahn identities at arbitrary rank.
\end{theorem}
\begin{proof}
Factor $B=B_0(I+B_0^{-1}(B-B_0))$. Its inverse belongs to
$\BB{Y}{X}$, has integral valuation, and has residue $B_0^{-1}$.
The off-diagonal blocks in \eqref{hgeo:eq:LR} are positive order.
Changing their signs gives the respective inverses of $L,R$.
Direct multiplication first clears the upper-right block by $R$
and then the lower-left block by $L$, leaving precisely
$A_{22}-A_{21}B^{-1}A_{12}$ in the lower-right corner. This proves
\eqref{hgeo:eq:schuridentity}.

Because $L,R$ are bijections and $B$ is an isomorphism, the kernel
and cokernel of $A$ are isomorphic to those of the finite matrix $S$.
Finite rank-nullity gives \eqref{hgeo:eq:kerdim}--\eqref{hgeo:eq:index}.
For support validity, all inverse expansions use the positive
monoid generated by the support of $E$; multiplication and the
finite block decompositions preserve well-ordered support. No
analytic or norm-convergent infinite process is involved.
\end{proof}

This is a standard Schur-complement calculation applied in the
specified Hahn algebra. Its role is to isolate a finite part where
orthogonal linear algebra is safe. In particular, ``Fredholm'' in
\eqref{hgeo:eq:index} is the finite algebraic kernel/cokernel condition
on these named spaces; it is not a compactoid-operator definition in
a different non-Archimedean category.

\begin{remark}[Comparison with the Hahn Fredholm alternative \mergetag]
\label{hgeo:rem:fredholm}
The spectral report's \replabel{ihs:fr:thm:fredholm} treats $I+C$ with $C$
a Hahn series of trace-class operators of nonnegative valuation on one
space $H$. It proves $\Det(I+C)\ne0\iff(I+C)^{-1}\in\BH\iff I+C$ is
bijective, that $I+C$ is algebraically Fredholm of index zero, with kernel
and cokernel of the same finite dimension, and a factorization of the
determinant through one finite matrix. The determinant-free clauses follow
from \cref{hgeo:thm:schur,hgeo:thm:MP}: $C=C_0+C_+$ with $C_0$ ordinary
trace class and $v(C_+)>0$, so $T_0=I+C_0$ is an ordinary Fredholm operator
of index zero and $A=T_0+C_+$ satisfies the hypotheses with $m=n$; thus
$\dim\ker A=\dim\coker A=m-r$, and if $A$ is bijective then
$A^{-1}=A^\dagger\in\BH$. \Cref{hgeo:thm:schur,hgeo:thm:MP} are more general
in these clauses: $T_0$ may have any index and act between different spaces,
and the perturbation need only have bounded, not trace-class, coefficients.
They say nothing about the determinant, its finite factorization, or the
spectral questions of that report. Source~03 does not make this comparison.
\end{remark}

\subsection{Finite Moore--Penrose geometry}
A finite matrix $S:\K^m\to\K^n$ has a Moore--Penrose inverse over $\K$ with
its standard positive Hermitian form. The finite spectral-theory report
(\fn{docs/surcomplex/spectral-theory}) constructs it by the singular value
decomposition, proves that the four identities
\begin{equation}\label{hgeo:eq:penrose}
 SS^\dagger S=S,\quad S^\dagger SS^\dagger=S^\dagger,
 \quad (SS^\dagger)^*=SS^\dagger,
 \quad (S^\dagger S)^*=S^\dagger S
\end{equation}
characterize it, and proves the least-squares theorem
\replabel{thm:leastsquares}, first over a real closed field; its SVD over a
nondivisible Hahn field, \replabel{spec:thm:svd}, makes the same arguments
available over $\K$ for every $\Gamma$. Source~03 constructs $S^\dagger$
without the SVD: finite Gram--Schmidt, which needs only the vector-norm roots
of \cref{hgeo:prop:inner}, gives
\[
 \K^m=\ker S\oplus(\ker S)^\perp,\qquad
 \K^n=\ran S\oplus(\ran S)^\perp;
\]
the restriction of $S$ between $(\ker S)^\perp$ and $\ran S$ is a bijection
of finite-dimensional spaces; invert it there and set it to zero on
$(\ran S)^\perp$. Neither construction requires $\K$ algebraically closed
or $\F$ real closed, and no divisibility of $\Gamma$ is inserted. These
finite identities and their geometry are classical Moore--Penrose linear
algebra \cite{Penrose}, not proposed new results.

\section{Exact Moore--Penrose inverses and least squares}
\label{hgeo:sec:least}

\subsection{Finite defects restore orthogonal range decomposition}

\begin{theorem}[Fredholm-residue least-squares theorem \src{03}]
\label{hgeo:thm:MP}
For $A=T_0+E$ as in \cref{hgeo:sec:schur}, both $\ker A$ and
$\ran A$ are orthogonally split. There is a unique
$A^\dagger\in\BB{G}{H}$ satisfying the four Penrose identities
\eqref{hgeo:eq:penrose} with $S$ replaced by $A$.
If $P$ projects onto $\ker A$ and $Q$ onto $\ran A$, it is explicitly
\begin{equation}\label{hgeo:eq:MPformula}
 A^\dagger=(I-P)\,
 R\begin{pmatrix}B^{-1}&0\\0&S^\dagger\end{pmatrix}L\,Q.
\end{equation}
In particular
\begin{equation}\label{hgeo:eq:MPproj}
 AA^\dagger=Q,\qquad A^\dagger A=I-P.
\end{equation}
\end{theorem}
\begin{proof}
The kernel is finite dimensional by \cref{hgeo:thm:schur}, so
its orthogonal projection $P$ exists. To prove the exact range
assertion rather than assume a Hilbert closed-range theorem, use
$A=L^{-1}\mathcal D R^{-1}$ from \eqref{hgeo:eq:schuridentity}. Then
\[
 A^*=R^{-*}\mathcal D^*L^{-*},\qquad
 \ker A^*=L^*\ker\mathcal D^*.
\]
Here $L^{-*}=(L^*)^{-1}$ and similarly for $R$. On the finite Schur
block, ordinary finite orthogonal algebra over $\K$ gives
$(\ker S^*)^\perp=\ran S$. On the other block $B$ is invertible.
Therefore
\[
 (\ker\mathcal D^*)^\perp=\ran\mathcal D.
\]
For any subspace $Z$, the adjoint identity gives
$(L^*Z)^\perp=L^{-1}Z^\perp$. Consequently
\begin{equation}\label{hgeo:eq:closedrangeexact}
 (\ker A^*)^\perp=L^{-1}\ran\mathcal D=\ran A.
\end{equation}
The subspace $\ker A^*$ has finite dimension $n-r$, so it is split.
Its orthogonal complement \eqref{hgeo:eq:closedrangeexact} is split
as well; let $Q=I-P_{\ker A^*}$ be its projection. This proves
splitting using the finite reduction, not merely finite codimension.

Put
\[
 A^-=R\begin{pmatrix}B^{-1}&0\\0&S^\dagger\end{pmatrix}L.
\]
The block identity gives $AA^-A=A$, so $AA^-$ is the identity on
$\ran A$. Since $AP=0$, the map
$A'=(I-P)A^-Q$ satisfies $AA'=Q$. Moreover for every $x$,
$A(A^-Ax-x)=0$, so
\[
 A'Ax=(I-P)A^-Ax=(I-P)x.
\]
Thus $A'A=I-P$. The range of $A'$ lies in $(\ker A)^\perp$, and
$A'=A'Q$, which yields $A'AA'=A'$. The two product identities give
self-adjointness of $AA',A'A$ and hence all four Penrose equations.
Every factor in $A'$ is adjointable: the finite Schur inverse and
finite-rank projections are Hahn operators, and all remaining
factors were already constructed in the Hahn algebra.

For uniqueness, any Penrose inverse $A''$ makes $AA''$ a
self-adjoint idempotent with range $\ran A$, hence $AA''=Q$.
Similarly $A''A$ is a self-adjoint idempotent with kernel $\ker A$,
hence $A''A=I-P$. Also $A''=A''AA''$, so its values lie in
$(\ker A)^\perp$. Both $A''f$ and $A'f$ lie there and solve
$Ax=Qf$; their difference lies in $\ker A$, and is therefore zero.
\end{proof}

\begin{corollary}[All least-squares and minimum-norm solutions \src{03}]
\label{hgeo:cor:LS}
For every $f\in\HH{G}$ let $u_0=A^\dagger f$. Then
\begin{equation}\label{hgeo:eq:LSidentity}
 \norm{Ax-f}^2=\norm{(I-Q)f}^2+\norm{A(x-u_0)}^2
 \qquad(x\in\HH{H}).
\end{equation}
The least-squares minimizers are exactly $u_0+\ker A$.
Among them $u_0$ is the unique vector of minimum norm. The equation
$Ax=f$ is solvable exactly when $(I-Q)f=0$.
\end{corollary}
\begin{proof}
Write $Ax-f=A(x-u_0)-(I-Q)f$. Its two terms lie in the orthogonal
spaces $\ran A$ and $\ker A^*$, respectively. Pythagoras proves
\eqref{hgeo:eq:LSidentity}; positivity identifies equality with
$A(x-u_0)=0$. Since $u_0\in(\ker A)^\perp$, one also has
$\norm{u_0+k}^2=\norm{u_0}^2+\norm{k}^2$ for $k\in\ker A$.
This proves the minimum-norm and exact-solvability assertions.
\end{proof}

For a finite matrix this is the finite report's \replabel{thm:leastsquares};
the content here is its transport through the infinite Schur reduction.

\begin{remark}[Why this does not contradict the graph obstruction \src{03}]
\Cref{hgeo:ex:badgraph} has an infinite-dimensional ordinary range
defect, invisible to a finite matrix: its residue range is dense
and not closed. \Cref{hgeo:thm:MP} starts with an ordinarily closed
Fredholm range and leaves only finitely many kernel and cokernel
coordinates to resolve. A finite matrix over $\K$ always has the
orthogonal decompositions needed here. The two results therefore
identify different regimes, not competing definitions of solvability.
\end{remark}

\section{The exact amplification scale from finite minors}
\label{hgeo:sec:valuation}

\subsection{Finite valuation-ring reduction}
The valuation ring $\OO$ need not be Noetherian or discretely valued.
Nevertheless finite matrices admit the following reduction.

\begin{lemma}[Finite diagonal reduction over the valuation ring
\src{03}; = \replabel{thm:scales}]\label{hgeo:lem:smith}
Let $S\in M_{n\times m}(\OO)$ have rank $r>0$. There are
$U_1\in\GL_n(\OO)$, $V_1\in\GL_m(\OO)$ and exponents
$0\le s_1\le\cdots\le s_r$ such that
\begin{equation}\label{hgeo:eq:smith}
 U_1SV_1=\diag(t^{s_1},\ldots,t^{s_r},0).
\end{equation}
If every nonzero entry of $S$ has positive valuation, then $s_1>0$.
For $1\le j\le r$, set
\[
 \nu_j(S)=\min\{v(\det S_{I,J}):\abs I=\abs J=j,
                         \det S_{I,J}\ne0\},\qquad \nu_0(S)=0.
\]
Then
\begin{equation}\label{hgeo:eq:minors}
 \nu_j(S)=s_1+\cdots+s_j,
\end{equation}
and the exponents are invariants of $S$ under integral invertible row and
column operations.
\end{lemma}

This is the finite spectral-theory report's theorem ``singular scales are
determinantal scales'', \replabel{thm:scales}, with its invariance lemma
\replabel{lem:DeltaInvariant}: there $\nu_j=\Delta_j$ and $s_j=\gamma_j$ is the
$j$th singular valuation, and the proof goes through the SVD
(\replabel{spec:thm:svd} for nondivisible $\Gamma$). Source~03 proves it by
elimination, without any spectral theorem; the route is kept because it
makes the generality of the inverse-scale formula explicit.
\begin{proof}[Second route by elimination \src{03}]
Choose a nonzero entry of minimum valuation and permute it to the
upper-left corner. Write it as $p=t^s u$ with $u\in\OO^\times$.
Multiplying its row by $u^{-1}$ makes the pivot $t^s$; all entries
still have valuation at least $s$. Because the pivot has minimum
valuation, it divides every entry in $\OO$. Subtract multiples of
the first row to clear the first column, then multiples of the
first column to clear the first row. These are invertible integral
row and column operations. The remaining lower-right block still
has all entries of valuation at least $s$. Repeat on this finite block.
The process terminates after $r$ nonzero pivots and produces
nondecreasing exponents. This proves \eqref{hgeo:eq:smith} and the
strict positivity assertion. Cauchy--Binet shows that the least nonzero
$j$-minor valuation cannot decrease under an integral left or right
multiplication, and applying the integral inverse gives the reverse
inequality. For the diagonal matrix in \eqref{hgeo:eq:smith}, the least
nonzero $j$-minor valuation is $s_1+\cdots+s_j$.
\end{proof}

The diagonal exponents in this lemma are valuation invariants, not
asserted eigenvalue or singular-value exponents of an arbitrary
infinite operator. No spectral theorem is being used for infinite operators.

\subsection{A sharp formula for the Moore--Penrose inverse}

\begin{theorem}[Exact finite-defect amplification \src{03}]\label{hgeo:thm:scale}
Let $A=T_0+E$ and $S$ be as in \cref{hgeo:thm:MP}. If
$r=\rank S>0$, then
\begin{equation}\label{hgeo:eq:scale}
 v(A^\dagger)=-s_r
       =-\bigl(\nu_r(S)-\nu_{r-1}(S)\bigr)<0,
\end{equation}
where the $s_j$ are the exponents in \cref{hgeo:lem:smith}.
If $S=0$ and $T_0\ne0$, then
\begin{equation}\label{hgeo:eq:nolift}
 v(A^\dagger)=0,
 \qquad [t^0]A^\dagger=T_0^\dagger.
\end{equation}
If $S=0=T_0$, then $A=A^\dagger=0$.
\end{theorem}
\begin{proof}
Combine the Schur reduction, multiplication of the regular block
by $B^{-1}$, and the integral finite transformations in
\cref{hgeo:lem:smith}. There are bounded-coefficient Hahn
isomorphisms $\mathsf L$ on the codomain side and $\mathsf R$ on the domain
side, with integral inverses, such that
\begin{equation}\label{hgeo:eq:fullnormal}
 \mathsf LA\mathsf R=\mathcal E
  =\diag(I_X,t^{s_1},\ldots,t^{s_r},0).
\end{equation}
To write the regular identity $I_X$ we identify $Y$ with $X$ by
$B_0$ or $B$ as appropriate; this is an isomorphism between the
named coefficient summands, not a dimension-changing operation.
Every change and its inverse is integral. Zero regular summands
are allowed.

Assume first $r>0$. The diagonal map
\[
 Z_0=\diag(I_X,t^{-s_1},\ldots,t^{-s_r},0)
\]
is an inner generalized inverse of $\mathcal E$ and has valuation
$-s_r$. Therefore
$A^-_1=\mathsf RZ_0\mathsf L$ satisfies $AA^-_1A=A$ and has valuation $-s_r$,
because integral invertible multipliers preserve operator valuation.
The proof of \cref{hgeo:thm:MP} works with any such inner inverse and gives
\[
 A^\dagger=(I-P)A^-_1Q.
\]
Both orthogonal projections are integral by
\cref{hgeo:lem:integralproj}. Hence
$v(A^\dagger)\ge-s_r$.

For the converse bound put
$Z=\mathsf R^{-1}A^\dagger\mathsf L^{-1}$. The identity
$AA^\dagger A=A$ implies $\mathcal E Z\mathcal E=\mathcal E$.
At the $j$th nonzero finite pivot, its diagonal entry says
\[
 t^{s_j}Z_{jj}t^{s_j}=t^{s_j},
 \qquad\text{so}\qquad Z_{jj}=t^{-s_j}.
\]
A nonzero scalar coordinate of valuation $-s_r$ forces
$v(Z)\le-s_r$. The integral invertible transformations preserve
valuation, so $v(Z)=v(A^\dagger)$. This proves equality; the
minor formula follows from \eqref{hgeo:eq:minors}.

If $S=0$, the inner inverse
$A^-_0=R\diag(B^{-1},0)L$ is integral and its residue
is $T_0^\dagger$ relative to the ordinary orthogonal decompositions.
Thus $A^\dagger=(I-P)A^-_0Q$ is integral. To identify the projection
residues, observe from the Schur formulas with $S=0$ that
\[
 \ker A=R(0\oplus\HH{N}),\qquad
 \ker A^*=L^*(0\oplus\HH{N_*}).
\]
Because $R,L^*$ and their inverses are near the identity, their
integral residue spaces are $N,N_*$. \Cref{hgeo:cor:residue}
therefore gives $P_0=P_N$ and $Q_0=P_Y$. Taking residues in
\eqref{hgeo:eq:MPformula} yields
$[t^0]A^\dagger=(I-P_N)T_0^\dagger P_Y=T_0^\dagger$.
This is nonzero if $T_0\ne0$, proving \eqref{hgeo:eq:nolift}.
If $T_0=0$ is Fredholm, both $H,G$ are finite dimensional and the
regular block is empty; then $S=A$, so $S=0$ gives $A=0$.
\end{proof}

At the finite level, \eqref{hgeo:eq:scale} for $A=S$ also follows from the
finite report: in an SVD $S=U\Sigma V^*$ the unitary factors are integral
with integral inverses (\replabel{lem:unitaryintegral}), so
$v(S^\dagger)=v(\Sigma^\dagger)=-\gamma_r$, which is
$-(\Delta_r-\Delta_{r-1})$ by \replabel{thm:scales}. The content of
\cref{hgeo:thm:scale} is the transport of this formula through the infinite
Schur reduction.

\begin{corollary}[Exactly when the inverse stays finite \src{03}]
\label{hgeo:cor:finiteinverse}
For a nonzero ordinary Fredholm $T_0$, the following are equivalent:
$A^\dagger$ is integral; $S=0$; $\dim_\K\ker A=\dim_\C\ker T_0$.
They are also equivalent to preservation of the ordinary cokernel
dimension. If even one finite defect direction is lifted, the
Moore--Penrose inverse acquires a strictly negative valuation,
exactly as measured by \eqref{hgeo:eq:scale}.
\end{corollary}
\begin{proof}
Use \eqref{hgeo:eq:kerdim}, \eqref{hgeo:eq:cokerdim}, and
\cref{hgeo:thm:scale}; all $s_j$ are strictly positive.
\end{proof}

\subsection{An off-diagonal term can dominate the inverse scale}

\begin{example}[A fourth-order pole from second- and third-order diagonals
\src{03}]\label{hgeo:ex:pole}
Put $q=t^\eta$, and on $H_0\oplus\C^2$ take
\[
 T_0=I_{H_0}\oplus0_2,
 \qquad A=I_{H_0}\oplus
 S,\qquad S=\begin{pmatrix}q^2&q\\0&q^3\end{pmatrix}.
\]
Here $H_0$ may be infinite dimensional. The finite Schur matrix is
exactly $S$, with
\[
 \nu_1(S)=\eta,\qquad \nu_2(S)=5\eta,
 \qquad s_1=\eta,
 \quad s_2=4\eta.
\]
Since $A$ is bijective, $A^\dagger=A^{-1}$, and explicitly
\[
 S^{-1}=\begin{pmatrix}q^{-2}&-q^{-4}\\0&q^{-3}\end{pmatrix}.
\]
The inverse valuation is $-4\eta$, not $-2\eta$ or $-3\eta$.
Off-diagonal coupling is measured by minors and cannot be recovered
merely by inspecting the orders on the original diagonal.
\end{example}

More generally, for positive $\alpha,\beta,\gamma\in\Gamma$,
\[
 S=\begin{pmatrix}t^\alpha&t^\beta\\0&t^\gamma\end{pmatrix}
 \quad\Longrightarrow\quad
 v(S^{-1})=-\bigl(\alpha+\gamma-
                       \min\{\alpha,\beta,\gamma\}\bigr).
\]
This follows either from the two-minor formula or by writing its
inverse. The expression makes sense in any ordered abelian group.
For example, in $\Q\oplus\Q$ with lexicographic order take
$\alpha=(1,0)$, $\beta=(0,1)$, $\gamma=(1,1)$. Then the inverse
valuation is $-(2,0)$. There is no need to collapse these scales to
real numbers or choose an Archimedean embedding.

\section{Coherent extension of scales and the full surreal class}
\label{hgeo:sec:global}

\subsection{Set-sized scalar extension preserves the formulas}
An ordered embedding $\Gamma\hookrightarrow\Delta$ sends a
well-ordered support to a well-ordered support, and defines embeddings
of scalar fields, vector spaces, and bounded-coefficient Hahn
operator algebras. It preserves products, conjugation, adjoints,
and the strong sums used in this report. A projection therefore
extends to a projection on the larger Hahn space.

\begin{proposition}[Coherence of the constructive results \src{03, 06}]
\label{hgeo:prop:extension}
Under any ordered embedding $\Gamma\hookrightarrow\Delta$:
\begin{enumerate}[label=\textup{(\roman*)}]
\item the projection normal form, its direct rotation, the graph formulas
and the finite simultaneous straightening extend by coefficient transport;
\item for a Fredholm-residue $A$, its transported Moore--Penrose
inverse is the Moore--Penrose inverse of the transported operator;
\item the Schur rank is unchanged, its nonzero minor valuations are
transported, and the exact amplification formula is unchanged under
that transport.
\end{enumerate}
No cofinality assumption on the embedding is required.
\end{proposition}
\begin{proof}
The normal form and the other constructions in (i) are made from finite
algebraic expressions and positive-support binomial series, all preserved
by the embedding, since they use exactly the same coefficients and finite
contributions at each exponent. The transported Moore--Penrose inverse still
satisfies the four Penrose identities on the larger space, since they are
operator Hahn identities. Uniqueness in \cref{hgeo:thm:MP} identifies it.
A finite determinant is preserved by an injective field map, so the
largest nonzero minor size, namely the rank, is unchanged. Its
valuation is mapped by the ordered embedding, which preserves
minima of finite sets; this proves the final assertion.
\end{proof}

These are algebraic and strong-summation coherence statements. They do not
assert that valuation completions or topological closures commute
with a noncofinal scale enlargement. The three-duals report explains
why those topological operations behave differently.

\subsection{A literal vector space over all surcomplex scalars}
For readers who want all of $\No[i]$, rather than a specified set
subfield, there is a precise class formulation. Work here in
von Neumann--Bernays--G\"odel class theory without assuming global
choice, and keep the ordinary Hilbert spaces $H,G$ as sets. Define
\[
 \mathscr H_{\No}(H)=
 \left\{\sH_{\gamma\in S}h_\gamma t^\gamma:
 S\subset\No\text{ is a set and well ordered},\ h_\gamma\in H\right\},
 \qquad t^\gamma=\omega^{-\gamma}.
\]
Vectors can be encoded as sets of nonzero coefficient pairs. The
collection of all these vectors is a proper class when $H\ne0$.
By Conway normal forms the scalar class
$\C((t^{\No}))$ is exactly $\No[i]$ in this notation, and $\R((t^{\No}))$ is
$\No$ \cite{Conway}. The vector operations and inner product remain the same
coefficient formulas. This gives a compatible system of spaces, not the
generally smaller algebraic tensor product of $H$ with the full scalar field.

Every set of vectors and scalars belongs to a common set-sized workspace
\src{06}: collect its exponents and let $\Gamma$ be the divisible ordered
subgroup of $(\No,+)$ generated by them together with $1$. The set of finite
rational combinations of a set is a set, and this $\Gamma$ contains all the
original supports. The real-closed or algebraically closed Hahn scalar fields
can therefore be used locally. All inner-product identities, finite
Gram--Schmidt arguments, and positive-order formal series involving finitely
or set-many specified inputs can be verified in a common workspace. This
supplies a literal positive inner-product vector class over all surcomplex
scalars, rather than calling a proper class a set. A separate issue is
whether a globally defined operator needs a proper class of coefficient
exponents; adjointability gives a strong answer.

\begin{theorem}[Class extension and set support of adjoints \src{03, 06}]
\label{hgeo:thm:class}
In this class formulation the following hold.
\begin{enumerate}[label=\textup{(\roman*)}]
\item The inner product, vector norm, and all support-admissible
set-indexed strong sums are well defined over $\No[i]$.
\item Every everywhere-defined surcomplex-linear map
$A:\mathscr H_{\No}(H)\to\mathscr H_{\No}(G)$ with an
everywhere-defined class adjoint has a single \emph{set-supported}
Hahn expansion $A=\sum_{\gamma\in S}A_\gamma t^\gamma$ in ordinary bounded
maps $A_\gamma:H\to G$, with $S\subseteq\No$ a well-ordered set. For any
set-sized ordered subgroup $\Gamma$ containing $S$, $A$ and its adjoint
restrict to the corresponding Hahn workspaces.
\item Every orthogonal projection, every amplified-graph criterion
of \cref{hgeo:thm:threshold}, and every Fredholm-residue
Moore--Penrose construction above applies on these full vector
classes. Each resulting operator has its coefficients in one
set-sized Hahn workspace, although it acts on vectors from all
workspaces.
\end{enumerate}
\end{theorem}
\begin{proof}
For (i), each operand has set support. The sum, product, and inner
product use only finite unions, support sums, and finite coefficient
convolutions of those sets. Their validity is the already proved
set-sized argument in the subgroup they generate. A set-indexed
admissible family has set union of supports by replacement and union;
the same reduction applies to its strong sum. The positive-root
construction for norms also has set support.

For (ii), restrict $A$ to the set $H$ of constant vectors. Class
replacement says that the image $\{Au:u\in H\}$ is a set. Therefore
\[
 S_A=\bigcup_{u\in H}\supp(Au)
\]
is a set of surreal exponents. Do the same for the adjoint restricted
to $G$. Extract coefficients $A_\gamma u=\coeff{\gamma}{A(u)}$ as in
\cref{hgeo:thm:adjoint}; the ordinary closed graph argument makes them
bounded, and \cref{hgeo:lem:baire} makes the nonzero coefficient support well
ordered; it is a subset of the set $S_A$, and the adjoint coefficients are
the $A_\gamma^*$, with the same support. That proof only uses the two set
coefficient spaces and the support sets just exhibited. The series defines
an operator $\widetilde A$ on every global vector, because its support and
that vector's support fit in one set-sized subgroup, where Hahn convolution
is defined. Finally, for arbitrary $x$ and constant $w\in G$,
$\ip{Ax}{w}=\ip{x}{A^*w}=\ip{\widetilde Ax}{w}$, and testing against constants
identifies $A=\widetilde A$ on every class vector. This proves one
set-supported expansion for the entire operator, not a separate expansion
for each input; convolution with support in $\Gamma$ preserves the
workspace, which proves the restriction assertion.

For (iii), an orthogonal projection is everywhere adjointable, so
(ii) applies and the integral-projection and direct-rotation proofs
use its set support. For a given surcomplex scalar $a$ choose a
set-sized group containing its normal-form support; the residue
proof of the graph threshold applies without alteration, as does
its explicit inverse in the closed-range case. For a positive-order
adjointable perturbation $E$, part (ii) places all its coefficients
in one set-sized group. The Schur and Moore--Penrose formulas can
be formed there. Their identities then hold on every class vector:
include that vector's support in a larger set-sized group and apply
\cref{hgeo:prop:extension}. Necessity in the graph obstruction
also persists, since any hypothetical global projection would have
an ordinarily closed residue by the same argument as before.
\end{proof}

Source~06 states the descent (its Theorem~11.1) in NBG \emph{with} global
choice. Its proof, like source~03's, uses choice only on the sets $H$ and
$G$ (the dependent choice in \cref{hgeo:lem:baire} and the ordinary closed
graph theorem), together with class replacement; global choice is not
needed, and the theorem is stated here without it. The quantifier is over a
\emph{given} class map and its given class adjoint; it does not require a
set containing all class maps, and Zorn's lemma is not applied to a proper
class of subspaces. The ordinary set of constants is crucial: it converts
initially class-valued operator data into one set of possible exponents
before the support argument begins.

\begin{corollary}[Global geometric classification and persistent defects
\src{06}]\label{hgeo:cor:globalgeometry}
Every global orthogonal projection and every global unitary descends to a
set-sized Hahn workspace. The residue, graph and unitary classification of
\cref{hgeo:sec:splitting,hgeo:sec:unitary} holds over literal $\No$ and
$\No[i]$. The unbounded-graph counterexamples, the finite nonlattice
obstruction, the missing countable join, and the metric rigidity criterion
also hold globally.
\end{corollary}
\begin{proof}
An orthogonal projection is its own everywhere-defined adjoint; a unitary
has its inverse as adjoint. Apply \cref{hgeo:thm:class}(ii) and perform the
finite algebraic and positive-order formal constructions in a workspace
containing the resulting support. For unbounded graphs the adjoint-domain
coefficient proof is unchanged globally. If such a graph had a global
orthogonal projection, its ordinary residue would be closed, contradicting
the same domain computation. The finite-meet and countable-join proofs use
only that residue fact and the existence of rank-one projections, so they
also apply without alteration. In the metric proof, a hypothetical bijective
isometry again has adjoint $\Theta J^{-1}$; use \cref{hgeo:thm:class}(ii) in
place of \cref{hgeo:thm:adjoint} and repeat the residue argument.
\end{proof}

If a global split subspace is $M=\ran P$ and $\Gamma$ contains $\supp P$, then
\begin{equation}\label{hgeo:eq:stageintersection}
 M\cap\HH{H}=\ran\bigl(P|_{\HH{H}}\bigr).
\end{equation}
The nontrivial inclusion follows from $m=Pm$ for $m\in M$. Consequently the
global range is the compatible union of its workspace ranges \src{06}.

This extension explains why the operator theorems genuinely concern
surcomplex scalars, not only a convenient notation for ordinary
series. It does not turn the global vector class into an ordinary
set-sized Hilbert space. Nor is a ``class of all subspace
classes'' formed. Orthogonal projections can instead be represented by
their set Hahn codes; the non-meet example of
\cref{hgeo:thm:lattice} still obstructs the corresponding
projection-code order in infinite coefficient dimension.

\subsection{Global fine topology: set-sized discreteness}
On a global space use the fine norm neighborhoods
$\{x:\norm{x-a}<\varepsilon\}$ for all positive surreal $\varepsilon$.

\begin{proposition}[Set-sized discreteness in the fine topology
\src{03, 06}]\label{hgeo:prop:globalnets}
In $\mathscr H_{\No}(H)$, every set of vectors is closed and discrete in the
fine norm topology, and its distinct points have a common positive surreal
separation bound. A net indexed by a set that converges in the fine topology
is eventually equal to its limit, and one that is Cauchy for \emph{all}
positive surreal radii is eventually constant. Each fixed Hahn workspace,
viewed as a subset of the global space, is therefore discrete in the induced
fine topology.
\end{proposition}
\begin{proof}
This is the vector form of the analysis report's
\replabel{a:thm:discrete}, and of the foundations report's
\replabel{found:thm:discrete}, whose surcomplex clauses are formalized in
Lean according to \fn{docs/FORMALIZATION.md}; neither source cites them. The
mechanism is the surreal cut property. For a set $X$ of vectors, the nonzero
distances between distinct members form a set $D\subseteq\No_{>0}$, and the
cut $\{0\mid D\}$ gives $0<\varepsilon<d$ for every $d\in D$ (any positive
number if $D$ is empty). Thus each point of $X$ is isolated. For $a\notin X$,
apply the same argument to $\{\norm{a-x}:x\in X\}$; a ball around $a$ misses
$X$, so $X$ is closed. For a set-indexed net converging to $a$, apply the
separation argument to its range together with $a$: eventual membership in
the resulting ball around $a$ is eventual equality to $a$. For a Cauchy net,
the nonzero pairwise differences form a set of vectors; choose a positive
radius below all their norms, and the Cauchy condition at that radius makes
a tail constant \src{03}. A workspace is a set, so the last assertion
follows.
\end{proof}

When $\Gamma\ne0$ and $H\ne0$, the \emph{intrinsic} valuation topology on
$\HH{H}$ is not discrete: every valuation neighborhood of zero contains
nonzero monomial vectors, and nonconstant sequences such as $t^n\to0$
converge inside $\C((t^\Q))$. Those sequences only pass the \emph{workspace}
tests; in the full surreal field, a radius smaller than all their nonzero
terms prevents convergence. Thus the spherically complete set-valued model,
the class-valued fine model, and strong formal summation should not be
conflated. The global operator theorem is algebraic; it neither asserts
useful set-indexed fine convergence nor needs it. No claim is made about
arbitrary class-indexed Cauchy families.

\section{The report in the collection, the literature, and open questions}
\label{hgeo:sec:provenance}

\subsection{Relation to neighbouring reports}\label{hgeo:sub:neighbours}
The following repository results are predecessors, not contributions of this
report.
\begin{itemize}[leftmargin=*]
\item \emph{Infinite-Dimensional Hahn Spectral Theory}
(\fn{docs/surcomplex/infinite-dimensional-hahn-spectral-theory}). Part~I
constructs $H((t^\Gamma))$ and proves the positive inner product
(\replabel{ihs:hh:prop:inner}), the automatic adjoint theorem and bound
(\replabel{ihs:hh:thm:adjoint}, \replabel{ihs:hh:cor:autobounded}), the
norm-attainment criterion (\replabel{ihs:hh:thm:norm}), positive-order defect
rigidity (\replabel{ihs:hh:thm:defect}), the coercive operator $D^2+s^2I$
(\replabel{ihs:hh:thm:coercive}) and a closed proper range with zero
orthogonal complement (\replabel{ihs:hh:thm:closedrange}). Part~IV proves
the support-controlled direct rotation (\replabel{ihs:fr:lem:rotation}) and the
Hahn Fredholm alternative (\replabel{ihs:fr:thm:fredholm}). This report uses
the first group as background, derives its normal form from the rotation lemma
(\cref{hgeo:thm:normalform}), and contains the determinant-free clauses of the
Fredholm alternative (\cref{hgeo:rem:fredholm}). It does not reproduce the
spectral classifications, claim a new discovery of nonrepresented continuous
functionals, or rename a Fredholm determinant as a Moore--Penrose inverse.
The spectral report's Part~I remarks that one ``could change the vector
construction \ldots\ or impose a stronger orthogonal-completeness axiom'';
by Sol\`er's theorem (\cref{hgeo:sub:literature}) the stronger axiom cannot be
orthomodularity itself for a set-sized space over $\K$ with an infinite
orthonormal sequence, and \cref{hgeo:thm:unboundedgraph} exhibits an
explicit orthoclosed witness. No construction is changed here. The missing
countable join (\cref{hgeo:thm:nojoin}) does not contradict that report's
complete Boolean algebra of coordinate projections
(\replabel{ihs:rf:thm:commutant}), which concerns constant coordinate
projections in its row- and column-finite model; the $P_n$ here are not
constant. Source~07 answers the first further question of Part~IV
(\replabel{ihs:fr:sec:questions}) in exact-criterion form
(\cref{hgeo:us:sub:status}); Part~II's row- and column-finite calculus is a
different category from the bounded-coefficient one used there
(\cref{hgeo:us:sub:why}).
\item \emph{Three Duals at Surreal Scales}
(\fn{docs/surcomplex/three-duals-of-hahn-vector-spaces}): spherical
completeness (\replabel{duals:prop:spherical}), the coefficientwise Riesz
criterion (\replabel{duals:thm:riesz}), the closed hyperplane with zero
orthogonal complement (\replabel{duals:thm:hyperplane}) and the exact
distance cut (\replabel{duals:thm:distance}), which
\cref{hgeo:prop:distancecut} repeats for an orthoclosed subspace. That report
also distinguishes algebraic scalar extension, valuation completion and
strong closure.
\item \emph{Spectral theory} (\fn{docs/surcomplex/spectral-theory}), the
finite report: finite Moore--Penrose inverses and least squares
(\replabel{thm:leastsquares}), singular scales as determinantal scales
(\replabel{thm:scales}, \replabel{lem:DeltaInvariant}) and the SVD over a
nondivisible Hahn field (\replabel{spec:thm:svd}); \cref{hgeo:lem:smith} and
the finite part of \cref{hgeo:sec:schur} are these results
\cite{RepoFinite}. Its non-claim that it has ``no theorem \ldots\ about
infinite-dimensional Hilbert spaces over $\No[i]$'' is continued by the
spectral report and by this one.
\item \emph{Analysis} and \emph{Foundations}: fine discreteness of sets
(\replabel{a:thm:discrete}, \replabel{found:thm:discrete}), of which
\cref{hgeo:prop:globalnets} is the vector form \cite{RepoAnalysis,RepoFound}.
\item \emph{Hidden negative directions}
(\fn{docs/surcomplex/hidden-negative-hermitian-directions}): its two-scale
positivity criterion for $P+\varepsilon Q$ is related in spirit to
\cref{hgeo:thm:metric}; neither report answers a question of the other.
\end{itemize}
The relevant repository portions were inspected by the sources, not the
complete proof corpus: source~03 read the relevant parts of the spectral and
duals reports; source~06 read the root README, the opening of the
formalization ledger, the spectral report's README and article opening, and
the duals report's README; source~07 read the catalogues and the guides of
this report, the spectral report and two omnific reports
(\cref{hgeo:us:app:provenance}). No source claims inherited Lean
verification, and \fn{docs/FORMALIZATION.md} maps no label of this report to
Lean.

\subsection{Primary literature and the novelty boundary}
\label{hgeo:sub:literature}
Both literature searches were targeted, not exhaustive. Source~03 searched
the phrases ``Hahn Hilbert orthogonal graph'', ``Hahn Moore--Penrose'',
``Hahn Fredholm least squares'', non-Archimedean orthogonal complementation
and the works below; source~06 searched Hahn/Hilbert orthogonal graphs,
projection lattices and isometric positive forms; source~07's check was also
targeted (\cref{hgeo:us:app:provenance}). Failure to identify an
exact predecessor is evidence for a candidate contribution, not a proof that
none exists under other terminology.

\emph{Sol\`er's theorem} \src{06}. For a Hermitian space $E$ call
$M\subseteq E$ orthoclosed if $M=M^{\perp\perp}$, and call $E$ orthomodular
if $E=M\oplus M^\perp$ for every orthoclosed $M$. Sol\`er's theorem says
that an orthomodular Hermitian space with an infinite orthonormal sequence is
a classical real, complex or quaternionic Hilbert space; see Holland's
exposition, Theorem~1.3 \cite{Holland,Soler}. Consequently a nonclassical
scalar field such as $\F$ or $\K$ cannot retain this entire package. The
theorem applies to set-sized scalar fields and spaces; it is not applied here
to proper classes. An abstract prediction that some Hilbert axiom must fail is
therefore not an original contribution; the proposed contribution is the
explicit infinitesimal graph mechanism, its exact boundedness criterion, its
polynomial positive-kernel realization, and the concrete failure of binary
lattice operations and of a countable join. Modern categorical consequences
of Sol\`er's theorem are actively studied; for example Lack and Tobin
\cite{LackTobin} characterize the category of ordinary Hilbert spaces.

\emph{Non-Archimedean Hilbert-like spaces.} Aguayo--Nova study
non-Archimedean Hilbert-like $c_0$ spaces and orthogonal-complementation
obstructions \cite{AN}. Their model is not the layered coefficient space
retaining all of an ordinary Hilbert $H$. Source~03's comparison used the
publication record and available author-posted text; a blocked publisher
download was not treated as a completed full-text audit. Ishizuka
\cite{Ishizuka} studies compactoid self-adjoint spectral theory over a
non-Archimedean field with formally real residue field; those hypotheses and
that Banach-space category differ from the coefficient Hahn space with an
order-valued inner product, and a result in either setting is not
automatically a theorem in the other. Neither source claims a first
non-Archimedean Hilbert space, a first failure of Riesz representation, or a
first non-Archimedean spectral theorem.

\emph{Projections and perturbations.} Bursztyn--Waldmann develop formal
deformations of Hermitian vector bundles and finite projective modules,
including uniqueness up to isometric equivalence \cite{BW}. Simon discusses
unitaries relating projections and records the classical direct-rotation
formula \cite{Simon,SimonKato}; Kato's book is the standard source for the
two-projection theory and for closed operators \cite{Kato}. Neither
projection deformation nor that formula is claimed to originate here.
Corach and Maestripieri relate products of projections to polar
decompositions \cite{CorachMaestripieri}; source~07 credits them and Kato for
projection geometry and polar decomposition, and claims neither.
Avrachenkov--Lasserre investigate analytic perturbations of generalized
inverses \cite{Avrachenkov}, and M\"uller--Strohmaier develop
Hahn-meromorphic functions and a holomorphic Fredholm theorem \cite{MS}. They
are important analytic precedents. The present proofs use instead
unrestricted well-ordered supports over an arbitrary ordered group,
finite-defect orthogonal geometry, and exact leading valuation from minors;
their analyticity and convergence hypotheses are not imported, and neither
their methods nor classical generalized inverses are presented as new.
Source~03's comparison with \cite{Avrachenkov} used its abstract and the
coauthor's publication description, not the full text.

\emph{Foundations.} Conway supplies the normal-form and small-cut
foundations \cite{Conway}; Neumann supplies the classical support mechanism
\cite{Neumann}. Formal infinite-sum spaces and strong linearity have an
established contemporary literature, including \cite{Bagayoko,Freni}; those
constructions and that terminology are not new here.

\begin{table}[!ht]
\centering\small
\caption{Attribution and the boundary of the proposed contributions, as
stated by the sources.}\label{hgeo:tab:attribution}
\begin{tabularx}{\textwidth}{@{}>{\raggedright\arraybackslash}p{.25\textwidth}YY@{}}
\toprule
Result here & Prior input & Candidate contribution or status\\
\midrule
Construction, positivity, completeness, automatic adjoints, norm bounds
 & Spectral and duals reports
 & Reproved background; not new\\
Residue normal form, graph chart
 & Automatic adjoints; direct rotation (\replabel{ihs:fr:lem:rotation})
 & Global classification of split subspaces in this model (03, 06)\\
Unitary group, orbits, finite straightening
 & Formal exponential calculus; two-projection intertwiner
 & Exact statements in this Hilbert setting (06)\\
Amplified graphs
 & Normal equations; ordinary closed-range inverse
 & Exact finite/infinite threshold and residue obstruction (03)\\
Infinitesimal unbounded graphs; positive kernel; countable join
 & Closed-operator adjoints; Sol\`er's theorem
 & Explicit orthoclosed mechanism and its consequences (06)\\
Scaled closed relations
 & The two graph theorems
 & Added by the merge\\
Split-subspace poset
 & Orthomodular identities; finite projection theorem
 & Lattice if and only if coefficient dimension is finite (03; one direction 06)\\
Metric rigidity
 & Coercive operator of \replabel{ihs:hh:thm:coercive}
 & Exact ordinary-residue threshold (06)\\
Fredholm least squares
 & Schur reduction; finite Penrose geometry (\replabel{thm:leastsquares})
 & Arbitrary-rank existence and exact infinite-space solutions (03)\\
Sharp inverse valuation
 & Finite valuation-ring elimination; \replabel{thm:scales}
 & Exact finite-Schur minor formula for the full pseudoinverse (03)\\
Full surcomplex extension
 & Class replacement; normal forms; automatic adjoints
 & Explicit set support of global adjointable operators (03, 06)\\
Infinite simultaneous straightening, atomic calculus, joins
 & Polar decomposition and projection geometry
 \cite{Kato,CorachMaestripieri}; support calculus;
 \cref{hgeo:thm:finite,hgeo:sub:countable}
 & Exact criterion, jets, calculus equivalence, uniformly controlled
 counterexample (07)\\
Omnific matrix rigidity
 & Omnific normal forms \cite{LInnocenteMantova}
 & Elementary boundary corollary (07)\\
\bottomrule
\end{tabularx}
\end{table}

Several positive components are natural extensions of classical formulas, so
their priority should be assessed at the level of the precise hypotheses and
combined classification, not the formula alone. Independent mathematical
review and a wider bibliographic check remain necessary before any
publication claim of priority.

\subsection{Open questions}\label{hgeo:sub:questions}
The sources' questions are collected here, with what the other source
settles. None is a named published problem. Notes marked ``Status (batch 34)''
record what source~07 settles; its own questions are in
\cref{hgeo:us:sub:questions}.

\begin{question}[Beyond Fredholm residues \src{03}]\label{hgeo:q:mp}
The Fredholm assumption is sufficient, not asserted necessary, for an
adjointable Moore--Penrose inverse. Which operators in $\BB{H}{G}$ have one?
A broader characterization would have to treat infinitely many interacting
defect directions without assuming that their inverses share a well-ordered
support of bounded coefficient operators; the graph threshold
(\cref{hgeo:thm:threshold}) is a precise obstruction any such
characterization must explain.
\end{question}

\begin{question}[Orthoclosed subspaces with closed residue \src{03, 06}]
\label{hgeo:q:orthoclosed}
Source~03 asks when an independently presented valuation-closed subspace
admits a positive-order graph chart, and makes no claim that closed residue
suffices. \Cref{hgeo:prop:closedresidue} (completing a remark of source~06)
settles that negatively for valuation-closed subspaces, with a subspace that
is not orthoclosed. For orthoclosed subspaces the question is open; within
the scaled relations of \cref{hgeo:thm:relation}, closed residue is
equivalent to splitting. Source~06 asks, more specifically, which kernels of
regular Hahn operators have closed residues and which higher-order conditions
make them complemented; the graph chart gives a local model, but closed
residue is not asserted to settle every orthoclosed subspace.
\end{question}

\noindent\emph{Status (batch 34).} Source~07 restates this question
(\cref{hgeo:us:q:closedresidue}) and does not settle it; it stays open.

\begin{question}[Infinite orthogonal synthesis \src{06}]\label{hgeo:q:synthesis}
For infinite orthogonal families, a coherent common support and bounded
ordinary coefficient synthesis are separate necessary conditions
(\cref{hgeo:cor:nosimultaneous,hgeo:prop:incoherent}). Which natural
sufficient conditions, stable under composition and scale enlargement, make
a countable family simultaneously straightenable, or give it a join?
\end{question}

\noindent\emph{Status (batch 34).} Answered by source~07. For any set-indexed
orthogonal family, complete or not, over any set-sized ordered group, a
simultaneous unitary exists if and only if the canonical columns are bounded
at every exponent (B) and their active support is well ordered (H)
(\cref{hgeo:us:thm:main}); solvability is unchanged by every ordered scale
enlargement (\cref{hgeo:us:cor:extension}), and families satisfying the
criterion have joins (\cref{hgeo:us:prop:joins,hgeo:us:rem:joins}). The
criterion is exact, not an effective test, and stability under composition is
not treated separately (\cref{hgeo:us:sub:status}).

This bears on the first further question of the spectral report's Part~IV
(\replabel{ihs:fr:sec:questions}): what uniform spectral-gap and
coefficient-ideal hypotheses ensure that all cluster rotations assemble into
a bounded-coefficient Hahn unitary? Assembling the cluster rotations into one
unitary would conjugate all the cluster projections simultaneously, so the
two conditions above are necessary for it; \cref{hgeo:cor:nosimultaneous}
and \cref{hgeo:prop:incoherent} give families for which each fails. Sources~03
and~06 prove no sufficient hypothesis. \emph{Status (batch 34):} source~07's
criterion answers that question in exact-criterion form, for the family of
cluster projections, but not in the spectral-gap and coefficient-ideal form in
which it is posed; in that form it stays open (\cref{hgeo:us:sub:status}).
Its other two parts (regularized spectral products; compact residues with a
kernel) are addressed by no source.

\begin{question}[Singular positive forms \src{06}]\label{hgeo:q:metric}
\Cref{hgeo:thm:metric} suggests a stratification of positive Hahn metrics by
the ordinary lower-bound behaviour of their leading coefficient, with more
singular forms treated by a different coefficient space rather than by a
forced canonical isometry.
\end{question}

\begin{question}[Structures not developed \src{03}]\label{hgeo:q:notdeveloped}
Source~03 lists general tensor products, frame bounds, spectral measures,
unbounded operators with proper domains, and infinite-family projection
joins as not developed. Within this report source~06 treats two of them:
graphs of closed unbounded operators (\cref{hgeo:thm:unboundedgraph}), and
countable projection joins, negatively (\cref{hgeo:thm:nojoin}: a countable
orthogonal family need not have a join). Which countable families do have
joins is not decided. Tensor products, frame bounds and spectral measures
remain undeveloped; by \cref{hgeo:thm:nojoin} no construction can assign the
join of \emph{every} countable orthogonal family, but this does not forbid
spectral measures for particular operators or other coefficientwise
summation rules.
\end{question}

\noindent\emph{Status (batch 34).} Answered in part by source~07: every
family satisfying the criterion of \cref{hgeo:us:thm:main}, and each of its
subfamilies, has a join (\cref{hgeo:us:prop:joins,hgeo:us:rem:joins}), and
\cref{hgeo:us:thm:uniform} gives a further family, complete and with uniformly
bounded coefficients, with a subfamily that has no join. A characterization of
the countable families with joins is not given (\cref{hgeo:us:q:boolean}).

For computational use, the finite Schur matrix is exact but its entries may
themselves be infinite Hahn series. The minor formula specifies an invariant
and an algebraic reduction; it is not an algorithm for deciding equality or
leading terms of arbitrary uncomputably presented surreal numbers. Effective
implementations need an explicit support and coefficient representation.

\subsection{Proof status}
All central arguments are given in full above or in the cited reports. Their
main dependencies are the ordinary Hilbert closed-graph, Baire,
bounded-inverse and Riesz theorems, the adjoint theorem for closed densely
defined operators, ordinary positive square roots, the Hahn--Neumann support
lemma, Conway normal forms, and finite linear algebra; these are identified
inputs, not proved from set theory here. The class section additionally uses
class replacement in NBG without global choice.
\Cref{hgeo:us:sec:intro,hgeo:us:sec:criterion,hgeo:us:sec:consequences,%
hgeo:us:sec:calculus,hgeo:us:sec:uniform,hgeo:us:sec:support,%
hgeo:us:sec:surreal,hgeo:us:sec:questions} use in addition ordinary bounded
extension from dense subspaces, the diagonal representation of
$\ell^\infty$ on a Hilbert direct sum, and the matrix-amplification estimate
for completely bounded maps. No proof assistant was run,
no source has been independently refereed, and no assertion rests on
numerical evidence. The mathematical claims should be evaluated by checking
the proofs and their hypotheses, not the novelty wording.

\section{Infinite simultaneous straightening: the source, its results and
conventions}
\label{hgeo:us:sec:intro}

\Cref{hgeo:us:sec:intro,hgeo:us:sec:criterion,hgeo:us:sec:consequences,%
hgeo:us:sec:calculus,hgeo:us:sec:uniform,hgeo:us:sec:support,%
hgeo:us:sec:surreal,hgeo:us:sec:questions} add a third manuscript,
source~07, to the report. They are placed after \cref{hgeo:sec:provenance}
so that no earlier number changes. Source~07 was written against this
report's guide and answers its \cref{hgeo:q:synthesis}: for a set-indexed
orthogonal family of Hahn projections it gives an exact criterion for one
unitary conjugating every member back to its residue.

\subsection{Source 07 and what it settles}\label{hgeo:us:sub:source}
Source~07 is manuscript~09 of the thirty-fourth batch, placed in commit
\texttt{a7a435f} from the archive \fn{surreal_unitary_synthesis} (delivered
in \texttt{a4611ad}). It continues this report's local file sequence after
\fn{03-} and \fn{06-}, so its shipped files carry the prefix
\fn{07-unitary-synthesis-}.

\begin{center}
\small
\renewcommand{\arraystretch}{1.15}
\begin{tabular}{>{\raggedright\arraybackslash}p{0.07\textwidth}>{\raggedright\arraybackslash}p{0.34\textwidth}>{\raggedright\arraybackslash}p{0.49\textwidth}}
\toprule
Source & Manuscript & Contribution to this report\\
\midrule
07 & \emph{Simultaneous Unitary Straightening at Surreal Scales: An exact
synthesis criterion, atomic spectral calculus, and uniformly controlled
counterexamples} (22 pp., dated September 23, 2026; pin
\texttt{fc27d87}) &
The exact criterion (B)$+$(H) for simultaneous straightening, with an
explicit unitary for complete and incomplete residual families; no new
scales, gauge classification, scale-enlargement invariance; the finite-jet
theorem and the first-order equation; the equivalent atomic functional
calculus, joins and weighted reconstruction; a complete family with uniform
coefficient bounds at every order and no straightener, with a subfamily
lacking a join; bounded columns with incoherent support; surreal descent;
omnific matrix rigidity; twelve questions.\\
\bottomrule
\end{tabular}
\end{center}

The pin \texttt{fc27d87} (committed September 23, 2026 at 17:27:33 in
America/Los\_Angeles) is 61 commits before the placement. Between the pin
and the placement neither this report nor the infinite-dimensional spectral
report changed, except that the placement added source~07's files here. The
report numbers source~07 quotes (Question~15.3, Part~IV of the spectral
report) are therefore still correct.

\paragraph{What source~07 settles here.}
\begin{itemize}[leftmargin=*]
\item \Cref{hgeo:q:synthesis} is answered: the two necessary conditions of
\cref{hgeo:sub:countable}, in the precise form of bounded synthesized
coefficients (B) and well-ordered active support (H), are together
sufficient (\cref{hgeo:us:thm:main}). Families satisfying them have joins
(\cref{hgeo:us:prop:joins,hgeo:us:rem:joins}).
\item \Cref{hgeo:q:notdeveloped} is answered in part: a sufficient condition
for joins and a further family with a subfamily lacking a join
(\cref{hgeo:us:thm:uniform}). A characterization of the countable families
with joins remains open.
\item The first further question of the spectral report's Part~IV
(\replabel{ihs:fr:sec:questions}) is answered in exact-criterion form, not in
the spectral-gap form in which it is posed (\cref{hgeo:us:sub:status}).
\item \Cref{hgeo:q:orthoclosed} stays open; source~07 says so itself
(\cref{hgeo:us:q:closedresidue}).
\end{itemize}

\paragraph{Where the merge had to choose.}
\begin{itemize}[leftmargin=*]
\item \emph{Printed once.} Source~07's integrality lemma is
\cref{hgeo:lem:integralproj} together with the unitary residue argument of
\cref{hgeo:cor:orbits}; its two small additions are kept
(\cref{hgeo:us:lem:integral}). Its support principle and formal functions
are \cref{hgeo:lem:support} and \eqref{hgeo:eq:formalcalculus}; its strong
Hahn summability is \cref{hgeo:def:strong}; its substitution
$t^\gamma=\omega^{-\gamma}$ is \eqref{hgeo:eq:embedding}. Its finite-family
corollary is \cref{hgeo:thm:finite}, re-derived independently from the
criterion (\cref{hgeo:us:rem:finite}).
\item \emph{Second routes kept.} The necessity direction of
\cref{hgeo:us:thm:main} contains \cref{hgeo:cor:nosimultaneous} and
\cref{hgeo:prop:incoherent} (\cref{hgeo:us:rem:instances}); source~06's
direct proofs, which argue from the support of a hypothetical unitary, are
kept. \Cref{hgeo:us:thm:support} contains the family of
\cref{hgeo:prop:incoherent}; both nonexistence proofs are printed.
\item \emph{Added by the merge} \mergetag, with complete proofs: the
identification of the construction with \cref{hgeo:thm:finite} in the finite
case (\cref{hgeo:us:rem:finite}); joins without completeness or complex
scalars, and the range projection as the join (\cref{hgeo:us:rem:joins});
the class-map form of the surreal descent, from \cref{hgeo:thm:class}
(\cref{hgeo:us:cor:classmaps}); the transport of \cref{hgeo:us:thm:support}
to every group containing $\Q\eta$ (\cref{hgeo:us:rem:transport}); the status
notes of \cref{hgeo:us:sub:status} and in \cref{hgeo:sub:questions}.
\item \emph{Renamed symbols.} Source~07's letters meet the report's reserved
letters in many places; \cref{hgeo:us:sub:conventions} and
\cref{hgeo:tab:renames} list every renaming.
\item \emph{A stale scope statement.} Source~07 deliberately does not use an
automatic class-operator representation theorem of a neighbouring report
(it means \cref{hgeo:thm:class}), and asks for one
(\cref{hgeo:us:q:classmaps}). For maps with an everywhere-defined class
adjoint that theorem is proved in this report, which gives
\cref{hgeo:us:cor:classmaps}.
\end{itemize}
Source~07's limitations are kept in \cref{hgeo:app:nonclaims}, its
provenance appendix is printed in \cref{hgeo:us:app:provenance}, and its suite
is described in \cref{hgeo:app:verification}.

\subsection{Why an infinite family is different}\label{hgeo:us:sub:why}
Let $(E_i)_{i\in\mathcal I}$ be an orthogonal family of ordinary projections
on $H$, and let $P_i$ be Hahn projections with residues $E_i$. For finitely
many $i$ one near-identity unitary removes the infinitesimal perturbation
(\cref{hgeo:thm:finite}); for infinitely many it need not
(\cref{hgeo:sub:countable}). The difficulty is not the algebraic polar
formula. It is whether the would-be intertwiner is an element of the operator
algebra at all \src{07}.

Two logically different failures occur. At a fixed exponent, infinitely many
harmless-looking columns may assemble into an unbounded operator. Across
exponents, individually legal Hahn supports may have a union that is not well
ordered. Neither failure is visible on finitely many projections. Source~07
turns these observations into an exact and constructive equivalence, rather
than a list of necessary conditions.

The setting is the concrete space $\HH{H}=H((t^\Gamma))$ with the concrete
algebra $\BH=\mathcal B(H)((t^\Gamma))$, not a Banach space over a
non-Archimedean field with an unspecified completion. Every coefficient is an
ordinary bounded operator, and membership in the algebra is part of each
theorem, not a suppressed convention.

This is not a repetition of Part~II of the spectral report
\cite{RepoSpectral}. That part uses row- and column-finite matrices on series
with finitely supported coordinate coefficients. Here the coefficient space
is the whole ordinary Hilbert space and the coefficients must be bounded:
infinite rows and columns are allowed when they define bounded operators, and
unbounded synthesis is forbidden. The two categories have different
infinite-dimensional obstructions.

The report already contains the projection-lattice counterexamples and the
support-incoherence obstruction (\cref{hgeo:thm:nojoin,hgeo:cor:nosimultaneous,%
hgeo:prop:incoherent}); source~07 does not claim either phenomenon for the
first time. Its strengthened coefficient example imposes simultaneously a
common discrete support and uniform bounds on every individual coefficient
at every order, and isolates the collective synthesis obstruction that these
stronger assumptions still fail to control.

Source~07 names as its proposed contributions the exact criterion with its
incomplete-family construction, invariance of solvability under scale
enlargement, the finite-jet equivalence, the atomic functional-calculus
equivalence, and the uniformly controlled counterexample with its explicit
missing join. They are candidate contributions relative to the inspected
sources, not a certified exhaustive novelty claim. Polar intertwiners are
classical; projection geometry and polar decomposition are developed, for
example, by Kato and by Corach--Maestripieri \cite{Kato,CorachMaestripieri};
well-ordered support calculus is Hahn--Neumann algebra \cite{Neumann}; the
omnific normal form used in \cref{hgeo:us:sec:surreal} is standard
\cite{LInnocenteMantova}. None of these is advertised as new. The proofs
supply the infinite-family membership and support arguments that a formal
appeal to polar decomposition would leave out. Source~07 needs no unreviewed
theorem of the collection: the report guides fixed its target and its
comparison, and its argument starts from ordinary Hilbert operator theory and
the support lemma.

\subsection{Conventions and renamed symbols of this part}
\label{hgeo:us:sub:conventions}
The following conventions hold in
\cref{hgeo:us:sec:intro,hgeo:us:sec:criterion,hgeo:us:sec:consequences,%
hgeo:us:sec:calculus,hgeo:us:sec:uniform,hgeo:us:sec:support,%
hgeo:us:sec:surreal,hgeo:us:sec:questions}. Source~07's notation is
renamed where it meets a letter the report already uses; each renaming is
also listed in \cref{hgeo:tab:renames}.
\begin{itemize}
\item \emph{Scalars and spaces.} $\kk\in\{\R,\C\}$ is the coefficient field of
$H$ (07: $\mathbb F$), and the Hahn scalars are $\F$ or $\K$ as in
\cref{hgeo:sub:conventions} (07: $K=\mathbb F((t^\Gamma))$).
The space is $\HH{H}$ (07: $\mathscr H$) and the algebra is $\BH$ (07:
$\mathscr A$); its integral and positive-order parts are $\Bnn$ and
$\Bpos{H}{H}$ (07: $\mathscr A_{\ge0}$, $\mathscr A_{>0}$). The identity is
$I$ (07: $I_H$). Inner products are linear in the first argument, as in the
rest of the report; source~07's are conjugate-linear in the first, and its
two formulas that depend on this (rank-one maps and matrix entries) are
converted.
\item \emph{Indices.} The family is indexed by a set $\mathcal I$ (07: $I$,
which here is the identity); finite subsets are $\mathcal J$ (07: $J$), Boolean
subsets $\mathcal S$ (07: $A,B$), the left-half subfamily of
\cref{hgeo:us:thm:uniform} is $\mathcal L$ (07: $J$), and the index set of an
orthonormal basis is $\mathcal N$ (07: $J$). The letter $J$ keeps its meanings
of \cref{hgeo:cor:graphchart,hgeo:thm:metric}.
\item \emph{Residues.} $E_i=(P_i)_0$, $E=\sSOT_iE_i$ and $H_E=\ran E$ are
source~07's. The letter $E$ in this part always denotes such a residue
projection, never the positive-order perturbation $E$ of
\cref{hgeo:sec:schur,hgeo:sec:least,hgeo:sec:valuation}. The residue
projections of \cref{hgeo:thm:finite,hgeo:cor:nosimultaneous}, written
$Q_j$ and $Q_n$ there, are $E_j$ here. The complementary projection is
$E^\perp=I-E$ (07: $F$, which here is the real Hahn field); it projects onto
$(H_E)^\perp$.
\item \emph{The construction.} The canonical column series $C$, its
coefficients $C_\gamma$ and $C_\gamma^\fin$, the active support $S_C$ and the
constants $b_\gamma$ are source~07's; $C$ is not the complex field $\C$, nor
the block map of the proof of \cref{hgeo:thm:adjoint}. The diagonal
compression is $\Dg_E$ (07: $\mathcal D_E$; $\mathcal D$ is the block matrix
of \cref{hgeo:sec:schur}). Source~07's $G=C^*C+F$ is $\Omega=C^*C+E^\perp$,
because $G$ is a Hilbert space here. Its range projection $Q=CG^{-1}C^*$ is
$P_\vee$: the $Q$'s of the report are residue projections, and $P_\vee$ is
indeed the join of the family (\cref{hgeo:us:rem:joins}). Its $T=C+(I-Q)F$ is
$W$: $T$ is $\diag(1/n)$ here, and in the finite case this $W$ is exactly the
$W$ of \cref{hgeo:thm:finite} (\cref{hgeo:us:rem:finite}). Its
$B=(T^*T)^{-1/2}$ is $\abs{W}^{-1}=(W^*W)^{-1/2}$, the formal binomial series
of \eqref{hgeo:eq:formalcalculus}, not an ordinary operator absolute value;
$B$ is the row operator \eqref{hgeo:eq:rowB}.
\item \emph{Other operators.} A generic positive-order series is $Y$ (07:
$R$, a closed linear relation here). The block unitaries $\exp(tA_m)$ and
$(I+q_nJ)(1+q_n^2)^{-1/2}$ are $V_m$ and $V_n$ (07: $R_m$, $R_n$), their
generator is $X_m$ (07: $A_m$; $A$ is an unbounded or Fredholm operator
here), and the $2\times2$ rotation is $\mathsf J$ (07: $J$). The monomials
$q_n=t^{1/n}$ are written out, since $q=t^\eta$ is fixed. The flat vectors
are $u_m,u_m'$ (07: $u_m,v_m$; $v$ is the valuation), and the block identity
projection is $\Pi_m$ (07: $F_m$). The bounded skew-adjoint solution of the
first-order equation is $Y$ (07: $K$, the complex Hahn field here), its
commuting part $Y_c$ (07: $D$), and a gauge unitary is $\Xi$ (07: $W$). The
jet rings are $\mathscr B^{[N]}$ (07: $\mathscr A_N$), and the weighted
operator is $N_\lambda$ (07: $A_\lambda$). A generic operator is $A$ (07:
$T$). Source~07's generated monoid $\langle S\rangle$ is $S^*$ as in
\cref{hgeo:lem:support}; it is not an adjoint.
\item \emph{Kept letters with a scoped meaning.} $\Phi,\Psi$ are atomic
calculi (not the map in the proof of \cref{hgeo:prop:closedresidue}); $S$,
$S_C$, $S_\lambda$ are supports, as in \cref{hgeo:prop:incoherent}, never the
Schur matrix $S$ of \cref{hgeo:sec:schur}; $U$ is a Hahn unitary and
$\Un$ keeps its meaning of \cref{hgeo:thm:unitary}; $m$ indexes blocks, not
$\dim\ker T_0$.
\end{itemize}
A tempting false reading: ``$P_\vee$ is a residue projection'' or ``$P_\vee$
is the formal sum of the $P_i$''. It is the Hahn projection with residue $E$
onto the range of $C$; it equals $\sum_iP_i$ only for finite families, and in
general no strong Hahn sum of the $P_i$ exists, since infinitely many of them
have nonzero coefficient at exponent zero.

\subsection{The Hahn operator algebra and three summation conventions}
\label{hgeo:us:sub:algebra}
Throughout this part $H$ is a nonzero real or complex Hilbert space and
$\Gamma$ is any set-sized ordered abelian group; it need not be divisible, and
its order need not be Archimedean. The algebra $\BH$ acts on $\HH{H}$ as in
\cref{hgeo:sec:adjoints}: multiplication is convolution, $t^\gamma$ is
central, coefficient products are compositions in the displayed order, and
$A^*=\sum_\gamma A_\gamma^*t^\gamma$. We write $v(A)=\min\supp A$,
$v(0)=+\infty$, and $A_0=\res A$ for the coefficient at zero. The general
inequality $v(AB)\ge v(A)+v(B)$ is not an equality, since operator
coefficients can multiply to zero; but $v(A^*A)=2v(A)$ for $A\ne0$
(\cref{hgeo:cor:vop}), because $A_a^*A_a\ne0$ for a nonzero bounded $A_a$.
Likewise, if $x\ne0$ has first exponent $a$, the first coefficient of
$\ip{x}{x}$ is $\norm{x_a}^2>0$, at $2a$.

Source~07 uses the support principle of \cref{hgeo:lem:support}: a fixed
group element has only finitely many representations as an ordered finite
word in a well-ordered $S\subseteq\Gamma_{>0}$, repeated copies of an exponent
counting as repeated letters. Hence every ordinary formal power series
$\sum_na_nz^n$ can be evaluated at $Y\in\Bpos{H}{H}$, each coefficient using
finitely many products of operator coefficients and no operator-norm
convergence. In particular $(I+Y)^{-1}$ and $(I+Y)^{-1/2}$ are the series of
\eqref{hgeo:eq:formalcalculus}. Identities between such series follow from the
one-variable formal power-series ring even though the coefficient algebra is
noncommutative, because every evaluated factor is a power of the same $Y$.
The inverse square root is self-adjoint when $Y$ is, and commutes with every
operator commuting with $Y$.

\begin{lemma}[Integral projections and unitaries \src{07}; see also
\src{03, 06}]\label{hgeo:us:lem:integral}
Every nonzero projection $P\in\BH$ has $v(P)=0$ and a nonzero ordinary
orthogonal projection $P_0$. Every unitary $U\in\BH$ has $v(U)=0$ and an
ordinary unitary $U_0$. A projection with residue zero is zero, and two
projections $Q\le P$ with $P_0=Q_0$ are equal.
\end{lemma}
\begin{proof}
The projection assertion is \cref{hgeo:lem:integralproj}. Source~07's proof
is source~03's route: $P=P^*P$ gives $v(P)=2v(P)$, and an ordered abelian
group has no torsion. For a unitary, $U^*U=I$ gives $2v(U)=0$, so $v(U)=0$, and
the zero coefficients of both unitary identities give $U_0^*U_0=U_0U_0^*=I$;
source~06 obtains the same from contractivity (\cref{hgeo:cor:orbits}). A
nonzero projection has nonzero residue, which gives the third assertion. For
the last, $Q\le P$ means $PQ=QP=Q$, so $P-Q$ is a projection with residue
zero.
\end{proof}

A \emph{near-identity unitary} is a unitary $U$ with $U_0=I$, equivalently
$U-I\in\Bpos{H}{H}$, including $U=I$; these form the group
$\Un_{1,\Gamma}(H)$ of \cref{hgeo:thm:unitary}.

\begin{warning}[Three summation conventions \src{07}]\label{hgeo:us:warn:sums}
Three notions of an infinite sum occur in this part and must not be
confused.
\begin{enumerate}[label=\textup{(\alph*)}]
\item \emph{Strong Hahn summability} (\cref{hgeo:def:strong}): the union of the
supports is well ordered and at each exponent only finitely many members are
nonzero; the sum is formed by finite coefficient sums.
\item A \emph{coefficientwise strong operator sum} $\sSOT$: at each exponent
separately, the ordinary strong-operator limit of the finite partial sums,
whose values converge in the norm of $H$ on every ordinary vector; infinitely
many operators may be nonzero at one exponent, and the Hahn support of the
result must then be checked. Such a family need not be strongly Hahn
summable.
\item \emph{Valuation convergence}: the first exponent of the error moves
above every prescribed scale. Ordinary convergence of each coefficient does
not imply it.
\end{enumerate}
\end{warning}

\section{The exact synthesis criterion}\label{hgeo:us:sec:criterion}

\subsection{Canonical columns and the theorem}
Let $(P_i)_{i\in\mathcal I}$ be a set-indexed family of pairwise orthogonal
projections in $\BH$ and put $E_i=(P_i)_0$. By \cref{hgeo:us:lem:integral}
the $E_i$ form an ordinary orthogonal family of projections; zero members may
be discarded. Let
\[
 H_E=\overline{\bigoplus_{i\in\mathcal I}E_iH},\qquad
 E=\sSOT_{i\in\mathcal I}E_i,\qquad E^\perp=I-E .
\]
The sum defining $E$ is an ordinary Hilbert-space sum. It is \emph{not}
required that $E=I$: incomplete residual families are allowed.

\begin{definition}[Canonical coefficient columns \src{07}]
\label{hgeo:us:def:columns}
For $\gamma\in\Gamma$ define, on the algebraic direct sum
$\bigoplus_i^{\mathrm{alg}}E_iH$, the linear map
\begin{equation}\label{hgeo:us:eq:columns}
 C_\gamma^\fin\Bigl(\sum_{i\in\mathcal J}x_i\Bigr)
 =\sum_{i\in\mathcal J}(P_i)_\gamma x_i,
 \qquad \mathcal J\subseteq\mathcal I\text{ finite},\quad x_i\in E_iH .
\end{equation}
Its \emph{active support} is $S_C=\{\gamma:C_\gamma^\fin\ne0\}$. The family
satisfies
\begin{itemize}[leftmargin=2.2em]
\item[\textup{(B)}] \emph{bounded coefficient synthesis} if for every
$\gamma$ there is a finite constant $b_\gamma$ with
\[
 \Bigl\|\sum_{i\in\mathcal J}(P_i)_\gamma x_i\Bigr\|
 \le b_\gamma\Bigl(\sum_{i\in\mathcal J}\norm{x_i}^2\Bigr)^{1/2}
 \quad\text{for all finite }\mathcal J\text{ and all }x_i\in E_iH;
\]
\item[\textup{(H)}] \emph{Hahn coherence} if $S_C$ is well ordered in
$\Gamma$.
\end{itemize}
\end{definition}

Condition (B) is equivalent to
\begin{equation}\label{hgeo:us:eq:finitesup}
 \sup_{\mathcal J\subseteq\mathcal I\text{ finite}}
 \Bigl\|\sum_{i\in\mathcal J}(P_i)_\gamma E_i\Bigr\|<\infty\qquad
 (\gamma\in\Gamma):
\end{equation}
the vectors $x_i=E_ix$ give one direction, and $x=\sum_ix_i$, whose norm is
$(\sum_i\norm{x_i}^2)^{1/2}$ by orthogonality, gives the other. Under (B),
$C_\gamma^\fin$ extends uniquely to a bounded map on $H_E$; extending by zero
on $E^\perp H$ gives $C_\gamma\in\mathcal B(H)$. The finite sums in
\eqref{hgeo:us:eq:finitesup} are $C_\gamma\sum_{i\in\mathcal J}E_i$, so they
converge strongly to $C_\gamma$. Note that
\[
 C_0=E,\qquad C_\gamma E^\perp=0,\qquad S_C\subseteq\Gamma_{\ge0}.
\]
The constants $b_\gamma$ need not be bounded uniformly in $\gamma$.

\begin{theorem}[Exact simultaneous straightening \src{07}]
\label{hgeo:us:thm:main}
For the family above the following are equivalent.
\begin{enumerate}[label=\textup{(\roman*)}]
\item There is a near-identity unitary $U\in\BH$ with $P_i=UE_iU^*$ for every
$i\in\mathcal I$.
\item There is a unitary $U\in\BH$, not necessarily near the identity, with
the same identities.
\item Conditions \textup{(B)} and \textup{(H)} both hold.
\end{enumerate}
Under these conditions let $C=\sum_\gamma C_\gamma t^\gamma\in\BH$. A
specific near-identity straightener is given by
\begin{align}
 \Omega&=C^*C+E^\perp, & P_\vee&=C\,\Omega^{-1}C^*,
 \label{hgeo:us:eq:constructiona}\\
 W&=C+(I-P_\vee)E^\perp, & U&=W\abs{W}^{-1},\quad
 \abs{W}^{-1}=(W^*W)^{-1/2}.
 \label{hgeo:us:eq:constructionb}
\end{align}
Here $\Omega_0=I$, $P_\vee$ is a projection with residue $E$, and $W_0=I$;
the inverse and the inverse square root are the formal series of
\eqref{hgeo:eq:formalcalculus}. If $E=I$, the formula reduces to
\begin{equation}\label{hgeo:us:eq:polar}
 U=C\,(C^*C)^{-1/2}.
\end{equation}
\end{theorem}

The test is exact, but it is not an effective decision procedure for
arbitrary infinite data. It replaces the unknown common unitary by explicitly
specified linear maps, one at each exponent, and one support condition.
\Cref{hgeo:us:thm:uniform,hgeo:us:thm:support} show that neither condition
can be dropped, and that bounds on the individual coefficients do not
replace (B).

\subsection{The diagonal compression}
For $A\in\mathcal B(H)$ put
\[
 \Dg_E(A)=\sSOT_{i\in\mathcal I}E_iAE_i .
\]
It acts by the diagonal block $E_iAE_i$ on $E_iH$ and by zero on $E^\perp H$.
For a finite orthogonal sum,
$\norm{\sum_iE_iAE_ix}^2=\sum_i\norm{E_iAE_ix}^2\le\norm A^2\sum_i\norm{E_ix}^2$,
so by completion $\Dg_E$ is a linear contraction of $\mathcal B(H)$. Applied
coefficient by coefficient it defines $\Dg_E$ on $\BH$ without enlarging
supports. It need not be multiplicative.

\subsection{Necessity}
Let $U$ be a near-identity straightener. The Hahn operator
\begin{equation}\label{hgeo:us:eq:necessity}
 C=U\,\Dg_E(U^*)
\end{equation}
exists, as a product of two members of $\BH$. It vanishes on $E^\perp H$, and
since $\Dg_E(U^*)E_i=E_iU^*E_i$ coefficientwise,
\[
 CE_i=UE_iU^*E_i=P_iE_i .
\]
Comparing coefficients identifies $C_\gamma$ with the bounded extension of
\eqref{hgeo:us:eq:columns}, which proves (B). The active support of the
columns is exactly $\supp C$, since a bounded coefficient vanishing on every
$E_iH$ and on $E^\perp H$ is zero; this proves (H). No infinite Hahn sum is
interchanged with a product: two bounded coefficients are identified on a
dense ordinary subspace.

If $U$ is an arbitrary straightener, its residue $U_0$ is an ordinary unitary
(\cref{hgeo:us:lem:integral}) with $U_0E_iU_0^*=E_i$, so $U_0$ commutes with
every $E_i$ and $UU_0^*$ is a near-identity straightener. This proves
(ii)$\Rightarrow$(i)$\Rightarrow$(iii).

\subsection{The intertwining identities}
Assume (B) and (H). Then $C\in\BH$, $C=CE$, $C_0=E$ and $CE_i=P_iE_i$. For
$x\in E_jH$,
\[
 P_iCx=P_iP_jx=\delta_{ij}P_jx=CE_ix .
\]
Both sides vanish on $E^\perp H$. Each coefficient of either side is a finite
sum of bounded operators, so equality on the dense subspace
$\bigoplus_j^{\mathrm{alg}}E_jH+E^\perp H$ extends to $H$:
\begin{equation}\label{hgeo:us:eq:intertwine}
 P_iC=CE_i,\qquad C^*P_i=E_iC^* .
\end{equation}
In particular
\begin{equation}\label{hgeo:us:eq:commute}
 C^*CE_i=C^*P_iC=E_iC^*C .
\end{equation}
The coefficients of $C^*C$ therefore commute with every $E_i$, hence with
their strong sum $E$ and with $E^\perp$; moreover $C^*CE^\perp=0=E^\perp C^*C$.
Thus $\Omega=C^*C+E^\perp$ commutes with every $E_i$, with $E$ and with
$E^\perp$. Since $\Omega_0=E+E^\perp=I$, it has a formal inverse in $\Bnn$,
and both $\Omega$ and $\Omega^{-1}$ are self-adjoint.

\subsection{Completing the partial range}
Define $P_\vee$ by \eqref{hgeo:us:eq:constructiona}. It is self-adjoint with
residue $E$. From $CE^\perp=0$, $E^\perp C^*=0$ and $\Omega E^\perp=E^\perp$,
hence $\Omega^{-1}E^\perp=E^\perp$,
\[
 P_\vee^2=C\Omega^{-1}(C^*C)\Omega^{-1}C^*
 =C\Omega^{-1}(\Omega-E^\perp)\Omega^{-1}C^*=P_\vee ,
\]
and the same calculation gives $P_\vee C=C$. By
\eqref{hgeo:us:eq:intertwine} and \eqref{hgeo:us:eq:commute},
\[
 P_iP_\vee=CE_i\Omega^{-1}C^*=C\Omega^{-1}E_iC^*=P_\vee P_i .
\]
Hence $P_iP_\vee$ is a projection with residue $E_iE=E_i$, and
$P_i-P_iP_\vee$ is a projection with residue zero, so
\cref{hgeo:us:lem:integral} gives
\begin{equation}\label{hgeo:us:eq:underP}
 P_iP_\vee=P_\vee P_i=P_i .
\end{equation}
This is where self-adjointness rules out a hidden positive-order remainder
projection.

Put $W=C+(I-P_\vee)E^\perp$. Its residue is $E+(I-E)E^\perp=I$, so it is
invertible by \eqref{hgeo:eq:formalcalculus}. From \eqref{hgeo:us:eq:underP}
and $E^\perp E_i=0$,
\begin{equation}\label{hgeo:us:eq:Wintertwine}
 P_iW=CE_i=WE_i .
\end{equation}
Taking adjoints and multiplying gives $E_iW^*W=W^*P_iW=W^*WE_i$. Thus
$W^*W=I+Y$ with $Y\in\Bpos{H}{H}$ self-adjoint, and
$\abs{W}^{-1}=(W^*W)^{-1/2}$ is self-adjoint and commutes with every $E_i$.
The operator $U=W\abs W^{-1}$ satisfies $U^*U=I$. Both factors are
invertible, so $U$ is invertible and $U^*U=I$ implies $UU^*=I$. Finally
\eqref{hgeo:us:eq:Wintertwine} gives $P_iU=UE_i$, hence $P_i=UE_iU^*$, and
$U_0=I$.

If $E=I$, then $C_0=I$, $C$ is invertible, $P_\vee=C(C^*C)^{-1}C^*=I$ and
$W=C$, which gives \eqref{hgeo:us:eq:polar} and completes the proof of
\cref{hgeo:us:thm:main}.\hfill$\square$

\subsection{The range projection}
The projection $P_\vee$ is the projection onto the range of $C$ in the
algebraic Hahn sense: $P_\vee C=C$ and $P_\vee=C(\Omega^{-1}C^*)$ show that
the two operators have the same range on $\HH{H}$. Moreover
$P_\vee W=C=WE$. Since $W^*W$ commutes with every $E_i$, each of its bounded
coefficients commutes with their strong sum $E$, and so does $\abs W^{-1}$.
Hence $P_\vee U=P_\vee W\abs W^{-1}=WE\abs W^{-1}=UE$, and
\begin{equation}\label{hgeo:us:eq:range}
 P_\vee=UEU^* .
\end{equation}
No ordinary topological closure of a range in a non-Archimedean completion
is used here \src{07}.

\subsection{Relation to the finite theorem and to the countable obstructions
\mergetag}\label{hgeo:us:sub:relation}

\begin{remark}[The finite case is \cref{hgeo:thm:finite} \mergetag]
\label{hgeo:us:rem:finite}
Every finite orthogonal family satisfies (B) and (H): a finite sum of bounded
coefficient maps is bounded and a finite union of well-ordered supports is
well ordered. \Cref{hgeo:us:thm:main} therefore re-derives
\cref{hgeo:thm:finite} \src{07}. The two constructions agree. For a finite
family, \eqref{hgeo:us:eq:range} gives
$P_\vee=\sum_jUE_jU^*=\sum_jP_j$, so
\[
 W=\sum_jP_jE_j+\Bigl(I-\sum_jP_j\Bigr)E^\perp ,
\]
which is the operator $W$ of \cref{hgeo:thm:finite} for the family with its
complementary projection $I-\sum_jP_j$, of residue $E^\perp$, adjoined; in
the complete case $W=C$ and \eqref{hgeo:us:eq:polar} is the formula of
\cref{hgeo:thm:finite}. What is new in \cref{hgeo:us:thm:main} is the infinite
case: then no strong Hahn sum $\sum_jP_j$ exists, and $P_\vee$ is obtained
from the synthesized columns alone. By \cref{hgeo:us:rem:joins}, $P_\vee$ is
still the join of the family.
\end{remark}

\begin{remark}[The countable obstructions are instances of necessity
\mergetag]\label{hgeo:us:rem:instances}
For the family $(P_n)$ of \eqref{hgeo:eq:Pn}, which is incomplete (its
residues $E_n$ project onto $\kk e_n$), the coefficient of $P_n$ at $\eta$ is
the block $\left(\begin{smallmatrix}0&n\\n&0\end{smallmatrix}\right)$, so
$C_\eta^\fin e_n=nf_n$ and (B) fails at $\eta$: this is
\cref{hgeo:cor:nosimultaneous}. For the projections $\widetilde P_n$ of
\cref{hgeo:prop:incoherent}, $C^\fin_{\eta/n}\ne0$ for every $n$, so $S_C$
contains $\eta,\eta/2,\eta/3,\ldots$ and (H) fails: this is
\cref{hgeo:prop:incoherent}. Source~06's proofs argue instead from the
support of a hypothetical unitary; they are kept as the first routes.
\Cref{hgeo:us:thm:uniform} shows that uniform bounds on every individual
coefficient still do not give (B), and \cref{hgeo:us:thm:support} adds
(B) with $b_\gamma\le1$ to a complete version of the family of
\cref{hgeo:prop:incoherent}.
\end{remark}

\section{Support control, gauge freedom, and finite jets}
\label{hgeo:us:sec:consequences}

\begin{corollary}[No new scales \src{07}]\label{hgeo:us:cor:support}
Assume the criterion and put $S=S_C\setminus\{0\}$. The construction
\eqref{hgeo:us:eq:constructiona}--\eqref{hgeo:us:eq:constructionb} satisfies
\[
 \supp(U-I)\cup\supp(P_\vee-E)\subseteq S^*\setminus\{0\},
\]
where $S^*$ is the additive monoid generated by $S$
(\cref{hgeo:lem:support}; 07: $\langle S\rangle$). In particular it needs no
divisible hull and no enlargement of $\Gamma$.
\end{corollary}
\begin{proof}
Write $C=E+C_+$ with $\supp C_+\subseteq S$. The positive supports of
$\Omega-I$, $\Omega^{-1}-I$, $P_\vee-E$, $W-I$ and $W^*W-I$ lie, successively,
in the positive part of $S^*$, and so do those of $\abs W^{-1}-I$ and $U-I$.
\Cref{hgeo:lem:support} justifies each infinite power-series operation; the
ordinary factors $E$, $E^\perp$, $I$ have support $\{0\}$.
\end{proof}

\begin{corollary}[Gauge classification \src{07}]\label{hgeo:us:cor:gauge}
If $U$ is one near-identity straightener, the near-identity straighteners are
exactly the products $U\Xi$, where $\Xi$ ranges over the near-identity
unitaries in $\BH$ commuting with every $E_i$.
\end{corollary}
\begin{proof}
For another straightener $U'$, the unitary $\Xi=U^*U'$ satisfies
$\Xi E_i\Xi^*=E_i$, hence $\Xi E_i=E_i\Xi$, and $\Xi_0=I$. The converse is
substitution.
\end{proof}

The formula of \cref{hgeo:us:thm:main} selects one member of this gauge
class; uniqueness is not claimed without a specified choice. A commuting gauge
also preserves $E$ and $E^\perp$, by boundedness of each coefficient and
strong convergence of the residual partial sums. For a single projection the
gauge freedom is the nonuniqueness recorded after \cref{hgeo:thm:normalform}.

\begin{corollary}[Scale enlargement cannot repair failure \src{07}]
\label{hgeo:us:cor:extension}
Let $\iota:\Gamma\hookrightarrow\Gamma'$ be an order-preserving group
embedding, and regard the family in $\mathcal B(H)((t^{\Gamma'}))$. It has a
simultaneous unitary there if and only if it has one in $\BH$.
\end{corollary}
\begin{proof}
One direction is coefficient transport (\cref{hgeo:prop:extension}).
Conversely, a straightener over $\Gamma'$ gives the two necessary conditions
for the canonical columns over $\Gamma'$. These have no coefficients outside
$\iota(\Gamma)$, and inside it they are exactly the original maps
$C_\gamma^\fin$. Their boundedness is an ordinary Hilbert-space fact, and a
subset of $\Gamma$ is well ordered exactly when its image under $\iota$ is.
The criterion applies over $\Gamma$.
\end{proof}

\Cref{hgeo:prop:extension} transports solutions upward; the new content of
\cref{hgeo:us:cor:extension} is the converse, that nonexistence persists under
every enlargement, cofinal or not.

\subsection{Finite jets for an entire infinite family}
In this subsection $\Gamma=\Z$, so an integral Hahn operator is an ordinary
formal power series in $\mathcal B(H)[[t]]$. For $N\ge0$ put
\[
 \mathscr B^{[N]}=\mathcal B(H)[t]/(t^{N+1}).
\]
A unitary modulo $t^{N+1}$ is an element whose two unitary identities hold in
this quotient, not a polynomial that is exactly unitary at ordinary real
values of $t$.

\begin{theorem}[Jetwise existence implies existence \src{07}]
\label{hgeo:us:thm:jets}
For a set-indexed orthogonal family $(P_i)$ of projections in
$\mathcal B(H)[[t]]$, a simultaneous near-identity formal unitary exists if and
only if for every $N$ there is a near-identity unitary
$U_N\in\mathscr B^{[N]}$ with
\[
 P_i\equiv U_NE_iU_N^*\pmod{t^{N+1}}\qquad\text{for every }i\in\mathcal I .
\]
No compatibility between the separately chosen $U_N$ is assumed, and no bound
uniform in $N$ is required.
\end{theorem}
\begin{proof}
Truncation proves necessity. For sufficiency fix $k\ge0$ and choose $N\ge k$.
The identity $C^{(N)}=U_N\Dg_E(U_N^*)$ modulo $t^{N+1}$ satisfies
$C^{(N)}E_i\equiv U_NE_iU_N^*E_i\equiv P_iE_i$, so its coefficient of $t^k$
is a bounded extension of the canonical column $C_k^\fin$. This is finite
coefficient algebra, valid modulo $t^{N+1}$, and it proves (B) at every $k$.
Condition (H) is automatic for a subset of $\N$. \Cref{hgeo:us:thm:main}
then constructs one actual series, independently of the choices of $U_N$.
\end{proof}

The order of quantifiers matters. The theorem asks for one finite jet
straightening \emph{all} members at once. It does not say that exact
straightening of every finite subfamily implies a common straightener: the
families of \cref{hgeo:us:thm:uniform,hgeo:thm:nojoin} straighten on every
finite subfamily and already fail the whole-family jet condition at their
first scale.

\subsection{The first-order obstruction is gauge invariant}
Assume $E=I$ and $\Gamma=\Z$. The following proposition makes the initial
linear obstruction explicit before the higher coefficients are known to
synthesize.

\begin{proposition}[First-order equation \src{07}]\label{hgeo:us:prop:first}
A bounded skew-adjoint $Y$ with
\begin{equation}\label{hgeo:us:eq:firstorder}
 [Y,E_i]=(P_i)_1\qquad(i\in\mathcal I)
\end{equation}
exists if and only if $C_1^\fin$ has a bounded extension $C_1$. Then $C_1$ is
skew-adjoint, $\Dg_E(C_1)=0$, and the solutions are exactly $Y=C_1+Y_c$ with
$Y_c$ bounded, skew-adjoint and commuting with every $E_i$.
\end{proposition}
\begin{proof}
The coefficient of $t$ in $P_i^2=P_i$ gives $E_i(P_i)_1E_i=0$ and
$(I-E_i)(P_i)_1(I-E_i)=0$. For $i\ne j$ the coefficient of $t$ in $P_iP_j=0$
gives $E_i(P_j)_1+(P_i)_1E_j=0$. Suppose $C_1$ is bounded; its block
$E_kC_1E_l$ is $E_k(P_l)_1E_l$, so the diagonal blocks vanish. For $k\ne l$,
multiplying $E_l(P_k)_1+(P_l)_1E_k=0$ by $E_l$ on the left and by $E_k$ on
the right gives $E_l(P_l)_1E_k=-E_l(P_k)_1E_k=-E_lC_1E_k$; since $(P_l)_1$ is
self-adjoint, the left side is $(E_kC_1E_l)^*$. So $C_1$ is skew-adjoint.
The blocks of $[C_1,E_i]$ are $E_kC_1E_i$ for $l=i\ne k$,
$-E_iC_1E_l$ for $k=i\ne l$, and zero otherwise; those of $(P_i)_1$ are the
same, by the two identities from $P_i^2=P_i$ and skew-adjointness. Equality
of all blocks on the dense algebraic direct sum gives
$[C_1,E_i]=(P_i)_1$.

Conversely, if $Y$ solves \eqref{hgeo:us:eq:firstorder}, then
$(Y-\Dg_E(Y))E_i=YE_i-E_iYE_i=(P_i)_1E_i$ for every $i$, so $Y-\Dg_E(Y)$ is
the required bounded extension $C_1$. The difference of two solutions
commutes with every $E_i$, which describes all solutions.
\end{proof}

In particular an unbounded canonical first coefficient cannot be made bounded
by a clever choice of diagonal gauge. When the full criterion holds, the
complete-family polar formula begins with
\[
 U_1=C_1,\qquad U_2=\frac{C_2-C_2^*}{2}-\frac12C_1^*C_1 ,
\]
since $(C^*C)_1=C_1+C_1^*=0$ and $(C^*C)_2=C_2+C_2^*+C_1^*C_1$, and the first
two coefficients of the inverse square root then give the formula.

\section{An exact atomic functional calculus, and joins}
\label{hgeo:us:sec:calculus}

In this section, except in \cref{hgeo:us:rem:joins}, $H$ is complex, every
$E_i$ is nonzero and $\sSOT_iE_i=I$. Let
\begin{equation}\label{hgeo:us:eq:rho}
 \rho(a)=\sSOT_{i\in\mathcal I}a_iE_i,\qquad a=(a_i)\in\ell^\infty(\mathcal I).
\end{equation}
This is the ordinary normal unital $*$-representation of the bounded scalar
families on the orthogonal Hilbert direct sum, with
$\norm{\rho(a)}=\sup_i\abs{a_i}$. The characteristic function of
$\{i\}$ is written $1_i$.

\begin{definition}[Coherent coefficientwise atomic calculus \src{07}]
\label{hgeo:us:def:calculus}
A \emph{calculus} for $(P_i)$ is a unital complex-linear $*$-homomorphism
$\Phi:\ell^\infty(\mathcal I)\to\Bnn$ with $\Phi(1_i)=P_i$ and $\Phi_0=\rho$,
satisfying the following regularity. There is one well-ordered
$S\subseteq\Gamma_{\ge0}$ containing $\supp\Phi(a)$ for every $a$, and for
each $\gamma$ the coefficient map
\[
 \Phi_\gamma:\ell^\infty(\mathcal I)\to\mathcal B(H),\qquad
 a\mapsto(\Phi(a))_\gamma ,
\]
is normal and completely bounded. Normal means ultraweakly continuous;
completely bounded means that its entrywise actions on matrices of all
finite sizes have one common finite norm bound $\norm{\Phi_\gamma}_{\cb}$.
These bounds may depend on $\gamma$.
\end{definition}

\begin{theorem}[Atomic calculus equivalence \src{07}]
\label{hgeo:us:thm:calculus}
For a complete residual family $(P_i)$, a calculus in the sense of
\cref{hgeo:us:def:calculus} exists if and only if \textup{(B)} and
\textup{(H)} hold. It is then unique and given by
\begin{equation}\label{hgeo:us:eq:Phi}
 \Phi(a)=U\rho(a)U^*
\end{equation}
for any simultaneous straightener $U$ of \cref{hgeo:us:thm:main}, and
\begin{equation}\label{hgeo:us:eq:cbbound}
 \norm{\Phi_\gamma}_{\cb}\le
 \sum_{\alpha+\beta=\gamma}\norm{U_\alpha}\,\norm{U_\beta}.
\end{equation}
The displayed sum is finite, and the common support may be chosen inside
$\supp U+\supp U$.
\end{theorem}
\begin{proof}
Given the criterion, choose $U$ and define \eqref{hgeo:us:eq:Phi}. This is a
unital $*$-homomorphism with the required atoms and residue. Each coefficient
is a finite sum of maps $a\mapsto U_\alpha\rho(a)U_\beta^*$. Such a map is
normal, and its matrix amplification is multiplication on either side by the
diagonal amplifications of $U_\alpha$ and $U_\beta^*$, so its completely
bounded norm is at most $\norm{U_\alpha}\norm{U_\beta}$. This proves
\eqref{hgeo:us:eq:cbbound} and the support claims.

Conversely, let a calculus be given and fix $\gamma$. For a finite
$\mathcal J=\{i_1,\ldots,i_n\}$, rectangular matrix multiplication gives
\[
 \sum_{i\in\mathcal J}(P_i)_\gamma E_i
 =\bigl[\Phi_\gamma(1_{i_1})\ \cdots\ \Phi_\gamma(1_{i_n})\bigr]
 \begin{bmatrix}E_{i_1}\\ \vdots\\ E_{i_n}\end{bmatrix}.
\]
The column of the $E_i$ has norm at most one, because its squared norm
operator is $\sum_{i\in\mathcal J}E_i\le I$. The row
$[1_{i_1}\ \cdots\ 1_{i_n}]$ has norm one in
$M_{1,n}(\ell^\infty(\mathcal I))$, since at each index at most one entry is
one. Complete boundedness therefore gives
\begin{equation}\label{hgeo:us:eq:synthesiscb}
 \Bigl\|\sum_{i\in\mathcal J}(P_i)_\gamma E_i\Bigr\|
 \le\norm{\Phi_\gamma}_{\cb},
\end{equation}
which is (B) by \eqref{hgeo:us:eq:finitesup}. The common support of the
calculus contains the active column support, which gives (H), and
\cref{hgeo:us:thm:main} supplies $U$.

For uniqueness there is a simpler algebraic argument than approximation by
finitely supported families. If $\Psi$ is \emph{any} unital complex-linear
homomorphism into $\BH$ with the prescribed atoms, then
$\Psi(a)P_i=\Psi(a1_i)=a_iP_i$, and conjugating by $U$ gives
$(U^*\Psi(a)U)E_i=a_iE_i$ for every $i$. Every coefficient of
$U^*\Psi(a)U-\rho(a)$ is a bounded operator vanishing on every $E_iH$, whose
algebraic sum is dense, so it is zero. Thus $\Psi(a)=U\rho(a)U^*$ for all
$a$; in particular the regular calculus is unique, and the gauge freedom of
\cref{hgeo:us:cor:gauge} does not change it.
\end{proof}

The last paragraph is a rigidity statement once a straightener exists. It
does not prove that an arbitrary algebra map defined pointwise has coherent
support or supplies a straightener without the hypotheses of the theorem;
that automatic-regularity question remains open
(\cref{hgeo:us:q:regularity}).

\subsection{Boolean projections and coefficientwise additivity}
For $\mathcal S\subseteq\mathcal I$ put
$E(\mathcal S)=\rho(1_{\mathcal S})$ and
$P(\mathcal S)=\Phi(1_{\mathcal S})=UE(\mathcal S)U^*$. Then
\[
 P(\mathcal S\cap\mathcal S')=P(\mathcal S)P(\mathcal S'),\qquad
 P(\mathcal I\setminus\mathcal S)=I-P(\mathcal S),\qquad P(\mathcal I)=I .
\]
If $\mathcal S_1,\mathcal S_2,\ldots$ are disjoint, then for every $\gamma$
\begin{equation}\label{hgeo:us:eq:additivity}
 \Bigl(P\bigl(\textstyle\bigcup_n\mathcal S_n\bigr)\Bigr)_\gamma
 =\sSOT_{n\ge1}\bigl(P(\mathcal S_n)\bigr)_\gamma ,
\end{equation}
because $\sSOT_nE(\mathcal S_n)=E(\bigcup_n\mathcal S_n)$ and each coefficient
of $UE(\mathcal S)U^*$ is a fixed finite sum of left and right multiplications
by bounded operators. The same argument works for the net of finite subsets of
any disjoint set-indexed family.

\begin{proposition}[Joins when the criterion holds \src{07}]
\label{hgeo:us:prop:joins}
The projection $P(\mathcal S)$ is the least upper bound of
$\{P_i:i\in\mathcal S\}$ in the full projection poset of $\BH$.
\end{proposition}
\begin{proof}
It is an upper bound. If $P'$ is another upper bound, put $P''=U^*P'U$. Then
$P''E_i=E_i$ for $i\in\mathcal S$. Each bounded coefficient of
$(P''-I)E(\mathcal S)$ vanishes on the dense algebraic sum of the $E_iH$,
$i\in\mathcal S$, and on the range of $I-E(\mathcal S)$, so
$P''E(\mathcal S)=E(\mathcal S)$; taking adjoints,
$E(\mathcal S)P''=E(\mathcal S)$. Thus $E(\mathcal S)\le P''$, that is
$P(\mathcal S)\le P'$.
\end{proof}

\begin{remark}[Joins for incomplete families and real scalars
\mergetag]\label{hgeo:us:rem:joins}
The proof of \cref{hgeo:us:prop:joins} uses only a straightener. Hence, for a
possibly incomplete family over a real or complex $H$ satisfying (B) and (H),
with straightener $U$, every subfamily $\{P_i:i\in\mathcal S\}$ has the join
$UE(\mathcal S)U^*$, where $E(\mathcal S)=\sSOT_{i\in\mathcal S}E_i$: it is an
upper bound since $UE(\mathcal S)U^*P_i=UE(\mathcal S)E_iU^*=P_i$ for
$i\in\mathcal S$, and the leastness argument above applies verbatim. By
\eqref{hgeo:us:eq:range}, the range projection $P_\vee=UEU^*$ of
\cref{hgeo:us:thm:main} is the join of the whole family. Every subfamily of
a family satisfying the criterion satisfies it too (its columns are
restrictions, its active support a subset), so this is a sufficient
condition for all subfamily joins at once. It is not necessary for the
join of a family to exist: the whole family of \cref{hgeo:us:thm:uniform} has
the join $I$ but no straightener.
This is the partial answer to \cref{hgeo:q:notdeveloped} recorded in
\cref{hgeo:us:sub:status}; \cref{hgeo:us:q:boolean} asks for the converse
direction.
\end{remark}

This is an atomic spectral calculus in a carefully stated sense \src{07}. It
does not make $(P_i)$ strongly Hahn summable, since infinitely many nonzero
$P_i$ have a nonzero coefficient at exponent zero. Nor does
\eqref{hgeo:us:eq:additivity} give valuation convergence: for the constant
coordinate projections on $\ell^2(\N)$ and a vector $x$ with infinitely many
nonzero coordinates, $E(\{1,\ldots,n\})x\to x$ in the ordinary Hilbert norm,
whereas each nonzero difference, as a constant Hahn vector, has valuation
zero. No continuous-spectrum spectral measure and no valuation-topologically
countably additive general spectral theorem is proved. This is consistent with
the non-claims of sources~03 and~06 on spectral measures
(\cref{hgeo:q:notdeveloped}).

\subsection{Coherent eigenvalue reconstruction}
Eigenvalue series need not be uniformly bounded as whole Hahn scalars,
provided they are bounded at each coefficient and have one common support.

\begin{proposition}[Weighted reconstruction \src{07}]
\label{hgeo:us:prop:weighted}
Suppose the criterion holds with $E=I$. Let $(\lambda_i)_{i\in\mathcal I}$ be
scalar Hahn series whose supports lie in one well-ordered
$S_\lambda\subseteq\Gamma$, and assume
$\sup_i\abs{(\lambda_i)_\gamma}<\infty$ for every $\gamma$. Then
\[
 D_\lambda=\sum_\gamma\rho\bigl(((\lambda_i)_\gamma)_{i\in\mathcal I}\bigr)
 t^\gamma,\qquad N_\lambda=UD_\lambda U^*
\]
are Hahn operators, $N_\lambda$ is normal, and
$N_\lambda P_i=P_iN_\lambda=\lambda_iP_i$. It is the unique operator
$A\in\BH$ with $AP_i=\lambda_iP_i$ for all $i$.
\end{proposition}
\begin{proof}
The coefficient bound makes every diagonal coefficient bounded, and the
support hypothesis makes their series a Hahn series. All diagonal
coefficients commute with one another and with their adjoints, so $D_\lambda$
is normal. Conjugation gives the equations and normality of $N_\lambda$. For
uniqueness, conjugate the difference of two solutions by $U$: every
coefficient annihilates every $E_iH$, hence is zero by density.
\end{proof}

This reconstructs an operator from coherent atomic data. It does not assert
that every normal Hahn operator admits such data, nor that its
algebra-relative spectrum consists only of the labels $\lambda_i$;
accumulation phenomena and inverse coherence are separate issues, as the
spectral report emphasizes \cite{RepoSpectral}.

\section{Uniform individual bounds do not suffice}\label{hgeo:us:sec:uniform}

The next construction strengthens the observation that an unbounded direct
sum of infinitesimal rotations need not be a Hahn operator
(\cref{hgeo:cor:nosimultaneous}, where the coefficients of the $P_n$ at $\eta$
have norm $n$). Here every individual projection is controlled uniformly, at
every fixed order, over the entire countable family \src{07}.

\begin{theorem}[A uniformly controlled nonsynthesizable family \src{07}]
\label{hgeo:us:thm:uniform}
There is a countable pairwise orthogonal family of rank-one projections
$(P_i)$ in $\mathcal B(\ell^2)[[t]]$, with complete coordinate residues
$(E_i)$, such that:
\begin{enumerate}[label=\textup{(\roman*)}]
\item all supports lie in $\N$;
\item for every $n\ge0$, $\sup_i\norm{(P_i)_n}<\infty$;
\item every finite subfamily has an exact near-identity simultaneous Hahn
unitary;
\item the whole family has no simultaneous Hahn unitary, even after any
ordered enlargement of the exponent group;
\item a specified subfamily has no least upper bound in the projection poset,
although the whole family has the least upper bound $I$.
\end{enumerate}
The same example works over real coefficients and after complexification.
\end{theorem}

\subsection{The block construction}
For $m\ge1$ put $d_m=4^{m^2}$ and
\[
 H_m=\kk^{d_m}\oplus\kk^{d_m},\qquad H=\bigoplus_{m\ge1}H_m ,
\]
a separable infinite-dimensional Hilbert space (identified with $\ell^2$ by
its coordinate basis). Let $u_m$ be the normalized all-ones vector of the left
half of $H_m$ and $u_m'$ the corresponding vector of the right half. With the
rank-one notation $xy^*:z\mapsto\ip{z}{y}x$, define the ordinary finite-rank
skew-adjoint operator
\begin{equation}\label{hgeo:us:eq:Xm}
 X_m=m\bigl(u_m'u_m^*-u_mu_m'^*\bigr)\quad\text{on }H_m .
\end{equation}
Then $X_mu_m=mu_m'$, $X_mu_m'=-mu_m$, $\norm{X_m}=m$, and $X_m$ annihilates
the orthogonal complement of $\Span\{u_m,u_m'\}$ in $H_m$. For each
coordinate vector $e\in H_m$ let $E_e=ee^*$ and define
\begin{equation}\label{hgeo:us:eq:Pe}
 V_m(t)=\exp(tX_m),\qquad P_e(t)=V_m(t)E_eV_m(t)^* .
\end{equation}
These are formal matrix power series on a \emph{finite} block, extended by
zero to the other blocks. Since $X_m^*=-X_m$, $V_m$ is formally unitary. Within
a block the $P_e$ are mutually orthogonal and sum to the ordinary projection
$\Pi_m$ onto $H_m$; different blocks are orthogonal; their residues are all the
coordinate projections of $H$. No global operator $\exp(t\bigoplus_mX_m)$ is
assumed in \eqref{hgeo:us:eq:Pe}: the direct sum of the $X_m$ is unbounded.
Only the separately legitimate finite-block exponentials define the
individual projections.

\subsection{Uniform estimates at every order}
For every coordinate vector $e\in H_m$ and $k\ge1$,
\begin{equation}\label{hgeo:us:eq:column}
 \norm{X_m^ke}=\frac{m^k}{\sqrt{d_m}} ,
\end{equation}
because the component of $e$ in the plane $\Span\{u_m,u_m'\}$ has length
$1/\sqrt{d_m}$ and $X_m$ acts on that plane as $m$ times an isometry. For
$n\ge1$, expanding \eqref{hgeo:us:eq:Pe} gives
\[
 (P_e)_n=\sum_{a+b=n}\frac{X_m^aE_e(-X_m)^b}{a!\,b!} .
\]
The two terms with $a=0$ or $b=0$ have norm $m^n/(n!\sqrt{d_m})$ each. Each
interior term is rank one with norm $m^n/(a!\,b!\,d_m)$ by
\eqref{hgeo:us:eq:column}, and $\sum_{a=1}^{n-1}n!/(a!\,b!)=2^n-2$. Hence
\begin{align}
 \norm{(P_e)_n}
 &\le\frac{m^n}{n!}\Bigl(\frac{2}{\sqrt{d_m}}+\frac{2^n-2}{d_m}\Bigr)
 \notag\\
 &\le\frac{2^nm^n}{n!\,2^{m^2}} .\label{hgeo:us:eq:allorder}
\end{align}
For fixed $n$ the right-hand side is bounded over all $m\ge1$; at order zero
the bound is one. This proves (ii) without a bound depending on the
individual projection. In fact the positive-order coefficients from far-out
blocks tend to zero, uniformly within each block. Every finite subfamily meets
only finitely many blocks, and the direct sum of the corresponding finitely
many $V_m$ with the identity on the other blocks is an actual
bounded-coefficient Hahn unitary straightening that subfamily; this proves
(iii) directly, and (i) is clear.

\subsection{The synthesis obstruction}
Since the diagonal of $X_m$ in the coordinate basis is zero,
\[
 (P_e)_1e=[X_m,E_e]e=X_me .
\]
Hence $C_1^\fin$ restricts to $X_m$ on the whole finite block $H_m$, and
$\norm{C_1^\fin u_m}=m$ with $\norm{u_m}=1$. The canonical first coefficient is
unbounded: (B) fails, although (H) and every individual estimate
\eqref{hgeo:us:eq:allorder} hold. \Cref{hgeo:us:thm:main} excludes a common
unitary, \cref{hgeo:us:cor:extension} excludes one at every enlarged scale,
and \cref{hgeo:us:prop:first} shows that no block-diagonal first-order gauge
can hide the failure. This proves (iv).

The mechanism is collective. Each coordinate vector has a very small
component in the active plane, but a coherent sum of many coordinates recovers
the unit vector $u_m$, on which the operator norm is $m$. An estimate on
individual columns does not bound synthesis on arbitrary unit vectors of their
Hilbert direct sum.

\subsection{A missing Boolean join}
Let $\mathcal L$ consist of the coordinate vectors in the left halves of all
blocks, and suppose a projection $Q\in\BH$ were the join of
$\{P_e:e\in\mathcal L\}$. For every coordinate $f\notin\mathcal L$, $I-P_f$ is
an upper bound for this family, since all its members are orthogonal to
$P_f$. Leastness would give
\[
 QP_e=P_e\ (e\in\mathcal L),\qquad QP_f=0\ (f\notin\mathcal L).
\]
Each block projection $\Pi_m$ is the sum of its finitely many $P_e$, so
$Q\Pi_m=\sum_{e\in\mathcal L\cap H_m}P_e$ is self-adjoint; therefore $Q$
commutes with $\Pi_m$, and its compression to $H_m$ is forced to be
\[
 Q|_{H_m}=V_mE_{L,m}V_m^* ,
\]
where $E_{L,m}$ is the ordinary left-half projection. Its first coefficient is
$[X_m,E_{L,m}]$, which sends $u_m$ to $mu_m'$ and $u_m'$ to $mu_m$ and
annihilates the orthogonal complement of the plane; its norm is exactly $m$.
But $Q_1$ would be a bounded operator on $H$, whose compressions to all $H_m$
have norm at most $\norm{Q_1}$. This is impossible. So the selected subfamily
has no join; it is not merely that a particular formula for the join
diverges.

The whole family does have the join $I$. Any upper-bound projection has a
residue dominating every coordinate projection, hence residue $I$; its
complementary projection has zero residue and is zero by
\cref{hgeo:us:lem:integral}. This proves (v) and
\cref{hgeo:us:thm:uniform}.\hfill$\square$

\begin{remark}[Total vectors are not a synthesis theorem \src{07}]
\label{hgeo:us:rem:total}
The vectors $V_me$, taken one at a time, form an orthonormal family in
$\HH{H}$ with $\Gamma=\Z$. If a nonzero Hahn vector were orthogonal to every one,
its leading ordinary coefficient would be orthogonal to every coordinate
vector, which is impossible. So the family is total in this orthogonality
sense. Nevertheless no bounded-coefficient Hahn unitary synthesizes it from
the ordinary basis: totality and a valid global synthesis operator are
different notions in this category.
\end{remark}

Compared with \cref{hgeo:thm:nojoin,hgeo:cor:nosimultaneous}, whose family is
incomplete and has coefficients of norm $n$ at $\eta$, the family here is
complete, has uniformly bounded coefficients at every order and common support
$\N$, and still has no straightener, and a subfamily still has no join
\src{07}.

\section{The independent support obstruction}\label{hgeo:us:sec:support}

The family of \cref{hgeo:us:thm:uniform} is perfectly support coherent. The
next one has all canonical coefficient maps bounded, with norm at most one,
but active support that is not well ordered.

\begin{theorem}[Bounded columns with incoherent support \src{07}]
\label{hgeo:us:thm:support}
Over $\Gamma=\Q$ there is a countable orthogonal family of rank-one Hahn
projections with complete coordinate residues such that \textup{(B)} holds
with $b_\gamma\le1$ for every $\gamma$, whereas \textup{(H)} fails. Every
finite subfamily straightens, but the whole family does not, even after
ordered scale enlargement.
\end{theorem}
\begin{proof}
Take $H=\bigoplus_{n\ge1}\kk^2$ with basis $e_n,f_n$ of the $n$th block, and
on each block put
\[
 \mathsf J=\begin{pmatrix}0&-1\\1&0\end{pmatrix},\qquad
 V_n=\bigl(I+t^{1/n}\mathsf J\bigr)\bigl(1+t^{2/n}\bigr)^{-1/2},
\]
with $I$ the $2\times2$ identity. The scalar inverse square root is a Hahn
series supported in $(2/n)\N$. Since $\mathsf J^*=-\mathsf J$ and
$\mathsf J^2=-I$, $V_n$ is unitary. Define the block projections
$P_{n,j}=V_nE_{n,j}V_n^*$, $j=1,2$, where $E_{n,1},E_{n,2}$ project onto
$\kk e_n,\kk f_n$, and extend them by zero off the block. Each has a
well-ordered support contained in $(1/n)\N$. The blockwise canonical column
expression is
\begin{equation}\label{hgeo:us:eq:rationalblock}
 C^{(n)}=P_{n,1}E_{n,1}+P_{n,2}E_{n,2}=V_n\,\Dg_E\bigl(V_n^*\bigr)
 =\frac{I+t^{1/n}\mathsf J}{1+t^{2/n}},
\end{equation}
where on a block $\Dg_E$ takes the diagonal of a $2\times2$ matrix. Its coefficient at
$k/n$ is $(-1)^{k/2}I$ for even $k$ and $(-1)^{(k-1)/2}\mathsf J$ for odd $k$.
At a fixed rational $\gamma>0$, the canonical coefficient on all of $H$ is the
block diagonal operator with this coefficient in the blocks where $n\gamma$ is
an integer and zero elsewhere. Every nonzero block has norm one, so the full
coefficient has norm at most one; this proves (B), and order zero is the
identity. The active support is all of $\Q_{\ge0}$, since every positive
rational $k/n$ carries a nonzero block coefficient; it contains the strictly
decreasing sequence $1,1/2,1/3,\ldots$ and is not well ordered. So (H) fails
and \cref{hgeo:us:thm:main} excludes a simultaneous unitary. Finite
subfamilies and scale enlargement are handled by
\cref{hgeo:us:rem:finite,hgeo:us:cor:extension}.
\end{proof}

The projections $P_{n,1}$ project onto $\K(e_n+t^{1/n}f_n)$: they are the
projections $\widetilde P_n$ of \cref{hgeo:prop:incoherent} with $\eta=1$.
The theorem adds their complements $P_{n,2}$, so that the residual family is
complete, and proves (B) with $b_\gamma\le1$. Source~06's proof of
\cref{hgeo:prop:incoherent} is a second route to nonexistence: a straightener
$U$ would give $\supp P_{n,1}\subseteq\supp U+\supp U$, a well-ordered set,
which cannot contain $1,1/2,1/3,\ldots$.

\begin{remark}[Any group containing $\Q\eta$ \mergetag]
\label{hgeo:us:rem:transport}
If $\Gamma'$ contains $\Q\eta$ for some $\eta>0$, the order embedding
$\Q\to\Gamma'$, $r\mapsto r\eta$, transports the family to the one built with
$t^{\eta/n}$ in place of $t^{1/n}$, and \cref{hgeo:us:cor:extension} shows that
it has no straightener over $\Gamma'$. This is the hypothesis of
\cref{hgeo:prop:incoherent}.
\end{remark}

Both counterexamples are rank one, countable, and already defined over real
coefficients \src{07}. The first cannot be repaired by bounding individual
coefficients at more Taylor orders: it fails collective boundedness at its
first order. The second cannot be repaired by improving any coefficient bound,
since every such bound already has the value one. Their obstructions belong to
different parts of the definition of a Hahn operator.

\section{Surreal workspaces and the omnific boundary}\label{hgeo:us:sec:surreal}

\subsection{From formal scales to actual surreal monomials}
Let $\Gamma$ be a set-sized ordered additive subgroup of $\No$. The
substitution $t^\gamma=\omega^{-\gamma}$ of \eqref{hgeo:eq:embedding}
identifies $\F$ with a normal-form subfield of $\No$ and $\K$ with a subfield
of $\No[i]$; well-ordered $t$-support becomes reverse well-ordered Conway
support \cite{Conway,LInnocenteMantova}. All operator identities of this part
therefore hold for real or surcomplex scalar series in such a workspace. They
are not asymptotic formulas modulo an unspecified error but equalities of full
Hahn normal forms. In \cref{hgeo:us:thm:uniform} one may take $t=\omega^{-1}$;
\cref{hgeo:us:thm:support} uses the actual monomials $\omega^{-1/n}$
\src{07}.

\begin{corollary}[Set-indexed surreal descent \src{07}]\label{hgeo:us:cor:surreal}
Consider a set-indexed family of orthogonal projections given as set-supported
series of ordinary bounded operators with exponents in $\No$. The criterion of
\cref{hgeo:us:thm:main} applies after all their supports are placed in one
set-sized ordered subgroup. If no straightener exists in that subgroup, none
exists in any larger bounded-coefficient surreal Hahn workspace.
\end{corollary}
\begin{proof}
The union of the supports of a set-indexed family is a set, and the subgroup
of $(\No,+)$ it generates by finite sums and integer multiples is a set-sized
ordered subgroup; the family belongs to the corresponding Hahn operator
algebra. Any proposed straightener given by a set-supported bounded-operator
series, even one using further surreal exponents, lies with the family in a
larger such subgroup. Apply \cref{hgeo:us:cor:extension}.
\end{proof}

Source~07 states this for its specified category of operators, not as an
identification of all class-linear maps on a proper-class vector space with
Hahn series; it uses only set-indexed support unions, and neither a
large-cardinal axiom nor global choice beyond the ordinary set theory used for
Hilbert spaces. It leaves automatic representation theorems for class-linear
adjointable maps outside its scope. Within this report such a theorem is
\cref{hgeo:thm:class}(ii), and it removes the restriction for maps with an
everywhere-defined class adjoint.

\begin{corollary}[Class unitaries \mergetag]\label{hgeo:us:cor:classmaps}
Work in the class formulation of \cref{hgeo:sec:global} (NBG without global
choice). Let $(P_i)_{i\in\mathcal I}$ be a set-indexed family of pairwise
orthogonal class projections on $\mathscr H_{\No}(H)$, that is,
everywhere-defined self-adjoint idempotent class-linear maps. Then the
$P_i$ have set-supported expansions with supports in one set-sized
ordered subgroup $\Gamma$. If the family fails \textup{(B)} or \textup{(H)}
there, no class-linear unitary of $\mathscr H_{\No}(H)$ conjugates every $E_i$
to $P_i$. If it satisfies both, the straightener of \cref{hgeo:us:thm:main}
acts on the whole class and conjugates every $E_i$ to $P_i$ there.
\end{corollary}
\begin{proof}
Each $P_i$ is its own everywhere-defined class adjoint, so
\cref{hgeo:thm:class}(ii) gives it a set support $S_i$; the union over the set
$\mathcal I$ is a set by replacement, and it generates a set-sized ordered
subgroup $\Gamma$. A class unitary $U$ has the everywhere-defined class
adjoint $U^{-1}$, so it too has a set support $S_U$. In the set-sized subgroup
$\Gamma'$ generated by $\Gamma$ and $S_U$ the identities $P_i=UE_iU^*$ hold on
the workspace vectors, hence as identities of Hahn series, because operators
of the algebra are determined by their values on constants. By
\cref{hgeo:us:cor:extension} the family then straightens over $\Gamma$. The
last assertion is \cref{hgeo:thm:class}(ii): the straightener has set support
and acts on every class vector, and the identities hold there by coefficient
transport (\cref{hgeo:prop:extension}).
\end{proof}

A class-linear map without an everywhere-defined class adjoint, such as a
nonsurjective isometry, is not covered; \cref{hgeo:us:q:classmaps} remains
open for such maps.

\subsection{Omnific entries remove all infinitesimal unitary freedom}
Recall the normal-form description of the omnific integers
\[
 \Oz=\Bigl\{\sum_xr_x\omega^x:\supp(r)\subseteq\No_{\ge0},\ r_0\in\Z\Bigr\},
\]
with reverse well-ordered support and real coefficients
\cite{Conway,LInnocenteMantova}. In the convention $t^\gamma=\omega^{-\gamma}$,
an omnific number has support in nonpositive exponents and an ordinary
integer as constant coefficient. Put $\Oz[i]=\Oz+i\Oz$, whose constant
coefficients are Gaussian integers. Only this description is used.

Fix an ordinary orthonormal basis $(e_j)_{j\in\mathcal N}$ of $H$. The scalar
matrix entry of $A\in\BH$ at $(j,k)$ is the Hahn series
\[
 A_{jk}=\sum_\gamma\ip{A_\gamma e_k}{e_j}\,t^\gamma ;
\]
the matrix need not have finite rows or columns.

\begin{theorem}[Omnific matrix rigidity \src{07}]\label{hgeo:us:thm:omnific}
In a surreal Hahn workspace:
\begin{enumerate}[label=\textup{(\roman*)}]
\item a projection all of whose entries lie in $\Oz$ in the real case, or in
$\Oz[i]$ in the complex case, is an ordinary coordinate projection relative
to $(e_j)$;
\item a unitary with such entries is an ordinary signed permutation in the
real case, or an ordinary permutation with phases in $\{1,-1,i,-i\}$ in the
complex case;
\item a near-identity unitary with such entries is the identity.
\end{enumerate}
\end{theorem}
\begin{proof}
By \cref{hgeo:us:lem:integral} the operator has no negative $t$-exponents,
and by the omnific condition none of its entries has a positive one. So for
every $\gamma>0$ all matrix entries of $A_\gamma$ vanish, and a bounded
operator with all entries zero relative to an orthonormal basis is zero, by
density of finite linear combinations. Hence $A=A_0$ is ordinary, with
integer or Gaussian-integer entries.

For an ordinary orthogonal projection $P$, each diagonal entry is real and
$P_{jj}=\norm{Pe_j}^2\in[0,1]$; being an integer, it is $0$ or $1$. If it is
$0$ then $Pe_j=0$; if $1$, the orthogonal decomposition
$e_j=Pe_j+(I-P)e_j$ gives $Pe_j=e_j$. So $P$ is the coordinate projection onto
the indices with $P_{jj}=1$.

For an ordinary unitary $U$, Parseval's identity in a column gives
$\sum_j\abs{U_{jk}}^2=1$ with nonnegative integer summands, so each column has
exactly one nonzero entry, equal to $\pm1$ or to a Gaussian unit. Orthogonality
of the columns makes the resulting map of basis indices injective, and
surjectivity of $U$ makes it bijective. This proves (ii); a constant unitary
with residue $I$ is $I$, which proves (iii).
\end{proof}

The possible omnific-entry unitary groups are therefore
$\{\pm1\}^{\mathcal N}\rtimes\operatorname{Sym}(\mathcal N)$ and
$\{\pm1,\pm i\}^{\mathcal N}\rtimes\operatorname{Sym}(\mathcal N)$. This is an
elementary boundary corollary, not a claimed new theorem of omnific
factorization. It explains why the infinitesimal rotations of the
constructions above cannot stay inside the omnific matrix ring: positive
$t$-exponents and the omnific support restriction are incompatible.

\section{Status of the questions; further questions of source 07}
\label{hgeo:us:sec:questions}

\subsection{Answered and re-scoped questions}\label{hgeo:us:sub:status}
\begin{itemize}[leftmargin=*]
\item \emph{\Cref{hgeo:q:synthesis}} (Infinite orthogonal synthesis) is
answered. The two conditions that sources~03 and~06 isolated as necessary are,
in the form (B) and (H), also sufficient (\cref{hgeo:us:thm:main}), for every
set-indexed family, complete or not, over every set-sized ordered group. They
are stable under scale enlargement in the strongest sense: solvability is
unchanged by any ordered embedding (\cref{hgeo:us:cor:extension}). Families
satisfying them have joins (\cref{hgeo:us:prop:joins,hgeo:us:rem:joins}).
The criterion is exact but not an effective test
(\cref{hgeo:us:q:intrinsic,hgeo:us:q:effective}), and neither condition can be
dropped (\cref{hgeo:us:thm:uniform,hgeo:us:thm:support}). Stability under
composition, also asked for in the question, is not treated separately.
\item \emph{\Cref{hgeo:q:notdeveloped}} (Which countable families have joins)
is answered in part: families satisfying the criterion, and all their
subfamilies, have joins (\cref{hgeo:us:rem:joins}), and
\cref{hgeo:us:thm:uniform} adds a complete family with uniformly bounded
coefficients some of whose subfamilies lack a join. A characterization of the
countable families with joins is not given; \cref{hgeo:us:q:boolean} asks
whether Boolean joins for all subfamilies force the criterion. Tensor
products, frame bounds and general spectral measures remain undeveloped;
\cref{hgeo:us:q:intrinsic} proposes frame-type bounds as a route to (B).
\item \emph{\Cref{hgeo:q:orthoclosed}} (Orthoclosed subspaces with closed
residue) stays open. Source~07 restates it (\cref{hgeo:us:q:closedresidue})
and notes that its criterion, which starts from existing projections, does not
settle it.
\item \emph{The spectral report's assembly question}
(\replabel{ihs:fr:sec:questions}, first question: which uniform spectral-gap
and coefficient-ideal hypotheses make all cluster rotations assemble into a
bounded-coefficient Hahn unitary?). One bounded-coefficient unitary
conjugating every residual cluster projection to its deformed cluster
projection is a simultaneous straightener of the family of cluster
projections, so for that family \cref{hgeo:us:thm:main} answers the question
in exact-criterion form: such a unitary exists if and only if the family
satisfies (B) at every order and (H). \Cref{hgeo:us:prop:first} is the
corresponding first-order equation, stated for the projection family rather
than for the operator. The spectral report does not assert that its
first-order condition is sufficient at higher orders; the criterion supplies
the missing conditions, namely (B) at every order together with (H). The
question as posed, in terms of spectral gaps and coefficient ideals of the
operator, is not answered in that form. The spectral report is not edited
here, and its other two further questions are not addressed.
\end{itemize}

\subsection{Further questions}\label{hgeo:us:sub:questions}
The questions of source~07 are kept separate from the established statements.
Some continue a question of this report; others ask for a stronger category or
a more intrinsic test. None is a named published problem.

\begin{question}[Closed residue versus orthogonal splitting
\src{07}]\label{hgeo:us:q:closedresidue}
Does every orthoclosed subspace $M=M^{\perp\perp}$ of $\HH{H}$ with closed
ordinary residue admit an orthogonal projection? This is
\cref{hgeo:q:orthoclosed}. The present criterion starts from existing
projections and does not settle it. A promising test case is the graph of a
closed densely defined ordinary operator after a positive-order perturbation
with closed residue.
\end{question}

\begin{question}[Boolean existence versus coherent synthesis
\src{07}]\label{hgeo:us:q:boolean}
Suppose that for every $\mathcal S\subseteq\mathcal I$ there is a projection
separating the atoms in $\mathcal S$ from those outside $\mathcal S$, for a
complete residual family. Must (B) and (H) follow? Such a separating
projection, when it exists, is unique: the difference of two candidates
annihilates all deformed atoms, and its leading coefficient then annihilates
all residual blocks. The difficulty is to derive a common support and
collective coefficient bounds from existence for all subsets, without assuming
them.
\end{question}

\begin{question}[Automatic regularity of an atomic representation
\src{07}]\label{hgeo:us:q:regularity}
Can the common-support and completely bounded coefficient assumptions of
\cref{hgeo:us:def:calculus} be deduced for a unital $*$-homomorphism
$\ell^\infty(\mathcal I)\to\BH$ with the specified atoms and complete residue?
\Cref{hgeo:us:thm:calculus} gives uniqueness once synthesis exists, not this
automatic existence. A bounded, rather than completely bounded, hypothesis on
each coefficient map is an intermediate problem.
\end{question}

\begin{question}[Intrinsic estimates for the synthesis constants
\src{07}]\label{hgeo:us:q:intrinsic}
Find usable conditions on the Gram blocks $E_i(P_i)_\gamma^*(P_j)_\gamma E_j$
that imply (B) with effective bounds, without first constructing $C_\gamma$.
Schur-type block estimates or weighted frame bounds may give tractable
sufficient tests. \Cref{hgeo:us:thm:uniform} shows why bounds on the
individual blocks alone cannot work without control of their correlations.
\end{question}

\begin{question}[Non-discrete jet compactness \src{07}]\label{hgeo:us:q:jets}
Which filtrations of a general ordered group admit an analogue of
\cref{hgeo:us:thm:jets}? A useful answer must specify the truncated rings and
must recover well-ordered support, not merely boundedness at each exponent;
\cref{hgeo:us:thm:support} shows that coefficientwise local solvability cannot
by itself impose that support property.
\end{question}

\begin{question}[Continuous spectral decompositions \src{07}]
\label{hgeo:us:q:continuous}
Replace $\ell^\infty(\mathcal I)$ by $L^\infty(X,\mu)$ and the atomic Hilbert
direct sum by an ordinary direct integral. Is there an exact Hahn criterion
for straightening a deformed projection-valued measure? There are no atomic
columns on which to perform the present dense test; a replacement may involve
bounded coefficient bimodules and a normal conditional expectation.
\end{question}

\begin{question}[Other coefficient algebras \src{07}]\label{hgeo:us:q:algebras}
For which involutive subalgebras $\mathcal M\subseteq\mathcal B(H)$ does the
criterion produce $U\in\mathcal M((t^\Gamma))$? The proof needs closure under
the relevant diagonal compression and the synthesized coefficients; formal
inversion near the identity then stays in any unital algebra. Von Neumann
subalgebras, block operator algebras and ideals with a unit adjoined give
distinct test cases.
\end{question}

\begin{question}[Non-self-adjoint idempotents \src{07}]
\label{hgeo:us:q:idempotents}
What replaces the theorem for pairwise annihilating, not necessarily
self-adjoint idempotents? Similarity may replace unitary equivalence, but
integrality and zero-residue rigidity cannot simply be imported from
\cref{hgeo:us:lem:integral}: positivity of $A_a^*A_a$ is a genuine ingredient
of the argument.
\end{question}

\begin{question}[Effective infinite families \src{07}]\label{hgeo:us:q:effective}
For computably presented countable families, identify useful decidable or
semidecidable subclasses of (B) and (H). A finite algorithm cannot simply
sample more atoms: \cref{hgeo:us:thm:uniform} hides an unbounded synthesis
constant in increasingly large finite blocks. Certificates based on an a
priori structural bound may be more realistic than unrestricted decidability.
\end{question}

\begin{question}[Beyond the bounded-coefficient class \src{07}]
\label{hgeo:us:q:classmaps}
Under which explicit foundational and continuity hypotheses is every
surreal-linear isometry or everywhere-defined adjointable map represented by
one set-supported series of ordinary bounded operators? Such an automatic
representation theorem would extend the negative examples to a larger class
of maps; \cref{hgeo:us:cor:surreal} alone does not establish it.
\end{question}

\noindent\emph{Status \mergetag.} For everywhere-defined class-linear maps
with an everywhere-defined class adjoint, in NBG without global choice and
with ordinary Hilbert coefficient spaces, this representation is
\cref{hgeo:thm:class}(ii) (sources~03 and~06), which gives
\cref{hgeo:us:cor:classmaps}. It remains open for isometries and other maps
without such an adjoint, and under other foundational hypotheses.

\begin{question}[Optimal gauges and perturbative costs \src{07}]
\label{hgeo:us:q:gauge}
Can the gauge freedom of \cref{hgeo:us:cor:gauge} be used to minimize
specified ordinary norms of successive coefficients, and does a minimizer
exist? The canonical first coefficient is fixed off the diagonal, but higher
coefficients interact nonlinearly. A satisfactory result should separate norm
minimization at each order from any unsupported notion of a single real-valued
norm of a whole Hahn operator.
\end{question}

\begin{question}[Machine-checked assembly \src{07}]\label{hgeo:us:q:lean}
Formalize \cref{hgeo:us:thm:main} for arbitrary ordered abelian exponent
groups, including the incomplete-family formula. This would turn the finite
rotations and support obstructions of the collection into a reusable verified
assembly interface. The infinite-family boundedness extension, not the finite
matrix identities, is the principal formalization milestone.
\end{question}

\subsection{What the finite computations test, and a formalization route}
\label{hgeo:us:sub:verification}
The suite \fn{code/07-unitary-synthesis-verify.py} performs exact rational
matrix power-series arithmetic with truncation, without floating point
\src{07}. It builds noncommuting finite-dimensional unitary series from
skew-adjoint generators, forms complete and incomplete projection families,
and independently executes
\eqref{hgeo:us:eq:constructiona}--\eqref{hgeo:us:eq:constructionb}. It checks
both unitary identities, the projection identities, intertwining, the
identities for $\Omega$, $P_\vee$ and $W$, agreement of the synthesized
columns with the necessity formula \eqref{hgeo:us:eq:necessity}, the first two
coefficients of the complete-family polar formula, and consistency of all
shorter jets of the canonical construction. Separate tests check the
two-dimensional identity \eqref{hgeo:us:eq:rationalblock} and the finite-block
rank-two identities behind \eqref{hgeo:us:eq:column} and the missing join. The
recorded run passes 159 named assertion groups; its log records the matrix
sizes and truncation orders; no randomness is used. Details and the rerun are
in \cref{hgeo:app:verification}.

These checks are diagnostic evidence for coefficient algebra. They cannot
prove well-ordered support in arbitrary ordered groups, the passage from a
dense ordinary subspace to bounded operators, the unboundedness across
infinitely many blocks, or novelty relative to published work; those claims
rest on the written proofs, not on the number of tests.

Source~07 proposes a formalization route. A Hahn-series algebra with
coefficients in bounded operators must expose the support lemma, the
involution and substitution into positive-order series; the ordinary
orthogonal direct-sum interface must provide extension of bounded maps from
the algebraic sum, the diagonal compression and equality by dense testing;
projection residue rigidity and the four formulas for $\Omega$, $P_\vee$, $W$
and $U$ then give the constructive proof, and necessity needs only the
compression identity \eqref{hgeo:us:eq:necessity}. The uniform counterexample
can be formalized without constructing matrices of dimension $4^{m^2}$ entry by
entry: it needs two orthonormal flat vectors in finite coordinate blocks, the
norm of their rank-two skew operator, and \eqref{hgeo:us:eq:allorder}. The
support example needs an ordered copy of $\Q$ and the sequence $1/n$. A
verified finite-jet package would be useful but would not replace the general
well-ordered support proof. No formalization is included or claimed, and the
collection's existing Lean coverage should not be read as covering these
labels.

\subsection{Summary of source 07}\label{hgeo:us:sub:summary}
A simultaneous infinite straightening problem has an exact answer in the
bounded-operator Hahn category: the canonical columns must synthesize to a
bounded operator at every exponent, and the active exponent set must be well
ordered. The conditions are necessary and sufficient, and the proof gives an
explicit unitary, a support bound and descent from every larger scale group.
In discrete scale, independently solvable finite jets for the whole family
already force one exact formal solution. The criterion is equivalent to a
precise atomic spectral calculus, whose countable additivity is coefficientwise
in the ordinary strong operator topology, not strong Hahn summability and not
valuation convergence. The uniformly controlled counterexample shows why the
distinctions matter: every individual coefficient can be uniformly bounded at
every order while collective synthesis, and even a selected Boolean join,
fail. The omnific-entry theorem marks the opposite boundary, where support
restrictions remove every nonconstant unitary deformation. Independent
correctness review, priority assessment and formal verification remain
separate tasks \src{07}. The limitations of source~07 are listed in
\cref{hgeo:app:nonclaims}.

\clearpage
\appendix
\section{Provenance}\label{hgeo:app:provenance}

\subsection{Section by section}

{\small\renewcommand{\arraystretch}{1.15}
\begin{longtable}{>{\raggedright\arraybackslash}p{0.2\textwidth}>{\raggedright\arraybackslash}p{0.74\textwidth}}
\caption{Provenance of each part of the report. Source section and theorem
numbers refer to the manuscripts, which are not shipped.}\label{hgeo:tab:provenance}\\
\toprule
Part of this report & Sources and choices\\
\midrule
\endfirsthead
\toprule
Part of this report & Sources and choices\\
\midrule
\endhead
\Cref{hgeo:sec:intro} & The merge, from 03 \S1 and 06 \S1 and abstracts.\\
\Cref{hgeo:sec:scalars} & 03 \S2; 06 \S\S2.1--2.2 (support lemma,
strong summability, the lexicographic example). Neumann inverse and binomial
series cited from the spectral report.\\
\Cref{hgeo:sec:construction} & 03 \S3 and 06 \S\S2.3, 3 printed once:
the positive form (03 Prop.~3.1, 06 Prop.~2.2) and spherical completeness
(03 Thm.~3.2, 06 Prop.~3.1) by citation of \replabel{ihs:hh:prop:inner} and
\replabel{duals:prop:spherical} with sketches; 06 Prop.~3.3 (the $c_{00}$
example); 03 Prop.~3.4.\\
\Cref{hgeo:sec:coordinates} & 03 \S4 and 06 \S2.4 printed once; 06 \S5:
Prop.~5.1 kept, Prop.~5.2 by citation of \replabel{duals:thm:riesz},
Thm.~5.3 by citation of \replabel{ihs:hh:thm:norm} with a rectangular
reduction \mergetag.\\
\Cref{hgeo:sec:adjoints} & 03 \S5 and 06 \S4: automatic adjoints (03
Thm.~5.1, 06 Thm.~4.2, Cor.~4.3) by citation of \replabel{ihs:hh:thm:adjoint}
and \replabel{ihs:hh:cor:autobounded}, with the reduction of the rectangular
case \mergetag; 06 Lem.~4.1 (Baire) and \S4.1 (calculus) kept; 03 \S5.2.\\
\Cref{hgeo:sec:splitting} & 03 \S6 and 06 \S6 printed once: 03 Prop.~6.2;
03 Lem.~6.3 = 06 Lem.~6.1 with both proofs; 03 Cor.~6.4; 03 Thm.~6.5 = 06
Thm.~6.2 derived from \replabel{ihs:fr:lem:rotation}; 03 Cor.~6.6 = 06
Thm.~6.3. \Cref{hgeo:prop:closedresidue} completes 06's remark after its
Thm.~6.3 and answers 03 Rem.~6.7 \mergetag.\\
\Cref{hgeo:sec:unitary} & 06 Cor.~6.4 and \S7 (Thms.~7.1, 7.2).\\
\Cref{hgeo:sec:graphs} & 03 \S7 (Prop.~7.1, Thm.~7.2, Rem.~7.3, Ex.~7.4,
Cor.~7.5); \cref{hgeo:lem:finitegraph} is 06 Lem.~9.2 and the step in 03's
proof of its Thm.~8.2; 06 \S8 (Thm.~8.1, Cor.~8.2, Prop.~8.3, the last
credited to \replabel{duals:thm:distance}); \cref{hgeo:sub:relation} is the
merge's, from an observation of the placement dossier.\\
\Cref{hgeo:sec:lattice} & 03 \S8 (Prop.~8.1, Thm.~8.2, Rem.~8.3), which
contains 06 Thm.~9.3; 06 \S9 (Thm.~9.1 in 03's coordinates, the Rickart
remark, Thm.~9.4, Cor.~9.5, Prop.~9.6).\\
\Cref{hgeo:sec:metric} & 06 \S10.\\
\Cref{hgeo:sec:schur,hgeo:sec:least,hgeo:sec:valuation} & 03 \S\S9--11. The
comparison with \replabel{ihs:fr:thm:fredholm} is the merge's; the finite
Moore--Penrose paragraph and Lem.~11.1 are cited from the finite
spectral-theory report, 03's elimination proof kept as a second route.\\
\Cref{hgeo:sec:global} & 03 \S12 and 06 \S11 printed once: 03 Prop.~12.1
with 06's coherence remarks; 03 Thm.~12.2 = 06 Thm.~11.1, stated without
global choice; 06 Cor.~11.2 and \eqref{hgeo:eq:stageintersection}; 03
Prop.~12.4 and 06 Prop.~11.3 printed once as the vector form of
\replabel{a:thm:discrete}.\\
\Cref{hgeo:sec:provenance} & 03 \S13 and 06 \S12, Appendix~A; the relations
to the collection and the status of the questions are the merge's.\\
\Cref{hgeo:app:checklist} & 03 Appendix~A, with rows for 06's results and the
merge; rows for 07 and 07 Appendix~B added in batch 34.\\
\Cref{hgeo:us:sec:intro} & 07 \S1 and abstract; 07 \S2 in
\cref{hgeo:us:sub:algebra}, its Lem.~2.1 printed once with
\cref{hgeo:lem:integralproj,hgeo:cor:orbits} (\cref{hgeo:us:lem:integral}),
its support principle and strong summability by citation of
\cref{hgeo:lem:support,hgeo:def:strong}. \Cref{hgeo:us:sub:source} and
\cref{hgeo:us:sub:conventions} are the merge's.\\
\Cref{hgeo:us:sec:criterion} & 07 \S\S3--4 (Def.~3.1, Thm.~3.2 and its
proof); \cref{hgeo:us:sub:relation} is the merge's.\\
\Cref{hgeo:us:sec:consequences} & 07 \S5: Cors.~5.1--5.3, Thm.~5.5,
Prop.~5.6. 07 Cor.~5.4 is \cref{hgeo:thm:finite}
(\cref{hgeo:us:rem:finite}).\\
\Cref{hgeo:us:sec:calculus} & 07 \S6: Def.~6.1, Thm.~6.2, Props.~6.3
and~6.4; \cref{hgeo:us:rem:joins} is the merge's.\\
\Cref{hgeo:us:sec:uniform} & 07 \S7: Thm.~7.1 and Rem.~7.2.\\
\Cref{hgeo:us:sec:support} & 07 \S8: Thm.~8.1, whose family contains that of
06's \cref{hgeo:prop:incoherent}; \cref{hgeo:us:rem:transport} is the
merge's.\\
\Cref{hgeo:us:sec:surreal} & 07 \S9: Cor.~9.1 and Thm.~9.2;
\cref{hgeo:us:cor:classmaps} is the merge's, from \cref{hgeo:thm:class}.\\
\Cref{hgeo:us:sec:questions} & 07 \S10 (verification and formalization
route), \S11 (Questions~11.1--11.12) and \S12 (conclusion);
\cref{hgeo:us:sub:status} and the status note after
\cref{hgeo:us:q:classmaps} are the merge's.\\
\Cref{hgeo:us:app:provenance} & 07 Appendix~A.\\
\bottomrule
\end{longtable}}

Where the sources overlap, the statement with the weakest hypotheses is
printed. The proofs attached to printed statements were checked for hidden
uses of stronger hypotheses: 06's descent proof uses choice only on sets, so
the statement without global choice (03's) is correct for both; 03's no-meet
proof and its transfer to an arbitrary infinite-dimensional $H$ use no
property of $\ell^2$ beyond containing three orthogonal copies of it.

\subsection{Labels}
Every label carries the prefix \replabel{hgeo:}. The placed text was source~03
with 69 bare labels; nothing in the collection cited them, and each was
prefixed unchanged (for example \replabel{thm:threshold} is
\replabel{hgeo:thm:threshold}); the other labels were added for source~06's
results, the merge, and the appendices.

Source~07 (batch 34) arrived with bare labels (for example
\replabel{thm:main}); nothing in the collection cited them. Its material
carries the sub-prefix \replabel{hgeo:us:}, unused before, with names chosen
for this report (\replabel{thm:main} is \replabel{hgeo:us:thm:main},
\replabel{thm:uniform-counter} is \replabel{hgeo:us:thm:uniform}). No existing
label was renamed or removed, and no earlier section, theorem or equation
number changed.

\subsection{Renamed symbols}

{\small\renewcommand{\arraystretch}{1.15}
\begin{longtable}{>{\raggedright\arraybackslash}p{0.08\textwidth}>{\raggedright\arraybackslash}p{0.3\textwidth}>{\raggedright\arraybackslash}p{0.52\textwidth}}
\caption{Symbols renamed from the sources.}\label{hgeo:tab:renames}\\
\toprule
Source & Source symbol & Here, and why\\
\midrule
\endfirsthead
\toprule
Source & Source symbol & Here, and why\\
\midrule
\endhead
06 & $K=\R((t^\Gamma))$, $F=k((t^\Gamma))$ & $\F$ and $\K$ swapped back:
$\F$ real, $\K$ complex, as in 03, the spectral report and the finite
spectral-theory report; $k$ is $\kk$.\\
06 & $\mathcal E_\Gamma(H)$, $\mathcal E_{\No}(H)$ & $\HH{H}$,
$\mathscr H_{\No}(H)$ (03's), because $E_\Gamma(V)$ is the algebraic
extension $\Span_\K V$ in the three-duals report.\\
06 & $\mathcal A_\Gamma(H,G)$, $\mathcal A_{>0}$, $v_{\mathrm{op}}$ &
$\BB{H}{G}$, $\Bpos{H}{G}$, $v$ (03's).\\
06 & $\mathcal U(H)$, $\mathcal U_{1,\Gamma}(H)$ & $\Un(H)$,
$\Un_{1,\Gamma}(H)$, since the spectral report's $\mathcal U(T)$ is a set of
norm bounds.\\
06 & $\tau=t^\delta$ & $q=t^\eta$ (03's); $\delta\mapsto\eta$ throughout.\\
06 & $S$ ($Se_n=e_n/n$), $D$ ($\diag(n)$) & $T$ (03's) and
$\Lambda=T^{-1}$; the spectral report writes $D$ for $\diag(1/n)$, and $D$
was also 03's rotation denominator.\\
06 & $M_\tau(A)$ & $M_q(A)$; $M_b(R)$ for relations (merge).\\
06 & $T(x,y)=\tau x-Sy$, $L=T^*T$ & $B(x,y)=Tx-qy$ (03's, input coordinates
exchanged) and $L=B^*B$.\\
06 & $R_n$ & $P_n'$, since $R$ is a relation.\\
06 & $D=P-Q$, $A=PQ+(I-P)(I-Q)$, $Q=\widehat{P_0}$ & written out, $W$
(03's), $P_0$.\\
06 & $A=\sum P_jQ_j$ (finite straightening) & $W$.\\
06 & $G=\widehat{G_0}+R$ (metric) & $\Theta=\widehat{\Theta_0}+\Theta_+$,
since $G$ is a Hilbert space and $R$ a relation; $R'$ is $\Theta'$.\\
06 & $h$ in its Thm.~6.2, Cor.~8.2 & $u$, $h'$; $h=(1/n)$ is reserved.\\
06 & $D$ ($c_{00}$) & $c_{00}$.\\
06 & ``projector'', ``orthogonally complemented'' & ``projection'',
``split'' (both words are used; they are synonyms).\\
03 & $D=I-(P-P_0)^2$ & written out.\\
03 & $C=\ker T_0^*$ and blocks $C_1,D,F_1$ & $N_*$ and
$A_{12},A_{21},A_{22}$, since $F$ is the real field and $D$ was overloaded.\\
03 & $G=Y\oplus C$ in the threshold proof & $G=Y\oplus Y^\perp$.\\
03 & $G_0$, $F$, $F'$ (Penrose proof) & $A^-$, $A'$, $A''$, since $G$ is a
Hilbert space and $F$ the field.\\
03 & $J$, $K_1$, $G$ (scale proof) & $\mathsf L$, $\mathsf R$,
$A^-_1$; $A^-_0$ in the case $S=0$.\\
03 & $U,V$ (Lemma 11.1) & $U_1$, $V_1$, since $U$ is a unitary.\\
03 & $\eta=v(b)$ in the threshold proof & $\eta'$, since $\eta$ is fixed by
$q=t^\eta$.\\
03 & $R=(I+X^*X)^{-1}$ in the graph chart & $R_X$.\\
07 & $\mathbb F\in\{\R,\C\}$, $K=\mathbb F((t^\Gamma))$ & $\kk$; $\F$ or
$\K$.\\
07 & $\mathscr H$, $\mathscr A$, $\mathscr A_{\ge0}$, $\mathscr A_{>0}$,
$I_H$ & $\HH{H}$, $\BH$, $\Bnn$, $\Bpos{H}{H}$, $I$.\\
07 & index set $I$; finite subsets $J$; subsets $A,B$; left-half set $J$;
basis index set $J$ & $\mathcal I$, $\mathcal J$, $\mathcal S,\mathcal S'$,
$\mathcal L$, $\mathcal N$, since $I$ is the identity and $J$ the graph-chart
column and the metric isometry.\\
07 & $F=I-E$ & $E^\perp$, since $F$ is the real Hahn field.\\
07 & $\mathcal D_E$ & $\Dg_E$, since $\mathcal D$ is the Schur block
matrix.\\
07 & $G=C^*C+F$ & $\Omega$, since $G$ is a Hilbert space.\\
07 & $Q=CG^{-1}C^*$ & $P_\vee$, since the report's $Q_j$, $Q_n$ are residue
projections (07's $E_j$); $P_\vee$ is the join.\\
07 & $T=C+(I-Q)F$, $B=(T^*T)^{-1/2}$ & $W$ (the $W$ of
\cref{hgeo:thm:finite} in the finite case) and $\abs W^{-1}$, since
$T=\diag(1/n)$ and $B$ is the row operator.\\
07 & $R$ (generic), $R$ (upper bound), $S=U^*RU$ & $Y$, $P'$, $P''$, since
$R$ is a relation and $S$ a support.\\
07 & $R_m=\exp(tA_m)$, $A_m$, $u_m,v_m$, $F_m$ & $V_m=\exp(tX_m)$, $X_m$,
$u_m,u_m'$, $\Pi_m$; $v$ is the valuation.\\
07 & $J$ ($2\times2$), $q_n=t^{1/n}$, $R_n$ & $\mathsf J$, $t^{1/n}$, $V_n$;
$q=t^\eta$ is fixed.\\
07 & $K$, $D$ (first-order equation), $W$ (gauge), $V$ (residue of $U$;
second straightener) & $Y$, $Y_c$, $\Xi$, $U_0$, $U'$; $\K$ is the complex
field.\\
07 & $\mathscr A_N$, $A_\lambda$, generic $T$, $\langle S\rangle$ &
$\mathscr B^{[N]}$, $N_\lambda$, $A$, $S^*$.\\
07 & inner product conjugate-linear in the first argument;
$xy^*:z\mapsto x\ip{y}{z}$; entries $\ip{e_j}{T_\gamma e_k}$ & linear in the
first argument; $xy^*:z\mapsto\ip{z}{y}x$; entries $\ip{A_\gamma e_k}{e_j}$
(the same maps and entries).\\
\bottomrule
\end{longtable}}

\subsection{Stale and incomplete repository statements}\label{hgeo:app:stale}
The pins \texttt{9a385d3} (03) and \texttt{a45d722} (06) are kept as
provenance. None of the cited reports changed between the pins and the
placement, so the corrections below concern omissions of the sources rather
than later changes of the collection.
\begin{enumerate}[leftmargin=*]
\item Neither source cites the rotation lemma \replabel{ihs:fr:lem:rotation},
which already states the direct rotation for every projection $P$ of the
algebra with $v(P-P^\circ)>0$; 03 says only that Part~IV ``uses the classical
direct rotation'' for clusters. Given the lemma, the shared normal-form
theorem is a corollary of the integral-projection lemma and the automatic
adjoint theorem (\cref{hgeo:thm:normalform}).
\item Source~03 does not compare its Fredholm theorems with
\replabel{ihs:fr:thm:fredholm}; \cref{hgeo:rem:fredholm} does.
\item Source~03 does not cite the finite spectral-theory report, whose
\replabel{thm:leastsquares}, \replabel{thm:scales} and
\replabel{lem:DeltaInvariant} are its finite Moore--Penrose paragraph and its
Lemma~11.1. 03's remark that its finite construction needs neither real
closedness nor divisibility is correct, but at 03's pin that report already
proved the SVD over a nondivisible Hahn field (\replabel{spec:thm:svd}), so
the finite results were available for every $\Gamma$.
\item Source~06's distance cut (its Prop.~8.3) does not credit
\replabel{duals:thm:distance}, which has exactly the same cut for the closed
hyperplane (\cref{hgeo:prop:distancecut}).
\item Neither source cites \replabel{a:thm:discrete} (present at both pins) or
the Lean-formalized \replabel{found:thm:discrete} for fine discreteness
(\cref{hgeo:prop:globalnets}).
\item Source~06 states its global descent theorem in NBG with global choice;
the proof does not use it (\cref{hgeo:thm:class} and the paragraph after
it).
\item Source~06 cites the spectral report by theorem numbers (6.1, 11.3,
13.2) and the duals report by Proposition~5.1, Proposition~8.1 and Sections
8--9; these numbers are correct at the current build and are replaced here by
the labels \replabel{ihs:hh:thm:adjoint}, \replabel{ihs:hh:thm:coercive},
\replabel{ihs:hh:thm:norm}, \replabel{duals:prop:spherical},
\replabel{duals:prop:norm}, \replabel{duals:thm:riesz}. Its description of
the spectral report as 99 pages is still that report's page count.
\item Source~03 lists unbounded operators with proper domains and
infinite-family projection joins as not developed, and makes no claim that
closed residue suffices; within this report source~06 treats the first two
(the second negatively), and \cref{hgeo:prop:closedresidue} settles the third
negatively (\cref{hgeo:q:orthoclosed,hgeo:q:notdeveloped}).
\item Source~06's non-lattice theorem is stated for $\ell^2$ only; 03's
theorem covers every infinite-dimensional $H$ and adds the finite-dimensional
converse (\cref{hgeo:thm:lattice}).
\item Both sources describe their delivery archives (source PDFs, build
records). The source manuscripts and PDFs are not shipped with this report;
the shipped files are listed in \cref{hgeo:app:verification}.
\item Source~07 (pin \texttt{fc27d87}, 61 commits before its placement
\texttt{a7a435f}) quotes this report's Question~15.3 and Part~IV of the
spectral report correctly; neither report changed between the pin and the
placement, apart from the placement itself.
\item Source~07 says its finite-family corollary recovers a statement of the
repository without naming it; it is \cref{hgeo:thm:finite}, and the two
constructions agree (\cref{hgeo:us:rem:finite}). It likewise names the
support-incoherence obstruction without a label; its
\cref{hgeo:us:thm:support} contains the family of \cref{hgeo:prop:incoherent}
with $\eta=1$.
\item Source~07 does not use, and asks for (its Question~11.10), an automatic
representation theorem for class maps. For maps with an everywhere-defined
class adjoint this is \cref{hgeo:thm:class}(ii), present at its pin; the merge
applies it in \cref{hgeo:us:cor:classmaps}, and the status note after
\cref{hgeo:us:q:classmaps} re-scopes the question.
\item Source~07's delivered build scripts and build report refer to its own
manuscript and PDF, which are not shipped (\cref{hgeo:app:verification}).
\end{enumerate}

\subsection{Source 07: snapshot, scope and dependency ledger}
\label{hgeo:us:app:provenance}
This subsection prints source~07's own provenance appendix.

\emph{Repository snapshot.} Source~07 identified its target at repository
revision \fn{fc27d87646ea3b7ed5e04e0a9c00a307dfe12fe2}, committed on
September 24, 2026 at 00:27:33 UTC, that is September 23, 2026 at 17:27:33 in
America/Los\_Angeles; the manuscript's date uses that local date. Repository
content may change after the pin. It inspected the main and documentation
catalogues and the guides of four reports: this report;
\emph{Infinite-Dimensional Hahn Spectral Theory}; \emph{Set-Sized Quotients of
the Omnific Integers} (\fn{docs/surreal/set-sized-quotients-of-omnific-integers});
and \emph{Omnific Integers and Omnific--Diophantine Geometry}
(\fn{docs/surreal/omnific-diophantine-geometry}). The last two ruled out
normalization and basic Diophantine rigidity as useful nonduplicative
targets. The detailed proofs in the long report sources were not audited. The
target statements and the comparison with finite straightening come from the
two operator-theory guides, and source~07's proofs are not assertions that
all claims of those reports were independently checked. In particular it does
not use an automatic class-operator representation theorem of a neighbouring
report in its surreal descent (\cref{hgeo:us:cor:surreal}); in this report that
theorem is \cref{hgeo:thm:class}, used only in the merge's
\cref{hgeo:us:cor:classmaps}.

\emph{Dependency ledger.} The proof of the main criterion uses the
positive-support lemma, ordinary bounded extension from a dense subspace,
ordinary orthogonal projections and the elementary coefficient identities of
\cref{hgeo:us:sub:algebra}; no spectral theorem for Hahn operators is used.
The functional calculus uses in addition the matrix-amplification norm
estimate for ordinary bounded maps and the ordinary diagonal representation on
a Hilbert direct sum. The examples use finite-dimensional rank-one norm
identities and unboundedness across infinitely many finite blocks. The omnific
application uses only the normal-form support condition.

\begin{center}
\small
\begin{tabular}{@{}>{\raggedright\arraybackslash}p{0.3\textwidth}>{\raggedright\arraybackslash}p{0.62\textwidth}@{}}
\toprule
Result & Essential proof obligation\\
\midrule
\Cref{hgeo:us:thm:main} & Coefficientwise bounded synthesis and legal support
before polar normalization\\
\Cref{hgeo:us:cor:extension} & The canonical columns are unchanged by exponent
embedding\\
\Cref{hgeo:us:thm:jets} & Every finite jet concerns the entire family, not a
finite subfamily\\
\Cref{hgeo:us:thm:calculus} & A rectangular completely bounded estimate gives
collective synthesis\\
\Cref{hgeo:us:thm:uniform} & Small coordinate columns coexist with unbounded
block norms\\
\Cref{hgeo:us:thm:support} & Bounded coefficient operators have a
non-well-ordered active support\\
\Cref{hgeo:us:thm:omnific} & Integral operator support and omnific scalar
support meet only at zero\\
\bottomrule
\end{tabular}
\end{center}

\emph{Literature status and a current announcement.} The literature check was
targeted, not exhaustive. The older projection and polar-decomposition
literature supplies conceptual precedents, and the novelty claim is restricted
to the exact combination and strengthening stated in \cref{hgeo:us:sub:why}. A
user-supplied Wikipedia page was used for orientation, not as a source of a
theorem. During its current-status check source~07 found a September~18, 2026
author announcement of a Lean proof of Conway's refinement conjecture
\cite{Abramov}; the announcement itself distinguishes machine checking from
independent mathematical review. That proof was not audited, is used by no
argument, and no conclusion about its eventual acceptance is drawn: source~07
neither treats that conjecture as untouched nor claims to have resolved it.
The citation is provenance only; the target is the operator assembly problem.

\section{What this report does not claim}\label{hgeo:app:nonclaims}
Every limitation, disclaimer and priority caveat stated by a source is kept
here, numbered per source; source~07's are items 07.1--07.26, and the batch-34
addition's are A.1--A.7.

\subsection{Source 03 (34 items)}
\begin{enumerate}[leftmargin=*,widest=00,label=03.\arabic*]
\item An AI-assisted research draft, not independently refereed or formally
verified; no proof assistant was run, and there is no inherited Lean
verification.
\item The construction, automatic adjoints, the direct-rotation formula and
the classical perturbation tools are credited, not claimed as new.
\item No invention is claimed of non-Archimedean inner products,
non-Archimedean failures of Riesz representation or orthogonal
complementation, formal deformation of projections, the Schur complement, or
the Moore--Penrose inverse.
\item The proposed contributions (the exact graph threshold, the
finite-dimensional/lattice dichotomy, the combined finite-defect
least-squares and sharp-valuation package) were not located by a targeted
comparison; priority is not certified, and failure of the searches is
evidence for a candidate contribution, not proof that none exists.
\item No named published conjecture is asserted to be solved.
\item Finite symbolic checks are supplementary, not proofs; they can catch
sign or transpose mistakes but do not establish support well-ordering,
infinite-dimensional boundedness, spherical completeness, the class argument,
or bibliographic priority.
\item The model does not retain ordinary Hilbert-norm convergence as the
topology on the constant copy of $H$.
\item The classification is of orthogonally split subspaces, not of all
valuation-closed subspaces; near-identity unitaries representing a split
subspace are not unique.
\item Positive-order power series are identities of Hahn series and do not
presuppose topological convergence of finite partial sums.
\item Common workspaces for specified sets of scalars do not reduce every
class-sized problem to a single set field, nor identify the full fine
topology with the workspace topology.
\item Not every positive element of $\F$ has a square root; only the roots of
squared norms are asserted.
\item The rank-one ultrametric $\exp(-v(x-y))$ is not the $\F$-valued Hilbert
norm.
\item Spherical completeness does not imply the projection theorem and does
not force ordinary completeness on coefficient subspaces.
\item The constant orthonormal basis need not be a valuation Schauder basis;
Parseval holds coefficientwise only.
\item Neither Riesz observation is claimed as new.
\item The automatic adjoint theorem does not classify all valuation-continuous
maps, and uses ordinary Hilbert completeness of the coefficient spaces.
\item No general bounded-inverse theorem over $\K$ is assumed.
\item The rotation formula is not a new result.
\item Closed residue is necessary, not claimed sufficient, for splitting.
\item Positivity of $I+A^*A$ does not alone imply surjectivity.
\item The no-nearest-point example asserts failure of an attained minimum
only, not a statement about a distance infimum.
\item The orthomodular-poset axioms do not assert that arbitrary binary meets
or joins exist; the general non-Archimedean nonorthomodularity phenomenon is
not new \cite{AN}.
\item Nontrivial $\Gamma$ is essential to the lattice dichotomy.
\item ``Fredholm'' is the finite algebraic kernel/cokernel condition, not a
compactoid-operator notion; the Schur calculation is standard; the finite
Moore--Penrose identities are classical.
\item The diagonal exponents are valuation invariants, not asserted
eigenvalue or singular-value exponents of an arbitrary infinite operator; no
spectral theorem is used to diagonalize an arbitrary Hahn operator.
\item The coherence statements are algebraic and do not assert that
valuation completions or closures commute with a noncofinal scale
enlargement.
\item The class section does not make the global vector class an ordinary
set-sized Hilbert space, and forms no class of all subspace classes.
\item The fine-net result concerns set-indexed nets only; no claim is made
about class-indexed Cauchy families.
\item The Fredholm assumption is sufficient, not asserted necessary, for an
adjointable Moore--Penrose inverse.
\item General tensor products, frame bounds, spectral measures, unbounded
operators with proper domains and infinite-family projection joins are not
developed in source~03 (see \cref{hgeo:q:notdeveloped} for what source~06
adds); no classical infinite-dimensional spectral theorem, measurement
postulate or physical interpretation is claimed from the existence of a
positive surreal norm.
\item The minor formula is not an algorithm for uncomputably presented
surreal numbers.
\item The repository comparison inspected the relevant portions only, not
the complete proof corpus or all unintegrated manuscripts; no novelty is
claimed relative to uninspected sources.
\item The comparisons with Aguayo--Nova and Avrachenkov--Lasserre were not
full-text audits (publication record and author text; abstract and
coauthor's description); neither projection deformation nor the rotation
formula originates here, and neither the analytic methods of
\cite{Avrachenkov,MS} nor classical generalized inverses are presented as new.
\item No theorem of the main set-sized treatment assumes a divisible group;
no ordinary coordinate sum is exchanged with a strong Hahn family without
qualification; nearest-point statements compare exact norms and do not
replace minimization by an unproved infimum.
\end{enumerate}

\subsection{Source 06 (33 items)}
\begin{enumerate}[leftmargin=*,widest=00,label=06.\arabic*]
\item An AI-assisted research draft with complete written proofs; not
independently refereed or proof-assistant verified; no surreals-specific
theorem is labelled Lean-verified.
\item Candidate contributions are separated from repository results and
classical operator theory; the literature comparison is targeted, not a
certification of priority.
\item No named published conjecture or open problem is claimed solved.
\item ``Candidate'' qualifies originality only, not the logical status of the
proofs.
\item Inherited, not claimed: the construction, automatic adjointability,
the failure of Riesz representation, the norm-attainment criterion,
spherical completeness, and the distinctions among represented, strong and
continuous duals.
\item The Kato--Nagy intertwiner formula is emphatically not new; the formula
is canonical as a formula, not the only near-identity unitary.
\item Sol\`er's theorem is applied only to set-sized fields and spaces, not to
proper classes; an abstract prediction from it is not an original
contribution.
\item The normal-form framework is classical; divisibility is used only when
mentioned.
\item The positive-word support theorem is imported, not inferred from
convergence.
\item Strong summability is neither the strong operator topology nor ordinary
convergence; formal summability has an established literature.
\item The space is larger than the algebraic tensor product; no density of
the algebraic scalar extension is assumed.
\item The valuation topology is that of a fixed workspace, not the topology
of all positive tolerances in the full surreal field.
\item Spherical completeness concerns valuation balls only, not arbitrary
ordered-field norm balls, and implies neither ordinary completeness of the
coefficient space nor complementation of closed subspaces.
\item The $c_{00}$ example is not orthoclosed and is not a witness to the
failure of orthomodularity.
\item Positive-order operator series are not justified by
$n\delta\to+\infty$; formal smallness is not a real operator-norm estimate.
\item The order bound is not a construction of an $\F$-valued Banach
$C^*$-algebra norm.
\item ``Chart'' asserts no differential or manifold structure.
\item Closed residue alone is insufficient; the bounded-coefficient
hypothesis is part of the classification.
\item No real norm convergence of the Baker--Campbell--Hausdorff law is
asserted; formal exponential correspondences are established material.
\item The finite qualification in simultaneous straightening is essential.
\item The metric proof does not assume that an algebraically bijective
adjointable operator has an adjointable inverse.
\item The global theorem quantifies over a given class map and its given
class adjoint; no set of all class maps is formed and Zorn's lemma is not
applied to a proper class of subspaces.
\item The global operator theorem is algebraic; it neither asserts useful
set-indexed fine convergence nor needs it.
\item Workspace coherence is algebraic and strong-summation coherence, not a
claim that intrinsic completions commute with arbitrary enlargement.
\item ``Hahn--Hilbert'' names a construction with selected Hilbert features;
it does not assert orthomodularity or the full classical package.
\item The spectral report's spectra are not reproved or expanded, and are not
identified with a real Banach-algebra spectrum.
\item No unrestricted projection-valued spectral measure is established; the
missing countable join forbids assigning the join of every countable
orthogonal family, not all spectral measures for particular operators nor
alternative coefficientwise summation rules; physical interpretations would
need further choices, and no empirical claim is made.
\item No unrestricted Riesz, projection, countably additive spectral-measure,
or field-valued least-operator-norm theorem is asserted.
\item The further directions are extensions, not claimed resolutions of named
problems; closed residue is not asserted to settle every orthoclosed
subspace.
\item The repository comparison used READMEs, article openings and stated
labels; it is not a claim to have checked every theorem of the 99-page
spectral report, every manuscript or every Lean declaration, and it does not
override the distinctions of \fn{docs/FORMALIZATION.md}.
\item No first non-Archimedean Hilbert space, first failure of Riesz
representation or first non-Archimedean spectral theorem is claimed;
Ishizuka's setting is not automatically comparable.
\item Priority is to be assessed at the level of the precise hypotheses and
combined classification; independent review and a wider bibliographic check
remain necessary.
\item The identified inputs (Baire, closed graph, bounded inverse, Riesz,
closed-operator adjoints, positive square roots, Hahn--Neumann support
calculus) are not proved from set theory; class replacement and normal forms
are used only globally; no assertion rests on numerical evidence; the finite
checks do not verify unbounded domains, infinite-dimensional boundedness,
Baire category, Hahn support arguments, proper-class reasoning or the
nonexistence of projections, are not a substitute for the proofs and are not
a Lean formalization; successful typesetting does not prove the mathematics.
\end{enumerate}

\subsection{The merge (6 items)}
\begin{enumerate}[leftmargin=*,widest=00,label=M.\arabic*]
\item \Cref{hgeo:thm:relation} is stated for closed linear relations and
scalars of positive valuation only; it gives neither the finite-scalar clause
of \cref{hgeo:thm:threshold} nor the orthogonal complement in
\cref{hgeo:thm:unboundedgraph}.
\item \Cref{hgeo:prop:closedresidue} concerns a subspace that is not
orthoclosed; the orthoclosed case of \cref{hgeo:q:orthoclosed} is open.
\item The comparison with \replabel{ihs:fr:thm:fredholm} recovers only its
determinant-free clauses.
\item The spectral report's cluster-rotation question
(\replabel{ihs:fr:sec:questions}) receives only necessary conditions from
sources~03 and~06; after batch~34, source~07's criterion answers it in
exact-criterion form only, not in the spectral-gap form in which it is posed
(\cref{hgeo:us:sub:status}). Its other two parts, and the questions of
\cref{hgeo:sub:questions} not marked as answered there, remain open.
\item Results printed by citation are not reproved here; their proofs are
those of the cited reports, which are themselves unrefereed drafts.
\item Nothing in this report is formalized in Lean, and
\fn{docs/FORMALIZATION.md} maps none of its labels.
\end{enumerate}

\subsection{Source 07 (26 items)}
\begin{enumerate}[leftmargin=*,widest=00,label=07.\arabic*]
\item An AI-assisted, unrefereed proof draft; priority has not been
independently established; independent correctness review, priority
assessment and formal verification remain separate tasks.
\item The principal claims are proved in the explicitly defined Hahn operator
category; they propose a resolution of a repository research question, not of
a named classical conjecture, and are neither a general
infinite-dimensional spectral theorem nor a solution of every spectral problem
over surreal scalars.
\item The setting is the concrete algebra $\mathcal B(H)((t^\Gamma))$, not a
Banach space over a non-Archimedean field with an unspecified completion;
membership in the algebra is part of each theorem.
\item The proposed contributions (the exact criterion with its
incomplete-family construction, scale-enlargement invariance, the finite-jet
equivalence, the atomic calculus equivalence, the uniformly controlled
counterexample with its missing join) are candidate contributions relative to
the inspected sources, not a certified exhaustive novelty claim.
\item Polar intertwiners, projection geometry and polar decomposition
\cite{Kato,CorachMaestripieri}, well-ordered support calculus \cite{Neumann}
and the omnific normal-form description \cite{LInnocenteMantova} are credited,
not advertised as new.
\item No unreviewed theorem of the collection is used; the collection's
guides fixed only the target and the comparison.
\item The collection's projection-lattice counterexamples and
support-incoherence obstruction are not claimed for the first time.
\item The work does not repeat Part~II of the spectral report, whose
row- and column-finite category is different.
\item The criterion is exact but is not an effective decision procedure for
arbitrary infinite data.
\item A straightener is unique only up to the commuting gauge; uniqueness is
not claimed without a specified choice.
\item The finite-family corollary recovers a statement already in the
collection (\cref{hgeo:thm:finite}).
\item The jet theorem asks for one jet for the whole family; it does not say
that straightening every finite subfamily gives a common straightener.
\item Uniqueness of the atomic calculus is a rigidity statement once a
straightener exists; it does not show that an arbitrary pointwise-defined
algebra map has coherent support or supplies a straightener, and automatic
regularity is open.
\item The atomic calculus is not strong Hahn summability, and its
coefficientwise additivity is not valuation convergence; no continuous-spectrum
spectral measure and no valuation-topologically countably additive general
spectral theorem is proved.
\item Weighted reconstruction does not assert that every normal Hahn operator
admits atomic data, nor that its algebra-relative spectrum consists only of
the labels.
\item The uniform example uses no global exponential of the unbounded
direct-sum generator; totality of an orthonormal family is not a synthesis
theorem.
\item The surreal descent concerns the specified category of set-supported
bounded-coefficient series, not all class-linear maps; automatic representation
theorems for general class-linear adjointable maps are left outside its scope;
the set-workspace assertions use neither a large-cardinal axiom nor global
choice.
\item Omnific matrix rigidity is an elementary boundary corollary, not a new
theorem of omnific factorization.
\item The closed-residue question (\cref{hgeo:q:orthoclosed}) is not settled.
\item The finite computations check finite identities only, with small block
dimensions; they do not prove well-ordered support in arbitrary ordered groups,
bounded extension from a dense subspace, unboundedness across infinitely many
blocks, or novelty. The claims rest on the written proofs, not on the number of
tests.
\item No Lean formalization is included or claimed; the collection's Lean
coverage should not be read as covering its labels; a verified finite-jet
package would not replace the support proof.
\item The repository comparison inspected the catalogues and four report
guides, not the detailed proofs of the long report sources; its proofs are
not assertions that the claims of those reports were checked.
\item The literature check was targeted, not exhaustive; the novelty claim is
restricted to the stated combination and strengthening; a user-supplied
Wikipedia page was used for orientation only.
\item The September~18, 2026 announcement of a Lean proof of Conway's
refinement conjecture is cited for current-status provenance only: it was not
audited or used, and no conclusion about its acceptance is drawn; the
conjecture is neither treated as untouched nor claimed resolved.
\item The further questions are separated from the established statements
and are not claimed resolved.
\item The proof-checking checklist marks where a plausible but incorrect
infinite-dimensional argument could arise; the finite tests make no claim to
discharge those obligations.
\end{enumerate}

\subsection{The batch-34 addition (7 items)}
\begin{enumerate}[leftmargin=*,widest=00,label=A.\arabic*]
\item \Cref{hgeo:us:cor:classmaps} depends on \cref{hgeo:thm:class}(ii)
(sources~03 and~06, unrefereed) and covers only class maps with an
everywhere-defined class adjoint.
\item \Cref{hgeo:us:rem:joins} gives a sufficient condition for joins, not a
characterization; \cref{hgeo:q:notdeveloped} stays open as a
characterization question.
\item The spectral report's first further question is answered only in
exact-criterion form; that report is not edited, and its other two questions
are untouched.
\item \Cref{hgeo:us:rem:finite,hgeo:us:rem:instances} identify constructions
and counterexamples of different sources; they add no existence claim beyond
the sources.
\item \Cref{hgeo:q:orthoclosed} stays open.
\item Stability under composition, asked for in \cref{hgeo:q:synthesis}, is not
treated.
\item Nothing in the addition is formalized, and \fn{docs/FORMALIZATION.md}
maps none of its \replabel{hgeo:us:} labels.
\end{enumerate}

\section{A dependency and hypothesis checklist}\label{hgeo:app:checklist}

{\small
\begin{longtable}{@{}>{\raggedright\arraybackslash}p{.24\textwidth}>{\raggedright\arraybackslash}p{.35\textwidth}>{\raggedright\arraybackslash}p{.33\textwidth}@{}}
\caption{Hypotheses and proof mechanisms of the main statements.}
\label{hgeo:tab:checklist}\\
\toprule
Statement & Essential hypotheses & Proof mechanism\\
\midrule
\endfirsthead
\toprule
Statement & Essential hypotheses & Proof mechanism\\
\midrule
\endhead
Inner product and norms & Set support; ordinary positive Hilbert form
 & Leading coefficient and binomial root\\[.4em]
Spherical completeness & Set-sized group and nested set of balls
 & Compatible fixed truncations\\[.4em]
Automatic adjoints & Everywhere-defined adjoint; ordinary complete
coefficient spaces & Coefficient closed graph, Baire, separation by
constants\\[.4em]
Projection normal form & Self-adjoint idempotent; ordinary Hilbert
coefficients & Integral leading projection and direct rotation\\[.4em]
Unitary group, finite straightening & Adjointable unitary; finite
orthogonal family & Residue, formal exp/log, polar normalization\\[.4em]
Infinite graph threshold & Ordinary bounded $T$; scalar of negative valuation
 & Residue $\ker T\oplus\ran T$ and closed-range inverse\\[.4em]
Infinitesimal unbounded graphs & Closed densely defined $A$; monomial
$q=t^\eta$, $\eta>0$ & Coefficientwise adjoint domain; $A^{**}=A$;
residue $\Dom A\oplus0$\\[.4em]
Scaled closed relations & Closed relation; $v(b)>0$ & Residue
$\dom R\oplus\mul R$; closed graph theorem for the operator part\\[.4em]
Lattice dichotomy & Nontrivial group; ordinary Hilbert coefficient space
 & Three-coordinate graph intersection; all lines split\\[.4em]
Positive kernel, missing join & Nontrivial group; $\ell^2$ & Residue of
$M_q(\Lambda)$; closedness of projection residues\\[.4em]
Countable straightening & Common support; $\Q\eta\subseteq\Gamma$ for the
second obstruction & Finitely many pairs at exponent $\eta$; well-ordered
$S+S$\\[.4em]
Metric rigidity & Strictly positive ordinary residue; $H\ne0$ &
Automatic adjoint of a bijective isometry; ordinary bounded inverse\\[.4em]
Fredholm least squares & Ordinary closed range; finite kernel and cokernel;
positive-order global perturbation & Schur reduction, finite orthogonal
complements, Penrose compression\\[.4em]
Exact inverse scale & Nonzero finite Schur rank
 & Integral diagonal reduction and a forced inverse pivot\\[.4em]
Scalar extension & Ordered group embedding; no cofinality needed
 & Transport of finite algebra and admissible strong sums\\[.4em]
Full-class adjoints & Set coefficient spaces; class replacement;
everywhere-defined adjoint; no global choice
 & Set image of constants, then the automatic adjoint proof\\[.4em]
Fine discreteness & Set index and all surreal radii
 & Surreal cut below a set of distances\\[.4em]
Infinite straightening criterion & Set-indexed orthogonal projections;
complete or incomplete residues & Canonical columns, dense bounded extension,
residue rigidity, formal polar normalization\\[.4em]
Jets, first-order equation & $\Gamma=\Z$; one jet for the whole family
 & Jet column identity; block commutator identities\\[.4em]
Atomic calculus, joins & Complete residues and complex $H$ (joins: any
straightener) & Rectangular completely bounded estimate; density on residual
blocks\\[.4em]
Uniform counterexample & Blocks of dimension $4^{m^2}$ & Flat vectors;
$C_1=X_m$ on each block; compression of a would-be join\\[.4em]
Support counterexample & $\Gamma\supseteq\Q\eta$ & Coefficients $\pm I$,
$\pm\mathsf J$ at $k/n$; active support $\Q_{\ge0}$\\[.4em]
Omnific rigidity & Entries in $\Oz$ or $\Oz[i]$ & Integral support meets
omnific support only at zero; Parseval\\
\bottomrule
\end{longtable}}

No theorem in the main set-sized treatment assumes a divisible group.
No infinite-dimensional theorem uses a finite matrix spectral theorem
to diagonalize an arbitrary Hahn operator. No unqualified exchange
of an ordinary Hilbert coordinate sum with a strong Hahn family is
performed. All nearest-point statements compare exact ordered-field
norms; they do not replace minimization by an unproved infimum.

Source~07's proof-checking checklist names the highest-value checks for its
part: that the canonical column series exists \emph{before} any
infinite-family operator manipulation; in the incomplete case, $CE^\perp=0$,
$\Omega E^\perp=E^\perp$, $P_\vee^2=P_\vee$ and the zero-residue argument
giving $P_i\le P_\vee$; both unitary identities, since $U^*U=I$ alone is not
enough in a general operator algebra and the proof uses invertibility for the
other; in the uniform example, the difference between legal finite-block
exponentials and the illegal exponential of an unbounded direct-sum
generator; that the missing-join argument excludes \emph{every} projection,
not one suggested infinite sum; and the distinctions among the three
summation conventions (\cref{hgeo:us:warn:sums}). These are where a plausible
but incorrect infinite-dimensional argument could most easily arise; the
finite tests cover parts of the algebraic identities and do not discharge
these obligations.

\section{Verification suites}\label{hgeo:app:verification}
Three finite suites accompany the report, byte-identical to the deliveries
under prefixed names; none is a proof of any infinite-dimensional, Baire,
domain, lattice, support or class assertion.
\begin{itemize}[leftmargin=*]
\item \fn{code/03-orthogonal-splitting-checks.py} (Python with SymPy) runs 16
named exact symbolic checks: graph-projection idempotence and self-adjointness;
the direct-rotation identities $W^*W=WW^*=I-(P-P_0)^2$ and $WP_0=PW$, and
unitarity and conjugation through order $q^8$; the Schur identity
$LAR=\diag(B,S)$; the four Penrose equations for a complex rational rank-two
matrix; the pole example of \cref{hgeo:ex:pole} (two inverse identities and
$v(A^{-1})=-4\eta$); and the triangular inverse-valuation formula on 125
positive exponent triples. Its recorded run is
\fn{data/03-orthogonal-splitting-check_results.json} (the \texttt{.txt}
file is a byte copy); \fn{data/03-orthogonal-splitting-build_report.txt}
records the build of the source's own PDF, which is not shipped.
\item \fn{code/06-hilbert-foundations-verify_finite.py} (standard-library
Python with exact fractions) runs 234 identity checks in six groups: rank-one
projections (60), the two-projection intertwiner before normalization (60),
graph projections (60), graph intertwiners (40), the positive block operator
(12), and formal prefixes to order 14 (2). Its recorded run is
\fn{data/06-hilbert-foundations-verification.json}.
\fn{code/06-hilbert-foundations-build.sh} is the source's build script for
its own manuscript, which is not shipped; it is kept as delivered.
\item \fn{code/07-unitary-synthesis-verify.py} (Python with SymPy~1.14.0,
pinned in \fn{data/07-unitary-synthesis-requirements.txt}) runs 159 named
assertion groups of exact rational truncated matrix power-series arithmetic.
The construction
\eqref{hgeo:us:eq:constructiona}--\eqref{hgeo:us:eq:constructionb} is tested
on a complete rank-one family of size~3 to order~6 (38 groups), an incomplete
family of size~4 with residual blocks $\{1,2\},\{3\}$ to order~6 (28) and a
complete block family of size~5 to order~5 (37): both unitary identities, the
identities for $\Omega$, $P_\vee$ and $W$, the necessity formula, the atom
identities and consistency of all shorter jets, and for the complete families
the polar formula and its first two coefficients. The flat-block identities
are tested for $m=1,2,3$ with block dimensions $d=1,4,9$, not $4^{m^2}$ (14
groups each), and the support block to order~10 (14). The recorded run is
\fn{data/07-unitary-synthesis-verification_results.json}; it records the SymPy
version but not the Python version.
\fn{code/07-unitary-synthesis-build.sh} and
\fn{code/07-unitary-synthesis-build.ps1} run pdfLaTeX three times on an
\fn{article.tex} next to themselves, which is source~07's manuscript and is
not shipped; they cannot build this report and are kept as delivered.
\fn{data/07-unitary-synthesis-BUILD_REPORT.json} records source~07's own
22-page build and the suite's count; its file names (\fn{article.pdf},
\fn{verification_results.json}) are delivery names.
\end{itemize}
Both suites were rerun for this report on copies (Python 3.14.4, SymPy
1.14.0): 03 printed ``16 checks passed, including 125 exponent triples'', and
its written JSON equals the recorded one except for the recorded Python
version (3.13.5); 06 printed status PASS with 234 checks, and its written JSON
equals the recorded one up to line endings. Both scripts write their outputs
(03 into a \fn{verification} directory next to the script, 06 into
\fn{verification.json} in the working directory by default), so they should
be run on a copy; the README gives the commands. For the batch-34 addition
the 07 suite was rerun on a copy (Python 3.14.4, SymPy 1.14.0): it printed
``Passed 159 exact assertion groups'', and its written JSON equals the
recorded one up to line endings (the Windows run wrote CRLF). By default it
writes \fn{verification_results.json} in the working directory, and its
\texttt{--output} option names another path; it too should be run on a copy.

\clearpage
\begin{thebibliography}{99}

\bibitem{AN}
J.~Aguayo and M.~Nova,
\emph{Non-Archimedean Hilbert like spaces},
Bulletin of the Belgian Mathematical Society--Simon Stevin
\textbf{14} (2007), no.~5, 787--797.
\href{https://doi.org/10.36045/bbms/1197908895}{doi:10.36045/bbms/1197908895}.

\bibitem{Avrachenkov}
K.~E. Avrachenkov and J.-B. Lasserre,
\emph{Analytic perturbation of generalized inverses},
Linear Algebra and its Applications \textbf{438} (2013), no.~4,
1793--1813. Author-deposited record:
\url{https://inria.hal.science/hal-00926609}.
Coauthor's description:
\url{https://www.laas.fr/fr/homepages/lasserre/softwares/singular-perturbations/}.

\bibitem{Bagayoko}
V.~Bagayoko, L.~S. Krapp, S.~Kuhlmann, D.~Panazzolo, and M.~Serra,
\emph{Automorphisms and derivations on algebras endowed with formal
infinite sums}, preprint, 2024.
\url{https://arxiv.org/abs/2403.05827}.

\bibitem{BW}
H.~Bursztyn and S.~Waldmann,
\emph{Deformation Quantization of Hermitian Vector Bundles},
arXiv:math/0009170, 2000.
\url{https://arxiv.org/abs/math/0009170}.

\bibitem{Conway}
J.~H. Conway,
\emph{On Numbers and Games}, second edition,
A K Peters, 2001. ISBN 978-1-56881-127-7.
\url{https://www.routledge.com/On-Numbers-and-Games/Conway/p/book/9781568811277}.

\bibitem{Freni}
P.~Freni, \emph{On vector spaces with formal infinite sums},
Applied Categorical Structures \textbf{34}, article 15 (2026).
\url{https://doi.org/10.1007/s10485-025-09844-w}.

\bibitem{Holland}
S.~S. Holland, Jr., \emph{Orthomodularity in infinite dimensions;
a theorem of M. Sol\`er},
Bulletin of the American Mathematical Society \textbf{32} (1995), 205--234.
\url{https://doi.org/10.1090/S0273-0979-1995-00593-8}.
Preprint: \url{https://arxiv.org/abs/math/9504224}.

\bibitem{Ishizuka}
K.~Ishizuka, \emph{A spectral theorem for a non-Archimedean valued field
whose residue field is formally real},
Extracta Mathematicae \textbf{39}(2) (2024), 235--254;
online publication January 3, 2025.
\url{https://doi.org/10.17398/2605-5686.39.2.235}.

\bibitem{Kato}
T.~Kato, \emph{Perturbation Theory for Linear Operators},
second edition, Classics in Mathematics reprint, Springer, 1995.
\url{https://doi.org/10.1007/978-3-642-66282-9}.

\bibitem{LackTobin}
S.~Lack and S.~Tobin, \emph{A characterisation for the category of Hilbert
spaces}, Applied Categorical Structures \textbf{33}, article 13 (2025).
\url{https://doi.org/10.1007/s10485-025-09805-3}.

\bibitem{MS}
J.~M\"uller and A.~Strohmaier,
\emph{The theory of Hahn meromorphic functions, a holomorphic
Fredholm theorem and its applications},
Analysis \& PDE \textbf{7} (2014), no.~3, 745--770.
\href{https://doi.org/10.2140/apde.2014.7.745}{doi:10.2140/apde.2014.7.745};
\url{https://arxiv.org/abs/1205.0236}.

\bibitem{Neumann}
B.~H. Neumann, \emph{On ordered division rings},
Transactions of the American Mathematical Society \textbf{66} (1949),
202--252.

\bibitem{Penrose}
R.~Penrose, \emph{A generalized inverse for matrices},
Proceedings of the Cambridge Philosophical Society \textbf{51} (1955),
406--413.

\bibitem{Simon}
B.~Simon,
\emph{Unitaries Permuting Two Orthogonal Projections},
arXiv:1703.05437, 2017, version 2.
\url{https://arxiv.org/abs/1703.05437}.
Formula (4) records the classical direct rotation.

\bibitem{SimonKato}
B.~Simon, \emph{Tosio Kato's work on non-relativistic quantum mechanics:
part 1}, Bulletin of Mathematical Sciences \textbf{8} (2018), 121--232,
especially the discussion of two projections.
\url{https://doi.org/10.1007/s13373-018-0118-0}.

\bibitem{Soler}
M.~P. Sol\`er, \emph{Characterization of Hilbert spaces by orthomodular
spaces}, Communications in Algebra \textbf{23} (1995), 219--243.
\url{https://doi.org/10.1080/00927879508825218}.

\bibitem{RepoSpectral}
Surreal project, \emph{Infinite-Dimensional Hahn Spectral Theory},
AI-assisted research report, directory
\fn{docs/surcomplex/infinite-dimensional-hahn-spectral-theory}, label prefix
\replabel{ihs:}; read at the sources' pins \texttt{9a385d3} and
\texttt{a45d722}, unchanged at the placement \texttt{a4dcb91}.

\bibitem{RepoDuals}
Surreal project, \emph{Three Duals at Surreal Scales: Strong Hahn Operators,
Completion Defects, and Invisible Functionals}, AI-assisted research report,
directory \fn{docs/surcomplex/three-duals-of-hahn-vector-spaces}, label
prefix \replabel{duals:}; same pins.

\bibitem{RepoFinite}
Surreal project, \emph{Spectral theory} (finite matrices over real closed
fields and Hahn workspaces), AI-assisted research report, directory
\fn{docs/surcomplex/spectral-theory}; labels \replabel{thm:leastsquares},
\replabel{thm:scales}, \replabel{lem:DeltaInvariant},
\replabel{spec:thm:svd}.

\bibitem{RepoAnalysis}
Surreal project, \emph{Analysis} report, directory
\fn{docs/surcomplex/analysis}; label \replabel{a:thm:discrete}.

\bibitem{RepoFound}
Surreal project, \emph{Foundations} report and the collection's notation and
reading guide, directory \fn{docs/foundations-and-computation/foundations};
label \replabel{found:thm:discrete}.

\bibitem{CorachMaestripieri}
G.~Corach and A.~Maestripieri,
\emph{Products of orthogonal projections and polar decompositions},
Linear Algebra and its Applications \textbf{434} (2011), 1594--1609.
\url{https://doi.org/10.1016/j.laa.2010.11.033};
\url{https://arxiv.org/abs/1011.5237}.

\bibitem{LInnocenteMantova}
S.~L'Innocente and V.~Mantova,
\emph{A factorisation theory for generalised power series and omnific
integers}, arXiv:1710.07304, version~5, January 22, 2024.
\url{https://arxiv.org/abs/1710.07304v5}. Cited by source~07 for normal-form
conventions and the omnific support description only.

\bibitem{Abramov}
D.~Abramov, \emph{How I Vibed a Proof of Conway's Conjecture}, author
announcement, September 18, 2026.
\url{https://overreacted.io/how-i-vibed-a-proof-of-conways-conjecture/}.
Cited by source~07 solely for a current-status provenance note; the announced
proof is not audited or used.

\end{thebibliography}
\end{document}
